% Options for packages loaded elsewhere
\PassOptionsToPackage{unicode}{hyperref}
\PassOptionsToPackage{hyphens}{url}
\documentclass[
]{article}
\usepackage{xcolor}
\usepackage{amsmath,amssymb}
\setcounter{secnumdepth}{-\maxdimen} % remove section numbering
\usepackage{iftex}
\ifPDFTeX
  \usepackage[T1]{fontenc}
  \usepackage[utf8]{inputenc}
  \usepackage{textcomp} % provide euro and other symbols
\else % if luatex or xetex
  \usepackage{unicode-math} % this also loads fontspec
  \defaultfontfeatures{Scale=MatchLowercase}
  \defaultfontfeatures[\rmfamily]{Ligatures=TeX,Scale=1}
\fi
\usepackage{lmodern}
\ifPDFTeX\else
  % xetex/luatex font selection
\fi
% Use upquote if available, for straight quotes in verbatim environments
\IfFileExists{upquote.sty}{\usepackage{upquote}}{}
\IfFileExists{microtype.sty}{% use microtype if available
  \usepackage[]{microtype}
  \UseMicrotypeSet[protrusion]{basicmath} % disable protrusion for tt fonts
}{}
\makeatletter
\@ifundefined{KOMAClassName}{% if non-KOMA class
  \IfFileExists{parskip.sty}{%
    \usepackage{parskip}
  }{% else
    \setlength{\parindent}{0pt}
    \setlength{\parskip}{6pt plus 2pt minus 1pt}}
}{% if KOMA class
  \KOMAoptions{parskip=half}}
\makeatother
\usepackage{color}
\usepackage{fancyvrb}
\newcommand{\VerbBar}{|}
\newcommand{\VERB}{\Verb[commandchars=\\\{\}]}
\DefineVerbatimEnvironment{Highlighting}{Verbatim}{commandchars=\\\{\}}
% Add ',fontsize=\small' for more characters per line
\newenvironment{Shaded}{}{}
\newcommand{\AlertTok}[1]{\textcolor[rgb]{1.00,0.00,0.00}{\textbf{#1}}}
\newcommand{\AnnotationTok}[1]{\textcolor[rgb]{0.38,0.63,0.69}{\textbf{\textit{#1}}}}
\newcommand{\AttributeTok}[1]{\textcolor[rgb]{0.49,0.56,0.16}{#1}}
\newcommand{\BaseNTok}[1]{\textcolor[rgb]{0.25,0.63,0.44}{#1}}
\newcommand{\BuiltInTok}[1]{\textcolor[rgb]{0.00,0.50,0.00}{#1}}
\newcommand{\CharTok}[1]{\textcolor[rgb]{0.25,0.44,0.63}{#1}}
\newcommand{\CommentTok}[1]{\textcolor[rgb]{0.38,0.63,0.69}{\textit{#1}}}
\newcommand{\CommentVarTok}[1]{\textcolor[rgb]{0.38,0.63,0.69}{\textbf{\textit{#1}}}}
\newcommand{\ConstantTok}[1]{\textcolor[rgb]{0.53,0.00,0.00}{#1}}
\newcommand{\ControlFlowTok}[1]{\textcolor[rgb]{0.00,0.44,0.13}{\textbf{#1}}}
\newcommand{\DataTypeTok}[1]{\textcolor[rgb]{0.56,0.13,0.00}{#1}}
\newcommand{\DecValTok}[1]{\textcolor[rgb]{0.25,0.63,0.44}{#1}}
\newcommand{\DocumentationTok}[1]{\textcolor[rgb]{0.73,0.13,0.13}{\textit{#1}}}
\newcommand{\ErrorTok}[1]{\textcolor[rgb]{1.00,0.00,0.00}{\textbf{#1}}}
\newcommand{\ExtensionTok}[1]{#1}
\newcommand{\FloatTok}[1]{\textcolor[rgb]{0.25,0.63,0.44}{#1}}
\newcommand{\FunctionTok}[1]{\textcolor[rgb]{0.02,0.16,0.49}{#1}}
\newcommand{\ImportTok}[1]{\textcolor[rgb]{0.00,0.50,0.00}{\textbf{#1}}}
\newcommand{\InformationTok}[1]{\textcolor[rgb]{0.38,0.63,0.69}{\textbf{\textit{#1}}}}
\newcommand{\KeywordTok}[1]{\textcolor[rgb]{0.00,0.44,0.13}{\textbf{#1}}}
\newcommand{\NormalTok}[1]{#1}
\newcommand{\OperatorTok}[1]{\textcolor[rgb]{0.40,0.40,0.40}{#1}}
\newcommand{\OtherTok}[1]{\textcolor[rgb]{0.00,0.44,0.13}{#1}}
\newcommand{\PreprocessorTok}[1]{\textcolor[rgb]{0.74,0.48,0.00}{#1}}
\newcommand{\RegionMarkerTok}[1]{#1}
\newcommand{\SpecialCharTok}[1]{\textcolor[rgb]{0.25,0.44,0.63}{#1}}
\newcommand{\SpecialStringTok}[1]{\textcolor[rgb]{0.73,0.40,0.53}{#1}}
\newcommand{\StringTok}[1]{\textcolor[rgb]{0.25,0.44,0.63}{#1}}
\newcommand{\VariableTok}[1]{\textcolor[rgb]{0.10,0.09,0.49}{#1}}
\newcommand{\VerbatimStringTok}[1]{\textcolor[rgb]{0.25,0.44,0.63}{#1}}
\newcommand{\WarningTok}[1]{\textcolor[rgb]{0.38,0.63,0.69}{\textbf{\textit{#1}}}}
\usepackage{longtable,booktabs,array}
\newcounter{none} % for unnumbered tables
\usepackage{calc} % for calculating minipage widths
% Correct order of tables after \paragraph or \subparagraph
\usepackage{etoolbox}
\makeatletter
\patchcmd\longtable{\par}{\if@noskipsec\mbox{}\fi\par}{}{}
\makeatother
% Allow footnotes in longtable head/foot
\IfFileExists{footnotehyper.sty}{\usepackage{footnotehyper}}{\usepackage{footnote}}
\makesavenoteenv{longtable}
\setlength{\emergencystretch}{3em} % prevent overfull lines
\providecommand{\tightlist}{%
  \setlength{\itemsep}{0pt}\setlength{\parskip}{0pt}}
% =============================================================================
% LaTeX preamble for nox-mem paper
% Purpose: Configure Unicode-aware monospace font to improve xelatex build quality
%
% Context: PR #316 identified 18 Unicode character warnings from xelatex build.
% Root cause: lmmono10 (default monospace font in pandoc) lacks glyph coverage for
% mathematical operators (∈ ≈ ≤ ≥ ↔) and emoji (🚫 ❌), forcing substitution warnings.
%
% Solution: Replace lmmono10 with Courier New, a widely-available monospace font
% with better Unicode glyph coverage. This reduces monospace-related warnings from
% ~12 to ~4 (remaining emoji-only warnings are unavoidable PDF limitation).
%
% Technical trade-offs:
%   - Courier New: universally available (macOS, Linux, Windows via TeX Live)
%   - Character pitch: slightly wider than lmmono10 (0.9 scale applied for balance)
%   - Symbol rendering: Math operators render as text boxes in PDF (expected)
%   - Emoji: PDF core fonts don't support U+1F*** range; fallback is □ (acceptable)
%
% Result: Warning count reduced from 18 → 4-6 (emoji-only, inherent PDF limitation)
% References: PR #316 (flagged), PR #238 (xelatex wrapper), PR #264 (Postman v2.1.0)
% =============================================================================

\usepackage{fontspec}
\usepackage{amssymb}
\usepackage{amsmath}

% Configure monospace font with better Unicode support
% Courier New is bundled with all major TeX distributions (TeX Live, MacTeX, MiKTeX)
% and handles technical character requests more gracefully than Latin Modern Mono
\setmonofont{Courier New}[Scale=0.9]
\usepackage{bookmark}
\IfFileExists{xurl.sty}{\usepackage{xurl}}{} % add URL line breaks if available
\urlstyle{same}
\hypersetup{
  hidelinks,
  pdfcreator={LaTeX via pandoc}}

\author{}
\date{}

\begin{document}

\section{nox-mem: Pain-Weighted Hybrid Memory for LLM
Agents}\label{nox-mem-pain-weighted-hybrid-memory-for-llm-agents}

\textbf{nox-mem v3.8 --- §5 fifth revision May--June 2026 (Wave A +
EverMemBench 5-batch Phase D/G/H v2 + Lab Q1 standalones + Wave B/C
composability + Backbone Matrix Gemini-3-flash + LoCoMo dual + MuSiQue +
HotPotQA dual SOTA + LongMemEval cross-bench + Production operational
characteristics)}

\textbf{Paper version:} v1.0.0 (2026-06-30) --- frozen for arXiv
submission; changelog in \texttt{paper/CHANGELOG.md}

\textbf{Author:} Luiz Antonio Busnello (Toto) \textbf{Platform:}
OpenClaw Autonomous Agent Platform \textbf{Infrastructure:} Hostinger
KVM4, Tailscale VPN, Debian Linux

\begin{center}\rule{0.5\linewidth}{0.5pt}\end{center}

\subsection{Abstract}\label{abstract}

We introduce \textbf{nox-mem}, a persistent memory system for autonomous
LLM agents organized around a single design principle:
\textbf{pain-weighted hybrid memory with shadow discipline --- yours by
design}. Every retrieval and retention decision is governed by an
evidence-weighted salience formula in which \emph{pain} --- an
operator-assigned severity in {[}0.1, 1.0{]}, persisted on every chunk
--- is a first-class signal (its isolated retrieval effect is
directional and not yet statistically significant, §7.1, and
section-aware ranking is the dominant empirical driver, §5.1.3); the
formula itself evolved from multiplicative to additive form through the
system's own shadow-mode telemetry (§5.1.2), and every ranking change is
gated by a mandatory shadow phase before production activation. Compared
to existing memory systems (mem0, Letta, Zep, EverOS, LightRAG, and
MeMo), nox-mem offers several advantages: \textbf{(a)} \emph{live
writeback with sub-second indexing} (inotifywait-driven, no batch
retrain or daily reindex required), \textbf{(b)} \emph{typed temporal
decay} (per-\texttt{chunk\_type} retention windows with never-decay for
feedback/person --- see §2 and §8.2), \textbf{(c)} \emph{full provenance
to chunk and source} (every retrieval result carries \texttt{chunk\_id}
+ \texttt{source\_file}, every destructive op is wrapped in
\texttt{withOpAudit()} with a VACUUM INTO pre-snapshot), \textbf{(d)}
\emph{first-class self-evolution} through a
\texttt{crystallize}/\texttt{reflect}/\texttt{consolidate} triad that
promotes high-hit pending items to durable lessons and synthesizes
cross-session insights nightly (§3.4), and \textbf{(e)} \emph{zero
vendor lock-in}: a single SQLite file, MIT-licensed, with
provider-agnostic embeddings (Gemini default, swappable). Deployed in
production since March 14, 2026, the system manages 94.9k+ memory chunks
with \textasciitilde99.99\% vector coverage, \textasciitilde15.6k
knowledge graph entities with \textasciitilde21.5k relations (as of
2026-06-04, post corpus-hygiene pass), and serves 6 specialized AI
agents with isolated yet interconnectable memory spaces. On the
entity-flavored golden set (n=100), the full ablation stack reaches
\textbf{nDCG@10 = 0.6237} (+78.8\% over the G3 pre-Wave-A baseline;
§5.1.1). A \textbf{Conditional Hard Mutex} (G10d) gates section and
source-type boosts by query entity count, recovering multi-hop (+1.58\%
nDCG@10) and adversarial (+3.04\% nDCG@10, +6.25\% MRR) regressions
(§5.1.4). On the EverMemBench cross-system benchmark (5-batch, n=3,121),
nox-mem scores \textbf{62.22\%} with Gemini-2.5-flash (+2.95 pp over
MemOS; §5.1.5), \textbf{51.68\%} with GPT-4.1-mini (+9.13 pp over MemOS,
95\% CI {[}49.88, 53.49{]}; §5.1.6), and \textbf{63.28\% Overall +
88.42\% Memory Awareness composite with Gemini-3-flash --- both SOTA
versus the published MemOS Table 4 baseline (GPT-4.1-mini)} (+20.73 pp
Overall, +32.74 pp MA; §5.1.10). On classical multi-hop QA, nox-mem
achieves dual SOTA without specialized fine-tuning: \textbf{MuSiQue dev
F1 58.62\%} (+22.82 pp over IRCoT, +8.92 pp over EX(SA); §5.2.1) and
\textbf{HotPotQA dev distractor ans\_F1 73.37\%} (above DPR+FiD reader
SOTA; §5.2.2). LoCoMo cross-bench retrieval@10 strict reaches
\textbf{74.52\%}, above Mem0's \emph{published} end-to-end answer-F1
SOTA of 66.88\% (a different metric, not a same-corpus head-to-head;
§5.3.1); the F1 push is rank-5 (51.85\%, above Zep / LangMem)
constrained by a verbosity gap (§5.3.2). The §5.4 resolution of the
EverMemBench F\_MH paradox shows the 3--7\% absolute is a
corpus-structural property, not a multi-hop reasoning ceiling.
Production characteristics are strong across three self-hosted
operational axes: \textbf{KG path retrieval p50 = 2.9 ms} (sub-10 ms
class, §5.7.1), \textbf{cost \$0/query KG path and \textasciitilde667×
cheaper than Mem0 Cloud on hybrid} (§5.7.2), and \textbf{399 MB RSS}
(resident set size) \textbf{self-hosted single-process} footprint
(§5.7.3). The cross-bench LongMemEval validation (n=300, §5.6) confirms
the same per-category fingerprint across an orthogonal benchmark
distribution. Wave 2 (2026-05-31 -- 2026-06-02, §5.5.4--§5.5.8)
contributes an empirical per-backbone × per-knob composability study:
three Wave A single-stage retrieval knobs (KG path, AC, MQ) transfer at
only 0--40\% efficiency from gpt-4.1-mini to Gemini-3-flash (D75),
establishing that retrieval-stage compensation mechanisms attenuate on
stronger backbones; IterB ReAct (§5.5.2) remains the sole validated
F\_MH lever on Gemini-3-flash at +2.01 pp.~An architectural lock
discovery (§5.5.7) reveals that composability projections must account
for both backbone-portability and code-level architectural composability
--- two independent non-trivial requirements. The Q4 cross-system
comparison (§6) provides pre-registered, same-corpus results against
five competing memory systems (two, Mem0 and agentmemory, produced
head-to-head quality numbers; three were documented deployment
non-runs). Under each system's native embedder the two leaders
\textbf{split} --- Mem0 wins LoCoMo (nDCG@10 0.469 vs 0.426), nox-mem
wins LongMemEval (§6.3); controlling the embedding provider (both
Gemini-3072d, full n=2,482) \textbf{inverts the split}, with nox-mem
outperforming Mem0 on both datasets (LongMemEval 0.526 vs 0.406; LoCoMo
0.495 vs 0.441) and all five represented categories --- a task-type
ablation rules out the one asymmetry favoring nox-mem, with three
residual confounds declared (§6.3.2).

\begin{center}\rule{0.5\linewidth}{0.5pt}\end{center}

\subsection{1. Introduction}\label{introduction}

\subsubsection{1.1 Problem Statement}\label{problem-statement}

Large Language Model (LLM) agents operating in production environments
face a fundamental limitation: context window ephemerality. When a
conversation ends or context is compacted, the agent loses accumulated
knowledge, decisions, and operational state. For multi-agent systems
where specialized agents collaborate on complex tasks, this problem
compounds --- agents cannot learn from each other's experiences, cannot
recall past decisions, and cannot build institutional knowledge over
time.

\subsubsection{1.2 Design Goals}\label{design-goals}

nox-mem was designed with four core objectives:

\begin{enumerate}
\def\labelenumi{\arabic{enumi}.}
\tightlist
\item
  \textbf{Persistent Memory}: Survive context window resets, session
  boundaries, and agent restarts
\item
  \textbf{Intelligent Retrieval}: Return semantically relevant results,
  not just keyword matches
\item
  \textbf{Cross-Agent Intelligence}: Enable knowledge sharing across
  isolated agent workspaces
\item
  \textbf{Operational Autonomy}: Self-maintain through automated
  consolidation, pruning, and indexing
\end{enumerate}

\subsubsection{1.3 Scope}\label{scope}

The system operates within the OpenClaw platform, serving 6 AI agents
(Nox, Atlas, Boris, Cipher, Forge, Lex) on a single VPS with 4 vCPUs and
8GB RAM. Each agent has a distinct role and memory profile. The
workspace (shared memory) and individual agent databases form a
federated memory architecture.

\subsubsection{1.4 Related Memory Systems and the Six
Gaps}\label{related-memory-systems-and-the-six-gaps}

The published memory-for-LLM-agents literature spans roughly three
families: (i) \emph{vector-store wrappers with metadata layers} ---
\textbf{mem0} \footnote{mem0ai/mem0 --- open-source memory layer for LLM
  agents (PostgreSQL + Qdrant backend, OpenAI embeddings by default).
  github.com/mem0ai/mem0. Used in §1.4, §6.3, Table 2.}, \textbf{Letta}
\footnote{Letta (formerly MemGPT) --- agent-loop memory architecture
  with archival/recall memory separation. github.com/letta-ai/letta.
  Used in §1.4, §6.3, Table 2.}; (ii) \emph{temporal- and
provenance-aware memory services} --- \textbf{Zep} \footnote{Zep ---
  temporal knowledge-graph memory service with summarization.
  github.com/getzep/zep. OpenAI embedding is the hardcoded default in
  the OSS distribution. Used in §1.4, §6.3, Table 2.}, \textbf{memanto};
(iii) \emph{KG-augmented and graph-fused retrieval} ---
\textbf{LightRAG} \footnote{Guo et al., \emph{LightRAG: Simple and Fast
  Retrieval-Augmented Generation}, EMNLP 2025 (HKU).
  github.com/HKUDS/LightRAG (\textasciitilde35k stars, MIT). Cited in
  §1.4 as a KG-augmented baseline; §3.4 references its LLM-summarized
  incremental KG-merge pattern as a forward-looking optimization
  (LightRAG-style summarization parking-lotted until KG density
  \textgreater=10× current; see
  \texttt{docs/COMPETITIVE-ANALYSIS-2026-05-19.md}).} (HKU, EMNLP 2025),
\textbf{HippoRAG2} \footnote{HippoRAG2 --- graph-augmented retrieval
  with Personalized PageRank over an entity-relation graph. Cited as a
  graph-baseline peer in §1.4 and §6.}; and (iv) \emph{parametric-memory
paradigms}, most recently \textbf{MeMo} \footnote{arXiv 2605.15156v2,
  \emph{MeMo: Towards Language Models with Associative Memory
  Mechanisms} (parametric reflections folded into model weights). Cited
  as the design opposite of nox-mem's externalized, inspectable memory
  paradigm. §1.4, abstract.} (which folds reflections into model weights
via continued pretraining --- the design opposite of ours).
\textbf{EverMind-AI/EverOS} \footnote{EverMind-AI / EverOS ---
  github.com/EverMind-AI (\textasciitilde5k stars, Apache 2.0).
  Publishes EverMemBench dataset and an EvoAgentBench-framed evolution
  loop. The only memory-OS peer in our taxonomy that publishes its own
  benchmark; §3.4 is the direct narrative counter to EvoAgent framing.
  Honest-comparison action item: run nox-mem on EverMemBench (queued as
  F4 in §7.2).} occupies a distinct slot: it is the only memory OS in
this space that publishes its own benchmark dataset (EverMemBench) and
reports threshold numbers, raising the bar for honest cross-system
comparison (§6).

Across these systems, six recurring gaps appear in the design space.
Each gap motivates a concrete subsystem of nox-mem. Table 1 summarizes
who covers what:

\textbf{Table 1 --- Comparison of desirable memory-system properties
(Six Gaps).}

{\def\LTcaptype{none} % do not increment counter
\begin{longtable}[]{@{}
  >{\raggedright\arraybackslash}p{(\linewidth - 16\tabcolsep) * \real{0.0732}}
  >{\raggedright\arraybackslash}p{(\linewidth - 16\tabcolsep) * \real{0.0732}}
  >{\centering\arraybackslash}p{(\linewidth - 16\tabcolsep) * \real{0.1220}}
  >{\centering\arraybackslash}p{(\linewidth - 16\tabcolsep) * \real{0.1220}}
  >{\centering\arraybackslash}p{(\linewidth - 16\tabcolsep) * \real{0.1220}}
  >{\centering\arraybackslash}p{(\linewidth - 16\tabcolsep) * \real{0.1220}}
  >{\centering\arraybackslash}p{(\linewidth - 16\tabcolsep) * \real{0.1220}}
  >{\centering\arraybackslash}p{(\linewidth - 16\tabcolsep) * \real{0.1220}}
  >{\centering\arraybackslash}p{(\linewidth - 16\tabcolsep) * \real{0.1220}}@{}}
\toprule\noalign{}
\begin{minipage}[b]{\linewidth}\raggedright
\#
\end{minipage} & \begin{minipage}[b]{\linewidth}\raggedright
Gap
\end{minipage} & \begin{minipage}[b]{\linewidth}\centering
nox-mem
\end{minipage} & \begin{minipage}[b]{\linewidth}\centering
mem0
\end{minipage} & \begin{minipage}[b]{\linewidth}\centering
Letta
\end{minipage} & \begin{minipage}[b]{\linewidth}\centering
Zep
\end{minipage} & \begin{minipage}[b]{\linewidth}\centering
EverOS
\end{minipage} & \begin{minipage}[b]{\linewidth}\centering
LightRAG
\end{minipage} & \begin{minipage}[b]{\linewidth}\centering
MeMo
\end{minipage} \\
\midrule\noalign{}
\endhead
\bottomrule\noalign{}
\endlastfoot
1 & Static injection (memory frozen at session start) & PASS live
writeback (inotifywait \textless1s) & NB: partial & PASS & PASS & PASS &
FAIL batch & FAIL retrain \\
2 & No temporal decay (every chunk weighed the same forever) & PASS
salience × recency, typed retention (§8.2) & FAIL & FAIL & PASS & NB:
partial & FAIL & FAIL \\
3 & No provenance (cannot trace a result back to a source) & PASS
\texttt{chunk\_id} + \texttt{source\_file} (§4) & NB: partial & PASS &
PASS & PASS & PASS & FAIL baked in weights \\
4 & Flat memory (no hierarchy / sectioning / typed structure) & PASS KG
+ entity files + \texttt{section\_boost} (§3.1, §8) & FAIL & FAIL & PASS
& PASS hyper & PASS dual-level & FAIL \\
5 & No writeback (cannot promote findings into durable memory) & PASS
\texttt{crystallize} + \texttt{reflect} + \texttt{consolidate} (§3.4) &
NB: partial & PASS & PASS & PASS EvoAgent & FAIL & FAIL \\
6 & Indexing delay (changes invisible until nightly batch) & PASS
inotifywait \textless1s, idempotent re-ingest (§3.1) & NB: & PASS & PASS
& NB: & NB: batch & FAIL retrain \\
\end{longtable}
}

\textbf{Why each gap matters, and how nox-mem closes it:}

\begin{enumerate}
\def\labelenumi{\arabic{enumi}.}
\tightlist
\item
  \textbf{Static injection.} A memory system that only reads at the
  start of a session cannot observe what the agent does \emph{during}
  the session. nox-mem's watcher \footnote{\texttt{nox-mem-watcher}
    systemd service running \texttt{inotifywait} on \texttt{memory/}
    directories. See \texttt{docs/ARCHITECTURE.md} and §3.1 of this
    paper.} re-ingests every modified \texttt{.md} / \texttt{.json} file
  within \textasciitilde1 second of the write, with idempotent chunk
  replacement (§3.1).
\item
  \textbf{No temporal decay.} Treating a daily note from 91 days ago the
  same as a crystallized decision floods retrieval with stale noise.
  nox-mem assigns each \texttt{chunk\_type} a \texttt{retention\_days}
  window (V8 schema column; \texttt{feedback}/\texttt{person} =
  never-decay, \texttt{lesson} = 180d,
  \texttt{decision}/\texttt{project} = 365d, \texttt{daily} = 90d) and
  folds it into the salience formula (§8.2).
\item
  \textbf{No provenance.} Parametric and opaque-vector memories cannot
  answer ``\emph{why} did you return this?''. Every nox-mem result
  carries the originating \texttt{chunk\_id} and \texttt{source\_file},
  every destructive operation goes through \texttt{withOpAudit()} with a
  \texttt{VACUUM\ INTO} pre-snapshot to
  \texttt{/var/backups/nox-mem/pre-op/}, and the \texttt{ops\_audit}
  table is append-only.
\item
  \textbf{Flat memory.} Without structure, hybrid search collapses to
  ``more recent or more frequent wins''. nox-mem layers (a) FTS5 (SQLite
  full-text search) over chunk text, (b) typed retention, (c) an
  LLM-extracted knowledge graph (§8), and (d) an entity-file format
  (\texttt{frontmatter} / \texttt{compiled} / \texttt{timeline}
  sections) with section-aware \texttt{section\_boost} --- the dominant
  driver of the +78.8\% gain in §5.
\item
  \textbf{No writeback.} A system that can only \emph{be read} cannot
  grow. nox-mem's self-evolution triad --- \texttt{crystallize},
  \texttt{reflect}, \texttt{consolidate} --- is described in §3.4 and is
  the most direct counter to EverOS's EvoAgentBench framing \footnote{EverMind-AI
    / EverOS --- github.com/EverMind-AI (\textasciitilde5k stars, Apache
    2.0). Publishes EverMemBench dataset and an EvoAgentBench-framed
    evolution loop. The only memory-OS peer in our taxonomy that
    publishes its own benchmark; §3.4 is the direct narrative counter to
    EvoAgent framing. Honest-comparison action item: run nox-mem on
    EverMemBench (queued as F4 in §7.2).}.
\item
  \textbf{Indexing delay.} Daily-batch reindex makes ``new memory'' a
  24h-resolution operation. inotifywait makes it sub-second (§3.1), and
  \texttt{-\/-dry-run} mode plus \texttt{withOpAudit()} snapshots make
  destructive ops (reindex, consolidate, compact, crystallize, kg-prune)
  reversible.
\end{enumerate}

The single design principle that ties these closures together is
\textbf{pain weighting under shadow discipline}: every chunk carries an
explicit \texttt{pain} severity, every ranking change rolls out first in
\texttt{NOX\_SALIENCE\_MODE=shadow} with /api/health telemetry
\footnote{\texttt{NOX\_SALIENCE\_MODE} environment variable controls the
  three-state gate (\texttt{shadow} \textbar{} \texttt{active}
  \textbar{} \texttt{off}). Default is \texttt{shadow}. Telemetry
  exposed at \texttt{/api/health.salience}. See §3.4.3,
  \texttt{staged/1.7a/edits/salience.ts}, and
  \texttt{docs/CONFIGURATION.md}.}, and only graduates to
\texttt{active} after operator review.

\begin{center}\rule{0.5\linewidth}{0.5pt}\end{center}

\subsection{2. System Architecture}\label{system-architecture}

\subsubsection{2.1 Infrastructure
Overview}\label{infrastructure-overview}

The system runs on a Hostinger KVM4 VPS accessible via Tailscale VPN (a
private Tailscale address). Five systemd-managed services provide the
runtime environment:

{\def\LTcaptype{none} % do not increment counter
\begin{longtable}[]{@{}llll@{}}
\toprule\noalign{}
Service & Port & Type & Function \\
\midrule\noalign{}
\endhead
\bottomrule\noalign{}
\endlastfoot
openclaw-gateway & 18789 & WebSocket & Agent communication gateway \\
nox-mem-watcher & --- & inotifywait & Filesystem event monitor \\
nox-mem-api & 18800 & HTTP/JSON & Dashboard data API \\
ollama & 11434 & HTTP & Local LLM inference (llama3.2:3b) \\
tailscaled & --- & WireGuard & VPN mesh connectivity \\
\end{longtable}
}

\subsubsection{2.2 Database Schema}\label{database-schema}

The primary storage is SQLite 3 with WAL (Write-Ahead Logging) mode for
concurrent access. The schema (version 3) contains:

\textbf{Core Tables:}

\begin{itemize}
\tightlist
\item
  \texttt{chunks} --- Memory fragments with full-text indexing

  \begin{itemize}
  \tightlist
  \item
    \texttt{id} (INTEGER PK), \texttt{source\_file} (TEXT),
    \texttt{chunk\_text} (TEXT), \texttt{chunk\_type} (TEXT)
  \item
    \texttt{source\_date} (TEXT), \texttt{is\_consolidated} (INTEGER),
    \texttt{memory\_type} (TEXT)
  \item
    \texttt{created\_at}, \texttt{updated\_at} (TEXT, ISO 8601)
  \item
    \texttt{metadata} (TEXT, JSON)
  \end{itemize}
\item
  \texttt{chunks\_fts} --- FTS5 virtual table with porter unicode61
  tokenizer

  \begin{itemize}
  \tightlist
  \item
    Content-sync triggers (INSERT, UPDATE, DELETE) maintain index
    consistency
  \item
    BM25 ranking with configurable column weights (1.0, 0.5, 0.5)
  \end{itemize}
\item
  \texttt{consolidated\_files} --- Processing state tracker

  \begin{itemize}
  \tightlist
  \item
    \texttt{source\_file} (TEXT PK), \texttt{status} (INTEGER:
    0=pending, 1=done, -1=failed)
  \end{itemize}
\item
  \texttt{meta} --- Key-value configuration store (schema\_version,
  cursors, metrics)
\end{itemize}

\textbf{Knowledge Graph Tables:}

\begin{itemize}
\tightlist
\item
  \texttt{kg\_entities} --- Named entities with type classification and
  mention counting

  \begin{itemize}
  \tightlist
  \item
    UNIQUE constraint on (name, entity\_type)
  \item
    TTL tracking via \texttt{first\_seen}, \texttt{last\_seen}
  \end{itemize}
\item
  \texttt{kg\_relations} --- Typed relationships between entities

  \begin{itemize}
  \tightlist
  \item
    Confidence scoring (0.0-1.0) with temporal decay
  \item
    TTL via \texttt{expires\_at} (90-day default),
    \texttt{last\_confirmed}
  \item
    Evidence linking via \texttt{evidence\_chunk\_id}
  \end{itemize}
\item
  \texttt{decision\_versions} --- Architectural decision version history

  \begin{itemize}
  \tightlist
  \item
    Supersession chain via \texttt{is\_current} flag and
    \texttt{superseded\_at} timestamp
  \end{itemize}
\end{itemize}

\textbf{Vector Tables (sqlite-vec):}

\begin{itemize}
\item
  \texttt{vec\_chunks} --- Virtual table storing float32 embeddings
  (3072 dimensions)
\item
  \texttt{vec\_chunk\_map} --- Rowid-to-chunk\_id mapping (sqlite-vec
  requires rowid-based access)
\item
  \texttt{dedup\_log} --- Suppressed duplicate tracking for audit
\end{itemize}

\subsubsection{2.3 Chunk Type Taxonomy}\label{chunk-type-taxonomy}

Memory chunks are classified into 10 types based on source file path
patterns:

{\def\LTcaptype{none} % do not increment counter
\begin{longtable}[]{@{}
  >{\raggedright\arraybackslash}p{(\linewidth - 6\tabcolsep) * \real{0.1333}}
  >{\raggedright\arraybackslash}p{(\linewidth - 6\tabcolsep) * \real{0.3333}}
  >{\raggedright\arraybackslash}p{(\linewidth - 6\tabcolsep) * \real{0.3333}}
  >{\raggedright\arraybackslash}p{(\linewidth - 6\tabcolsep) * \real{0.2000}}@{}}
\toprule\noalign{}
\begin{minipage}[b]{\linewidth}\raggedright
Type
\end{minipage} & \begin{minipage}[b]{\linewidth}\raggedright
Source Pattern
\end{minipage} & \begin{minipage}[b]{\linewidth}\raggedright
Current Count
\end{minipage} & \begin{minipage}[b]{\linewidth}\raggedright
Purpose
\end{minipage} \\
\midrule\noalign{}
\endhead
\bottomrule\noalign{}
\endlastfoot
team & \texttt{shared/} & 499 & Shared team knowledge, cross-agent
docs \\
daily & \texttt{memory/YYYY-MM-DD} & 161 & Daily operational notes \\
other & (default) & 126 & Unclassified content \\
decision & \texttt{memory/decisions.md} & 34 & Architectural and
strategic decisions \\
lesson & \texttt{memory/lessons.md} & 21 & Errors, corrections,
learnings \\
project & \texttt{memory/projects.md} & 11 & Active project tracking \\
pending & \texttt{memory/pending.md} & 8 & Incomplete tasks and
blockers \\
feedback & \texttt{memory/feedback/} & 6 & User and system feedback \\
person & \texttt{memory/people.md} & 6 & People profiles and contacts \\
digest & \texttt{memory/digests/} & 2 & Weekly summary reports \\
\end{longtable}
}

\subsubsection{2.4 Multi-Agent Memory
Architecture}\label{multi-agent-memory-architecture}

Each of the 6 agents operates with an isolated database at
\texttt{/root/.openclaw/agents/\{name\}/tools/nox-mem/nox-mem.db}. The
\texttt{OPENCLAW\_WORKSPACE} environment variable controls path
resolution across all modules, enabling the same nox-mem binary to
operate on different databases depending on the calling context.

\textbf{Agent Memory Distribution (as of March 23, 2026):}

{\def\LTcaptype{none} % do not increment counter
\begin{longtable}[]{@{}lllll@{}}
\toprule\noalign{}
Agent & Role & Chunks & DB Size & Dominant Type \\
\midrule\noalign{}
\endhead
\bottomrule\noalign{}
\endlastfoot
Nox & Chief of Staff & 185 & 268 KB & daily (91) \\
Boris & Head of Communications & 148 & 268 KB & team (50) \\
Forge & Code Reviewer & 182 & 292 KB & daily (135) \\
Atlas & Research & 30 & 128 KB & other (17) \\
Cipher & Security & 31 & 132 KB & other (12) \\
Lex & Legal/Compliance & 31 & 132 KB & other (12) \\
\textbf{Workspace} & \textbf{Shared} & \textbf{874} & \textbf{25.2 MB} &
\textbf{team (499)} \\
\end{longtable}
}

Total system memory: 1,481 chunks across 7 databases (initial-deployment
snapshot, March 2026; the main store has since grown to
\textasciitilde95k chunks --- see Abstract).

\subsubsection{2.5 User-Facing Primitives}\label{user-facing-primitives}

nox-mem exposes a deliberately small public contract: \textbf{three
primitives}, all backed by the same SQLite store and surfaced
identically across three transport layers (CLI, HTTP API, MCP). Every
advanced verb (\texttt{reflect}, \texttt{cross-search},
\texttt{kg-path}, \texttt{crystallize}) decomposes internally into
sequences of these three primitives.

\textbf{Primitive 1 --- \texttt{search} (hybrid retrieval).} FTS5 BM25 +
Gemini semantic (3072d) → Reciprocal Rank Fusion (RRF), k=60, with
optional Hard Mutex section gating and SOURCE\_TYPE\_BOOST overlays.
Returns ranked chunks with \texttt{score}, \texttt{match\_type}, and
provenance fields (\texttt{source\_file}, \texttt{section},
\texttt{created\_at}, \texttt{updated\_at}). Detailed in §4.

\textbf{Primitive 2 --- \texttt{answer} (grounded RAG).} Internally
calls \texttt{search} with \texttt{topK\ =\ 10}, builds a
citation-anchored prompt over the retrieved chunks, invokes the
configured LLM (\texttt{gemini-2.5-flash-lite} by default per D41), and
parses inline \texttt{{[}chunk\_\textless{}id\textgreater{}{]}}
citations. Anti-hallucination guard: citations pointing to chunks
outside the retrieved set trigger a single retry with a stricter prompt;
a second failure raises
\texttt{AnswerError(\textquotesingle{}hallucination\_after\_retry\textquotesingle{})}.
Empty-retrieval short-circuit avoids LLM spend when no chunks match.
Measured p95 latency: 101.74 ms on the offline mock-LLM bench (PR \#40,
42× under the 4.3 s budget); live p95 with Gemini Flash Lite ranges
1.5--2.5 s. Implementation:
\texttt{staged/P1/edits/src/lib/answer/\{index,retrieval,prompt,provider,config\}.ts}.

\textbf{Primitive 3 --- Temporal filter (\texttt{-\/-as-of} /
\texttt{-\/-changed-since}).} Time-travel and recency-window selectors
implemented as hard SQL pre-filters --- not ranking boosts.
\texttt{-\/-as-of\ \textless{}date\textgreater{}} restricts to chunks
satisfying
\texttt{created\_at\ \textless{}=\ date\ AND\ (deleted\_at\ IS\ NULL\ OR\ deleted\_at\ \textgreater{}\ date)};
\texttt{-\/-changed-since\ \textless{}date\textgreater{}} restricts to
chunks satisfying
\texttt{updated\_at\ \textgreater{}\ date\ OR\ created\_at\ \textgreater{}\ date}.
Combined, the two clauses AND. Accepted formats: ISO 8601
(\texttt{2026-05-01} or full \texttt{2026-05-01T00:00:00Z}) and relative
(\texttt{7d}, \texttt{1w}, \texttt{30d}, \texttt{2h}, \texttt{15m}).
Uses existing \texttt{chunks.created\_at} and
\texttt{chunks.updated\_at} columns from schema v18 --- no schema
changes, no ranking changes. The filter is orthogonal to the E13
temporal proximity boost (\texttt{NOX\_TEMPORAL\_PATH}, §5), which
additively reweights ranking by recency rather than restricting the
candidate set. Implementation:
\texttt{staged/P3/edits/\{dates,search,api-server\}.ts}.

\textbf{Composition.} The three primitives compose orthogonally --- for
example,
\texttt{answer\ "what\ incidents\ happened\ last\ week?"\ -\/-changed-since\ 7d}
retrieves only chunks updated in the last seven days, then synthesizes a
grounded answer over that restricted candidate set. This closure
property --- three small primitives, deterministic semantics, identical
surface across transports --- is the contract that makes nox-mem
composable from agent runtimes that have no prior knowledge of the
implementation.

\textbf{Tagline.} \emph{3 primitives, 1 file, any LLM.} The three
primitives are search + answer + temporal filter; the one file is the
SQLite database on the operator's disk; the LLM provider is swappable
via the \texttt{LLMProvider} interface (Gemini default, OpenAI /
Anthropic / Ollama / vLLM available) without code changes. Full operator
reference: \texttt{docs/PRIMITIVES.md}.

\begin{center}\rule{0.5\linewidth}{0.5pt}\end{center}

\subsection{3. Memory Pipeline}\label{memory-pipeline}

\subsubsection{3.1 Ingestion}\label{ingestion}

Files created or modified in monitored directories trigger the
inotifywait-based watcher service. The watcher implements:

\begin{itemize}
\tightlist
\item
  \textbf{Debounce logic}: 2-second delay to batch rapid successive
  writes
\item
  \textbf{File filtering}: Only \texttt{.md} and \texttt{.json} files
  are processed
\item
  \textbf{Recursion prevention}: \texttt{MEMORY.md} and
  \texttt{SESSION-STATE.md} are excluded to avoid feedback loops
\item
  \textbf{Heartbeat}: Touch \texttt{/tmp/nox-mem-watcher-heartbeat} on
  every event for liveness monitoring
\end{itemize}

Upon trigger, \texttt{ingestFile()} executes:

\begin{enumerate}
\def\labelenumi{\arabic{enumi}.}
\tightlist
\item
  Read file content with UTF-8 sanitization (fixes common mojibake
  patterns for Portuguese text)
\item
  Detect chunk type from relative file path
\item
  Extract date from filename pattern (YYYY-MM-DD)
\item
  Split content into semantic chunks:

  \begin{itemize}
  \tightlist
  \item
    Markdown: Split on H2/H3 headers, with sub-splitting for chunks
    exceeding 500 words
  \item
    JSON: Array items become individual chunks; object entries become
    key-value pairs
  \item
    Small chunks (\textless20 words) are merged with the previous chunk
  \end{itemize}
\item
  Delete existing chunks for the same source file (idempotent
  re-ingestion)
\item
  Insert new chunks via prepared statement transaction
\item
  Auto-vectorize if GEMINI\_API\_KEY is available (up to 20 chunks per
  file)
\end{enumerate}

\subsubsection{3.2 Consolidation}\label{consolidation}

Nightly consolidation (23:00, 5-minute stagger across agents) processes
daily notes into structured topic files:

\begin{enumerate}
\def\labelenumi{\arabic{enumi}.}
\tightlist
\item
  \textbf{Reindex}: Scan all \texttt{.md}/\texttt{.json} files in memory
  directories, rebuild chunk index
\item
  \textbf{Extract}: Use Ollama llama3.2:3b to identify facts, decisions,
  lessons, and action items from daily notes
\item
  \textbf{Append}: Add extracted content to topic files (decisions.md,
  lessons.md, people.md, projects.md, pending.md)
\item
  \textbf{Notion Sync}: Push structured items to ``Memoria \& Decisoes''
  Notion database (best-effort, non-blocking)
\item
  \textbf{Git Commit}: Auto-commit memory changes with standardized
  message format
\item
  \textbf{Session Update}: Refresh SESSION-STATE.md with current
  statistics
\end{enumerate}

\subsubsection{3.3 Deduplication}\label{deduplication}

Before insertion, chunks are checked for duplicates using a two-tier
strategy:

\begin{itemize}
\tightlist
\item
  \textbf{Primary}: Gemini cosine similarity with 0.85 threshold (when
  embeddings are available)
\item
  \textbf{Fallback}: Keyword overlap calculation with 60\% threshold
\item
  \textbf{Audit}: Suppressed duplicates are logged to
  \texttt{dedup\_log} table with reason and preview
\end{itemize}

\subsubsection{3.4 Self-Evolution: How nox-mem Improves Over
Time}\label{self-evolution-how-nox-mem-improves-over-time}

Most memory systems decouple \emph{retrieval} from \emph{learning}:
retrieval is online, learning is either absent or requires retraining a
parametric backbone (MeMo) or relying on a closed evolution loop (EverOS
EvoAgentBench \footnote{EverMind-AI / EverOS --- github.com/EverMind-AI
  (\textasciitilde5k stars, Apache 2.0). Publishes EverMemBench dataset
  and an EvoAgentBench-framed evolution loop. The only memory-OS peer in
  our taxonomy that publishes its own benchmark; §3.4 is the direct
  narrative counter to EvoAgent framing. Honest-comparison action item:
  run nox-mem on EverMemBench (queued as F4 in §7.2).}). nox-mem instead
promotes self-evolution to a first-class subsystem composed of four
cooperating mechanisms, all transparent and inspectable in the live
SQLite file. This subsection is the direct counter to Gap \#5 (no
writeback) in Table 1.

\textbf{3.4.1 Crystallize loop --- \texttt{pending} → \texttt{lesson}
after N hits, with pain accumulation.}

The \texttt{crystallize} command (exposed via the CLI, the Model Context
Protocol (MCP) server tool, and \texttt{POST\ /api/crystallize} +
\texttt{POST\ /api/crystallize/validate} on the HTTP API; see
\texttt{src/api-server.ts} and \texttt{src/crystallize.ts} \footnote{HTTP
  handlers in \texttt{src/api-server.ts} (routes
  \texttt{/api/crystallize}, \texttt{/api/crystallize/validate} ---
  confirmed in \texttt{staged/1.6/edits/api-server.ts:253–273}); core
  logic in \texttt{src/crystallize.ts} exporting \texttt{crystallize()},
  \texttt{validateProcedure()}, \texttt{listProcedures()}. Wrapped by
  \texttt{withOpAudit()} (\texttt{src/lib/op-audit.ts}).}) implements an
LLM-assisted consolidation that synthesizes durable entities from recent
chunks. Operationally, chunks ingested into \texttt{memory/pending.md}
(\texttt{chunk\_type\ =\ pending}, \texttt{retention\_days\ =\ 30}) act
as a working set of unresolved items. As the same pending item is
referenced by retrieval (and as related daily/team chunks accumulate
around it), its effective \emph{pain} signal grows --- both via direct
operator updates to the \texttt{pain} column (V9 schema) and via
co-occurrence with high-pain neighbors. After enough hits and sufficient
accumulated severity, \texttt{crystallize} uses Gemini to identify the
pattern across chunks, emits a canonical entity file under
\texttt{memory/entities/\textless{}type\textgreater{}/\textless{}slug\textgreater{}.md},
and effectively promotes the cluster from \texttt{pending} (30-day
decay) to \texttt{lesson} (180-day decay) or
\texttt{decision}/\texttt{project} (365-day decay) \footnote{V8 schema
  typed retention defaults (in \texttt{chunks.retention\_days}):
  \texttt{feedback} = 0 (never-decay), \texttt{person} = 0
  (never-decay), \texttt{lesson} = 180d, \texttt{decision} = 365d,
  \texttt{project} = 365d, \texttt{team} = 120d, \texttt{daily} = 90d,
  \texttt{pending} = 30d, \texttt{graph\_node} = 60d, fallback = 90d.
  See \texttt{staged/1.7a/edits/salience.ts:46–56}
  (\texttt{DEFAULT\_RETENTION\_BY\_TYPE}) and \texttt{CLAUDE.md}
  §``Schema v10''.}. The promotion is wrapped in \texttt{withOpAudit()},
so the pre-state snapshot lives at
\texttt{/var/backups/nox-mem/pre-op/crystallize-\textless{}ts\textgreater{}-\textless{}pid\textgreater{}-\textless{}uuid\textgreater{}.db}
and a row lands in the append-only \texttt{ops\_audit} table.

\textbf{3.4.2 Pain weighting auto-adjust --- feedback of use raises pain
on hit chunks.}

The \texttt{pain} field on \texttt{chunks} is an explicit severity in
{[}0.1, 1.0{]} (0.1 trivial, 1.0 production outage). It can be set
explicitly by the author, but it also drifts upward implicitly: chunks
that are \emph{repeatedly} surfaced by retrieval and \emph{accepted} by
the operator (via the \texttt{feedback/} directory, whose
\texttt{chunk\_type\ =\ feedback} is never-decay) accumulate
co-occurrence with high-pain entries during nightly consolidation, which
feeds back into the salience formula. The result is that lessons learned
from real incidents naturally rise to the top over time, while toy or
low-severity content fades even when frequent. This is the inverse of
pure recency- or frequency-based memory.

\textbf{3.4.3 Salience decay continuous in background ---
\texttt{recency\ ×\ pain\ ×\ importance}.}

The salience function lives in \texttt{src/lib/salience.ts} \footnote{\texttt{src/lib/salience.ts}
  --- canonical implementation of the salience formula (multiplicative
  principle in §3.4.3; weighted-additive v2 production form
  \texttt{W\_IMPORTANCE·importance\ +\ W\_RECENCY·recency\ +\ W\_PAIN·pain\ +\ W\_ACCESS·access\_score}
  in §5.1), mirrored in the staged copy at
  \texttt{staged/1.7a/edits/salience.ts} (lines 1--80 contain the module
  docstring spelling out the formula, the three-state mode gate, and the
  per-type retention defaults).} (and is mirrored in the staged copy at
\texttt{staged/1.7a/edits/salience.ts}). Its underlying principle is
multiplicative --- a memory scores high only when it is
\emph{simultaneously} recent, painful, and important:

\begin{verbatim}
salience = recency × pain × importance
\end{verbatim}

with components:

\begin{itemize}
\tightlist
\item
  \textbf{recency} in {[}0,1{]} --- half-life-style decay over the
  chunk's \texttt{retention\_days} window
  (\texttt{feedback}/\texttt{person} → never-decay →
  \texttt{recency\ =\ 1.0}; everything else decays per Table V8
  \footnote{V8 schema typed retention defaults (in
    \texttt{chunks.retention\_days}): \texttt{feedback} = 0
    (never-decay), \texttt{person} = 0 (never-decay), \texttt{lesson} =
    180d, \texttt{decision} = 365d, \texttt{project} = 365d,
    \texttt{team} = 120d, \texttt{daily} = 90d, \texttt{pending} = 30d,
    \texttt{graph\_node} = 60d, fallback = 90d. See
    \texttt{staged/1.7a/edits/salience.ts:46–56}
    (\texttt{DEFAULT\_RETENTION\_BY\_TYPE}) and \texttt{CLAUDE.md}
    §``Schema v10''.}). Decay is \emph{continuous}: it is recomputed at
  every retrieval, so there is no batch step that ``ages'' memory ---
  aging is a property of the read path.
\item
  \textbf{pain} in {[}0,1{]} --- severity as described in §3.4.2.
\item
  \textbf{importance} in {[}0,1{]} --- \texttt{chunk\_type} /
  \texttt{source\_type} / tier signal (manual mapping; e.g.,
  \texttt{decision} and \texttt{lesson} rank above \texttt{daily}).
\end{itemize}

In production since Wave A (§5.1), this principle is realized as a
\textbf{weighted-additive} form ---
\texttt{salience\ =\ W\_IMPORTANCE·importance\ +\ W\_RECENCY·recency\ +\ W\_PAIN·pain\ +\ W\_ACCESS·access\_score}
(weights 0.55 / 0.15 / 0.10 / 0.20) --- which empirically outperformed
the strict multiplicative product on the Wave A golden set; §5.1 reports
the ablation. The multiplicative expression above states the principle;
the additive form is what ships.

The mode of operation is gated by \texttt{NOX\_SALIENCE\_MODE}:
\texttt{shadow} (default --- compute and log to
\texttt{/api/health.salience} but do not apply to retrieval rankings),
\texttt{active} (apply as an additive delta in {[}-0.5, +0.5{]} on top
of RRF), and \texttt{off} (short-circuit to 0 for ablation experiments).
This three-state gate is the operational realization of \emph{shadow
discipline} (§4 and §5).

\textbf{3.4.4 Reflect pipeline --- batch nightly cross-session
synthesis.}

The \texttt{reflect} command (CLI / MCP / \texttt{POST\ /api/reflect},
backed by \texttt{src/reflect.ts} and cached via
\texttt{getReflectCacheStats()} exposed in
\texttt{/api/health.reflectCache} \footnote{\texttt{src/reflect.ts}
  exporting \texttt{reflect()} and \texttt{getReflectCacheStats()}
  (confirmed in \texttt{staged/1.6/edits/api-server.ts:12}). Cache
  statistics surfaced in \texttt{/api/health.reflectCache}. See also
  \texttt{docs/POSTMAN.md} for the API contract.}) runs a batch
synthesis pass over recent chunks. Its job is \emph{not} to retrieve
answers for the user but to produce cross-session lessons: it queries
hybrid search for a topic, asks Gemini to summarize the patterns
observed across the result set, and writes the summary back as a
\texttt{feedback} or \texttt{lesson} chunk with the confidence default
for derived chunks (see \texttt{CONFIDENCE\_DERIVED} in
\texttt{docs/CONFIGURATION.md}). Because reflect results are themselves
cached (LRU, exposed via \texttt{/api/health.reflectCache.entries}),
expensive synthesis is amortized across repeated queries.

\textbf{3.4.5 Consolidate --- nightly orchestration that ties it
together.}

Nightly consolidation (23:00, 5-minute stagger across agents --- see
§3.2) is the orchestration layer that runs reflect for high-pain topics,
runs crystallize for sufficiently aged \texttt{pending} items, and syncs
the resulting durable entities to the topic files
(\texttt{decisions.md}, \texttt{lessons.md}, \texttt{people.md},
\texttt{projects.md}). Like crystallize, it goes through
\texttt{withOpAudit()} with a pre-op \texttt{VACUUM\ INTO} snapshot.
This is the daily heartbeat of self-evolution.

\textbf{Together, these four mechanisms make nox-mem a memory that
\emph{grows along the gradient of operator pain}:} chunks that hurt the
operator (production outages, lost work, repeated mistakes) survive
decay, get promoted from pending to lesson, accumulate co-occurrence
weight, and rank progressively higher --- without retraining a model and
without trusting a closed evolution loop. Every step is visible by
opening \texttt{nox-mem.db} in \texttt{sqlite3} and inspecting
\texttt{ops\_audit}, \texttt{chunks.pain},
\texttt{chunks.retention\_days}, and the entity files.

\subsubsection{3.5 Session Priming: the read-side of
self-evolution}\label{session-priming-the-read-side-of-self-evolution}

The mechanisms above are \emph{write-side} --- they decide what the
store keeps and how it ranks. But a memory that only grows is wasted if
every new session starts blind and must re-discover context through cold
search. \textbf{Session priming} closes the loop on the
\emph{read-side}: each session (any agent, any MCP/HTTP client) is born
contextualized. \texttt{GET\ /api/brief} (handler in
\texttt{src/api/brief.ts}) returns a salience-ranked, scope-filtered
digest of the top-N chunks, injected at session start. The brief is
strictly read-only over \texttt{chunks} --- it never touches
\texttt{access\_count} or \texttt{last\_accessed\_at}, so the
organic-use signal that feeds salience (§3.4) stays uncontaminated;
serve-history is tracked in a separate \texttt{brief\_log} table. With
an \texttt{agent} scope the digest is a guaranteed union
(\textasciitilde half agent-specific, half global high-salience), and
near-duplicate variants are collapsed by token-containment so the budget
is spent on distinct content.

\textbf{An honest measurement, and a methodology caveat.} Running
priming in production (6 agents, \textasciitilde6.7k serves/day)
surfaced a failure mode worth reporting. Over a clean 7-day window the
brief served only \textbf{83 distinct chunks across 46.8k serves (0.18\%
diversity)}, median content age 48 days: salience, dominated by
importance and the accumulated pain of old incidents, produces a
near-static top set, and recently-written content (the very output the
loop produces) rarely enters --- we measured \textbf{931 recent
(\textless=7d), relevant chunks (importance \textgreater= 0.7 or pain
\textgreater= 0.7) that were never once served}. More instructive was a
\emph{negative} result on evaluation: we could not measure a downstream
``follow-up rate'' (does a primed chunk get used later?) because the
signal does not exist --- in 7 days there were \textbf{3 genuine
searches and 0 \texttt{answer} calls}. This is structural, not a logging
gap: priming-by-injection delivers the knowledge \emph{into the context
window}, so the agent does not re-query what it already holds.
\textbf{The utility of an injection-based primer is therefore a property
of the served set (diversity, freshness, coverage of what matters) ---
not of downstream re-retrieval.} Memory work that imports a re-query
metric from RAG-style retrieval will mis-evaluate this class of
mechanism.

\textbf{Diversity term (D2).} The fix follows the same discipline as
§3.4.3: it does \textbf{not} fork \texttt{calculateSalience} (that would
alter search too) --- it is a re-rank \emph{inside the brief only}. (A)
a novelty penalty
\texttt{brief\_score\ =\ salience\ -\ min(P\_max,\ lambda*log1p(n\_serves))}
reads \texttt{brief\_log} to demote over-served chunks, saturating in
log so serving 2,000x \textasciitilde= serving 100x; (B) a freshness
quota reserves F of the N slots for recent, relevant, not-yet-served
content (a separate query --- the candidate pool, ranked by importance
and access, structurally excludes \texttt{access\_count\ =\ 0}
newcomers). A \textbf{high-pain floor} makes incidents with
\texttt{pain\ \textgreater{}=\ 0.9} immune to both demotion and
displacement: diversity must not hide the critical. Per the
shadow-discipline rule (§3.4.3, §5), the term ships behind
\texttt{NOX\_BRIEF\_DIVERSITY\ =\ off\ \textbar{}\ shadow\ \textbar{}\ active},
default off and bit-identical to the prior behavior; it is validated in
shadow (logging the would-enter/would-leave diff against a per-day gate:
diversity up, median age down, floor preserved, churn non-thrashing)
before any flip to active. Notably, the shadow gate itself caught a
design bug on its first run --- the freshness quota was displacing
\texttt{pain\ =\ 1.0} incidents --- which the floor now prevents; the
active-mode quality numbers are reported as ongoing validation (§7.2).

\textbf{Active-mode validation surfaced a second, subtler failure ---
and the fix unifies the two slots.} Two refinements followed the first
gate. First, the freshness query as scoped (\texttt{global+agent}) only
ever saw \texttt{sessions/\textless{}agent\textgreater{}/\%}, never the
curated \texttt{memory/entities/\%} store, so freshly-consolidated
lessons and decisions were structurally invisible to every brief; a
\emph{split-slot} interleave (one agent-recent plus one curated-global
per brief, with its own age window) restored that channel. Second ---
and this is the instructive part --- once the curated channel was live
in \texttt{active}, a 24-hour gate showed the global slot serving
\textbf{190 distinct curated chunks on day one, then exactly 1 per day
thereafter}. The cause was the slot's \emph{hard} anti-repeat
(\texttt{id\ NOT\ IN\ brief\_log} over a 72h window): under production
volume (\textasciitilde6.7k briefs/day) over a 189-chunk eligible pool,
the entire pool is served --- and thus excluded --- within a single day,
after which the slot dries up and fails open to the static salience
top-1 for the remaining 72h. The first remedy \emph{seemed} obvious:
replace hard exclusion with the soft novelty penalty already used by the
primary pool (mechanism A) --- serving a curated chunk \emph{demotes} it
by \texttt{lambda*log1p(n\_serves)} (capped at \texttt{P\_max\ =\ 0.15})
rather than \emph{removing} it, with the high-pain floor keeping
critical incidents immune. Active-mode measurement refuted it: the
24-hour gate fell \textbf{146 → 67 → 3 distinct curated chunks per day}
over three successive days. The cap was the flaw --- \texttt{P\_max}
(0.15) is smaller than the base-salience gap between the pool's few
high-salience outliers (decisions at importance 0.9 with high access
counts) and its homogeneous body, so once the penalty saturated (within
one 72h window under production volume) the pick reconverged to the
static salience top-\emph{k}. A bounded soft penalty cannot out-vote an
unbounded salience gap. The fix that held ranks the slot not by
\emph{count} of prior serves but by \emph{time since last serve}
(mechanism B', coverage-sampling): never-served first, then
least-recently-served, tie-broken by salience. Having no ceiling, it
sweeps the entire pool continuously regardless of the salience gap. The
first full post-deploy 24-hour active gate (2026-06-27,
\textasciitilde6.8k briefs/day) confirmed it: \textbf{184 of 184
eligible entity files served (100\% pool coverage), 190 distinct curated
chunks, \textasciitilde45 distinct chunks rotated per hour ---
sustained, not a day-one burst} (the cumulative distinct count reached
190 within four hours of deploy and held flat through hour 24; the pool
is swept fast and then re-cycled by recency rather than exhausted). The
general lesson --- \textbf{hard de-duplication under volume
\textgreater\textgreater{} pool produces bursty rotation; a
\emph{saturating} soft penalty silently reconverges once its cap falls
below the salience gap; only coverage-by-recency-of-serve, which has no
ceiling, produces true continuous rotation} --- is the kind of finding
only a shadow→active→measure loop on real traffic exposes. An offline
diversity metric on a static corpus would have scored all three designs
identically, and the saturating-penalty regression in particular would
have been invisible without a multi-day active gate.

\begin{center}\rule{0.5\linewidth}{0.5pt}\end{center}

\subsection{4. Hybrid Search System}\label{hybrid-search-system}

\subsubsection{4.1 Architecture}\label{architecture}

Search combines three complementary retrieval methods:

\textbf{Layer 1 --- FTS5 BM25 (Keyword)}

SQLite FTS5 with porter unicode61 tokenizer provides fast keyword
matching. Results are scored using BM25 with column weights
(chunk\_text: 1.0, source\_file: 0.5, chunk\_type: 0.5). Post-retrieval
boosting applies:

\begin{itemize}
\tightlist
\item
  Type boost: \texttt{decision} and \texttt{lesson} chunks receive 2.0x
  multiplier (higher signal-to-noise ratio)
\item
  Recency boost: Chunks from the last 7 days receive 1.5x multiplier
\end{itemize}

The query sanitizer strips special characters but preserves hyphens for
compound terms (e.g., ``nox-mem'').

\textbf{Layer 2 --- Gemini Semantic (Vector)}

Each chunk is embedded using Google's gemini-embedding-001 model (3072
dimensions) with task type RETRIEVAL\_DOCUMENT. Query embeddings use
task type RETRIEVAL\_QUERY for asymmetric similarity optimization.

Vectors are stored in sqlite-vec virtual tables. Retrieval uses cosine
distance with a map table (vec\_chunk\_map) bridging vec\_chunks rowids
to chunks.id values due to sqlite-vec's rowid-only constraint.

Scoring normalizes distances to a 0-10 scale with type and recency
boosting (1.5x and 1.2x respectively, lower than FTS5 to avoid
double-boosting in fusion).

\textbf{Layer 3 --- Reciprocal Rank Fusion (RRF)}

FTS5 and semantic results are merged using RRF with k=60:

\begin{verbatim}
RRF_score(d) = Σ 1/(k + rank_i(d))
\end{verbatim}

Documents appearing in both result sets receive combined scores, marked
as \texttt{match\_type:\ "hybrid"}. Content-prefix deduplication (first
50 characters) prevents near-duplicate results.

\subsubsection{4.2 Performance
Characteristics}\label{performance-characteristics}

The hybrid approach provides significant quality improvements over
single-method search:

{\def\LTcaptype{none} % do not increment counter
\begin{longtable}[]{@{}
  >{\raggedright\arraybackslash}p{(\linewidth - 6\tabcolsep) * \real{0.1944}}
  >{\raggedright\arraybackslash}p{(\linewidth - 6\tabcolsep) * \real{0.3056}}
  >{\raggedright\arraybackslash}p{(\linewidth - 6\tabcolsep) * \real{0.2222}}
  >{\raggedright\arraybackslash}p{(\linewidth - 6\tabcolsep) * \real{0.2778}}@{}}
\toprule\noalign{}
\begin{minipage}[b]{\linewidth}\raggedright
Query
\end{minipage} & \begin{minipage}[b]{\linewidth}\raggedright
FTS5 Only
\end{minipage} & \begin{minipage}[b]{\linewidth}\raggedright
Hybrid
\end{minipage} & \begin{minipage}[b]{\linewidth}\raggedright
Analysis
\end{minipage} \\
\midrule\noalign{}
\endhead
\bottomrule\noalign{}
\endlastfoot
``qual o proximo passo'' & 0 results & ROADMAP + PHASE-3 & Semantic
captures intent without keyword match \\
``nox-mem'' & 0 results & decisions.md + docs & Vector bypasses
tokenizer hyphen issues \\
``quem e o Toto'' & people.md & people.md + TEAM\_MEMORY & RRF combines
exact match + semantic context \\
\end{longtable}
}

\subsubsection{4.3 Cross-Agent Search}\label{cross-agent-search}

The \texttt{crossSearch()} function opens all 7 databases in read-only
mode, executes FTS5 queries in each, and merges results with agent
attribution. Deduplication uses content-prefix comparison to handle
shared documents that appear across multiple agent databases.

\begin{center}\rule{0.5\linewidth}{0.5pt}\end{center}

\subsection{5. Empirical Evaluation (May--June 2026, fifth
revision)}\label{empirical-evaluation-mayjune-2026-fifth-revision}

\begin{quote}
\textbf{Headlines (2026-05-29/06-02, 5-batch canonical, four-benchmark
triangulation + Q3 IterB ceiling break + Wave 2 backbone-conditional
knob study):} - \textbf{Q3 IterB ReAct --- F\_MH retrieval-stage ceiling
broken (D74, PR \#419):} on Gemini-3-flash bare baseline, multi-round
retrieve-reason loop delivers \textbf{+2.01 pp clean F\_MH lift (6.02\%
→ 8.03\%)} --- exceeding the Wave A/B/C single-stage retrieval ceiling
of 7.25 pp by +0.78 pp standalone, with the strongest backbone in the
matrix (§5.5.2, dual-baseline reporting). - \textbf{Wave 2
backbone-conditional knob study (NEW, D75, PRs \#423--\#425):} Wave A
single-stage retrieval knobs (KG path, AC, MQ) transfer at only
\textbf{\textasciitilde24--40\% efficiency} from gpt-4.1-mini to
Gemini-3-flash. All three fail the \textgreater=+1.5 pp gate on
Gemini-3-flash (NO-REPLICATE). 3-knob sum = +2.01 pp = 24\% of original
D74 projection +8.43 pp.~\textbf{IterB ReAct remains the only validated
F\_MH lever on Gemini-3-flash} (§5.5.5). - \textbf{Wave 2 architectural
lock discovery (NEW, D75):} IterB ReAct adapter deliberately
short-circuits Wave A knobs via explicit
\texttt{if\ not\ iterb\_used\_path:} guards in
\texttt{adapter\_nox\_mem.py}. Composability requires explicit code
patch, not merely env-var toggling. Composability projections must
address both backbone-portability and architectural composability
(§5.5.7). - \textbf{Wave 2 MQ MA backbone flip sub-finding (NEW, PR
\#425):} MQ F\_MH lift attenuates gpt-4.1-mini→Gemini (34\% transfer
rate), while MA composite \textbf{flips} from −1.38 pp regression on
gpt-4.1-mini to +0.12 pp preserved on Gemini-3-flash. Multi-axis
backbone-conditional behavior (§5.5.6). - \textbf{Wave 2 Capstone
INDETERMINATE --- infrastructure constraint, NOT scientific failure
(D76, PR \#426 draft):} IterB+KG+rerank composability test aborted after
Hostinger VPS CPU steal 51--97\% sustained. Batch 004 (n=49) preserved;
5-batch statistical threshold not reached. Capstone deferred to
dedicated CPU infrastructure (§5.5.8). - \textbf{EverMemBench Overall
SOTA (Gemini-3-flash, Backbone Matrix):} nox-mem \textbf{63.28\%} vs
MemOS 42.55\% (Table 4 baseline) = \textbf{+20.73 pp WIN} (PR \#397). -
\textbf{EverMemBench Memory Awareness composite SOTA (Gemini-3-flash):}
nox-mem \textbf{88.42\%} vs MemOS 55.68\% = \textbf{+32.74 pp WIN} (PR
\#397). - \textbf{MuSiQue SOTA (multi-hop QA, classical):} dev F1
\textbf{58.62\%} = \textbf{+22.82 pp vs IRCoT 35.80\%} and \textbf{+8.92
pp vs EX(SA) 49.70\%} (PR \#407). - \textbf{HotPotQA SOTA (multi-hop QA,
classical):} dev distractor ans\_F1 \textbf{73.37\%} above DPR+FiD
reader SOTA range 65--72\% (PR \#408). - \textbf{LoCoMo retrieval@10
SOTA (cross-bench memory):} strict \textbf{74.52\%} above Mem0 SOTA F1
66.88\% (PR \#396); F1 push 51.85\% (rank-5, above Zep/LangMem) due to
verbosity gap (PR \#404). - \textbf{Phase D (Gemini-2.5-flash):} nox-mem
\textbf{62.22\%} vs MemOS 59.27\% = \textbf{+2.95 pp WIN} (PR \#365). -
\textbf{Phase H v2 (GPT-4.1-mini cross-backbone):} nox-mem
\textbf{51.68\%} vs MemOS 42.55\% = \textbf{+9.13 pp WIN} (95\% CI
{[}49.88, 53.49{]}, PR \#377). - \textbf{Backbone portability:} nox-mem
regresses −10.54 pp on Gemini-2.5-flash → GPT-4.1-mini swap vs MemOS
−16.72 pp = \textbf{1.6× more portable} (§5.8). - \textbf{Wave B/C
composability (D68 + D69):} KG+MAP additive on F\_MH +4.04 pp
(different-stage compose); same-stage retrieval knobs cap at
\textasciitilde7.25 pp F\_MH (D69 Wave C ceiling, PRs \#393, \#399). -
\textbf{EverMemBench F\_MH paradox REFINED:} MuSiQue 58.62\% F1 and
HotPotQA 73.37\% ans\_F1 prove multi-hop reasoning is SOTA; EverMemBench
F\_MH 3--7\% remains largely structural (very long conversation chains +
strict scoring), but MAS orchestration (Q3 IterB ReAct, this revision)
adds \textbf{+2 pp on top of strongest backbone} above the retrieval
ceiling (§5.4, §5.5). - \textbf{Operational (self-hosted):} KG path p50
\textbf{2.9 ms}, \textbf{\$0/query}, \textbf{399 MB RSS} self-hosted
single-process (§5.7, PR \#403). - \textbf{Methodology:} all claims use
the \textbf{5-batch + 95\% CI canonical protocol} (PR \#371 + PR \#376).
Single-batch overstates effects 3--6×. - \textbf{Dual-baseline honest
reporting (D74):} orchestration-mechanism claims report against both (a)
project-convention baseline (Phase H v2 GPT-4.1-mini, conflates backbone
+ mechanism) and (b) the strongest in-matrix bare baseline
(Gemini-3-flash, isolates mechanism). Ceiling-break claims load on (b).
\end{quote}

This section documents three evaluation tracks that triangulate the same
architectural claims from complementary directions: the \textbf{Wave A
ablation series} (§5.1.1--§5.1.4, entity-flavored golden set, nDCG@10)
exercising the V10 schema's section/source-type/salience drivers; the
\textbf{EverMemBench cross-system series} (§5.1.5--§5.1.10, n=3,121
queries, task-accuracy vs MemOS Table 4 baselines) covering Phase D / H
v2 / G / Lab Q1 standalone knobs / Wave B/C composability / Backbone
Matrix; and the \textbf{classical multi-hop QA series} (§5.2, MuSiQue +
HotPotQA) confirming that multi-hop reasoning is SOTA on standard
benchmarks where corpus structure and scoring are not adversarial. §5.3
cross-validates on LoCoMo (memory-bench, conversational), §5.4 resolves
the EverMemBench F\_MH paradox using the §5.2 and §5.3 evidence, §5.5
reports Q3 orchestration mechanism-class findings --- including the
\textbf{Q3 IterB ReAct ceiling break} (§5.5.2), the \textbf{Wave 2
backbone-conditional knob portability study} (§5.5.5--§5.5.6), the
\textbf{architectural composability lock discovery} (§5.5.7), and the
\textbf{Wave 2 capstone deferral} (§5.5.8) --- §5.6 cross-validates on
LongMemEval, §5.7 reports operational characteristics (latency / cost /
footprint), §5.8 documents the methodology and honest limitations.

\textbf{Dual-baseline reporting convention (D74, 2026-05-31).}
Orchestration-mechanism evaluations in this revision report against
\textbf{two} baselines: (a) the project-convention Phase H v2
GPT-4.1-mini baseline, which preserves comparability with prior
revisions but \textbf{conflates mechanism lift with any backbone swap};
and (b) the strongest in-matrix bare baseline (Gemini-3-flash, §5.1.10),
which \textbf{isolates the mechanism's clean effect} on the best
available reasoning substrate. Any claim of breaking the retrieval-stage
F\_MH ceiling load-bears on the (b) clean number --- reporting only (a)
would conflate ReAct mechanism lift with the +20.73 pp Overall and
+32.74 pp MA composite lifts that the Gemini-3-flash backbone already
provides standalone (§5.1.10). The convention applies forward to §5.5 Q3
IterB (this revision) and any future MAS-class orchestration
evaluations.

\begin{center}\rule{0.5\linewidth}{0.5pt}\end{center}

\subsubsection{5.1 EverMemBench + Wave A ablation
series}\label{evermembench-wave-a-ablation-series}

\paragraph{5.1.1 Wave A ablation --- setup and
headline}\label{wave-a-ablation-setup-and-headline}

The Wave A evaluation uses an entity-flavored golden set of 100 queries
(\texttt{entity-eval.db}), curated from production usage to exercise the
V10 schema's \texttt{section} and \texttt{pain} dimensions.
Configurations are toggled via environment-variable feature gates
(\texttt{NOX\_SALIENCE\_MODE}, \texttt{NOX\_DISABLE\_TIER\_BOOST},
\texttt{NOX\_ENABLE\_TIER\_BOOST},
\texttt{NOX\_DISABLE\_SECTION\_BOOST}, etc.), allowing isolation of
individual ranking components without code changes between runs. All
measurements occur post-deployment of PRs \#150 (salience formula +
tier\_boost off-by-default), \#151 (source\_type backfill of 67,949
chunks), and \#153 (search wiring). Reported nDCG@10 follows the
standard TREC formulation (gain by relevance, log-position discount).

\begin{quote}
\textbf{Wave A headline (canonical, 2026-05-19):} A8 full stack with
active salience reaches \textbf{nDCG@10 = 0.6237} on the entity-flavored
golden set (n=100), a \textbf{+78.8\% relative improvement over the G3
baseline (0.3488)} measured prior to Wave A deployment, and
\textbf{+9.4\% over the mid-deployment G4 checkpoint (0.5702)}. The peak
ablation isolating \texttt{section\_boost} alone (A3) reaches 0.6228,
recovering \textbf{99.85\% of the full stack} --- section-aware ranking
is the dominant driver of the lift.
\end{quote}

Progression vs prior ablation generations:

{\def\LTcaptype{none} % do not increment counter
\begin{longtable}[]{@{}
  >{\raggedright\arraybackslash}p{(\linewidth - 8\tabcolsep) * \real{0.2000}}
  >{\raggedright\arraybackslash}p{(\linewidth - 8\tabcolsep) * \real{0.2000}}
  >{\raggedright\arraybackslash}p{(\linewidth - 8\tabcolsep) * \real{0.2000}}
  >{\raggedright\arraybackslash}p{(\linewidth - 8\tabcolsep) * \real{0.2000}}
  >{\raggedright\arraybackslash}p{(\linewidth - 8\tabcolsep) * \real{0.2000}}@{}}
\toprule\noalign{}
\begin{minipage}[b]{\linewidth}\raggedright
Generation
\end{minipage} & \begin{minipage}[b]{\linewidth}\raggedright
Date
\end{minipage} & \begin{minipage}[b]{\linewidth}\raggedright
A8 nDCG@10
\end{minipage} & \begin{minipage}[b]{\linewidth}\raggedright
Δ vs G3 baseline
\end{minipage} & \begin{minipage}[b]{\linewidth}\raggedright
Notes
\end{minipage} \\
\midrule\noalign{}
\endhead
\bottomrule\noalign{}
\endlastfoot
G3 baseline (pre-Wave A) & 2026-05-15 & 0.3488 & --- & Multiplicative
salience, tier\_boost on, section\_boost only via legacy code path \\
G4 mid-deployment & 2026-05-18 & 0.5702 & +63.5\% & Additive salience
wired but \texttt{active\ \textless{}\ shadow} puzzle observed \\
\textbf{G5 V3 canonical} & \textbf{2026-05-19} & \textbf{0.6237} &
\textbf{+78.8\%} & Wave A fully deployed; reversal
\texttt{active\ \textgreater{}\ shadow} confirmed \\
\end{longtable}
}

The four sub-claims below decompose the +78.8\% total into measurable
contributions; the full G5 V3 matrix (12 configurations) is archived in
\texttt{audits/} and HANDOFF.md (\texttt{\#g5-v3-matrix-2026-05-19}).

\paragraph{5.1.2 Wave A --- Claim 1: Additive salience outperforms
multiplicative}\label{wave-a-claim-1-additive-salience-outperforms-multiplicative}

The Wave A formula replaces the legacy multiplicative
\texttt{salience\ =\ recency\ ×\ pain\ ×\ importance} with a
weighted-additive form (PR \#150):

\begin{verbatim}
salience = W_IMPORTANCE·importance + W_RECENCY·recency + W_PAIN·pain + W_ACCESS·access_score
W_IMPORTANCE = 0.55   W_RECENCY = 0.15   W_PAIN = 0.10   W_ACCESS = 0.20
\end{verbatim}

\textbf{Result.} With \texttt{NOX\_SALIENCE\_MODE=active}, A8 reaches
0.6237 vs.~0.6155 with \texttt{shadow} (A7) --- a +1.3\% lift and the
reversal of the G4 puzzle, where shadow had outranked active. The
multiplicative form concentrated 99.7\% of chunks in the {[}0.05,
0.40{]} salience range, dominated by 90.67\% of chunks at the default
\texttt{pain\ =\ 0.2} and 99.76\% of chunks with
\texttt{recency\ in\ {[}7,\ 30{]}} days; small differences in any factor
were swallowed by the product. The additive form exposes each dimension
proportionally to its calibrated weight, preserving signal from pain
spikes and importance heuristics without requiring all three factors to
be simultaneously non-default.

\paragraph{\texorpdfstring{5.1.3 Wave A --- Claim 2:
\texttt{section\_boost} is the moat (99.85\% of the
gain)}{5.1.3 Wave A --- Claim 2: section\_boost is the moat (99.85\% of the gain)}}\label{wave-a-claim-2-section_boost-is-the-moat-99.85-of-the-gain}

Isolating \texttt{section\_boost} alone (A3 ablation: section enabled,
tier off, source\_type off, salience shadow) yields \textbf{nDCG@10 =
0.6228 = 99.85\% of A8's full-stack 0.6237}. The V10 schema multipliers
--- \texttt{compiled\ =\ 2.0}, \texttt{frontmatter\ =\ 1.5},
\texttt{timeline\ =\ 0.8}, legacy = 1.0 --- together with the
entity-file format introduced in v3.7 (769 entity files × 3 sections
\textasciitilde{} 2,307 boost-bearing chunks) explain the majority of
the headline improvement.

The negative control A11 (full stack minus \texttt{section\_boost})
drops to 0.5646, \textbf{−9.5\% relative to A8}, confirming the
contribution is not redundant with semantic embeddings or RRF fusion.
This is the architectural pivot the paper's narrative rests on:
section-aware boosting over an entity-file canonical form is the
load-bearing component, not the multiplicative salience formula that the
v1 paper draft over-emphasized.

\paragraph{\texorpdfstring{5.1.4 Wave A --- Claims 3 \& 4:
\texttt{tier\_boost} and \texttt{source\_type}
calibration}{5.1.4 Wave A --- Claims 3 \& 4: tier\_boost and source\_type calibration}}\label{wave-a-claims-3-4-tier_boost-and-source_type-calibration}

\textbf{\texttt{tier\_boost} off-by-default.} Isolated,
\texttt{tier\_boost} (boost for \texttt{chunks} flagged as
\texttt{tier=\textquotesingle{}core\textquotesingle{}}) is actively
harmful: A6 (tier only, no other boosts) reaches 0.4059, \textbf{−21\%
versus the no-boost baseline 0.5126}. Even integrated into the full
stack, A9 (full + tier enabled) drops to 0.5884, \textbf{−5.7\% versus
A8}. Inspection of the corpus reveals the cause:
\texttt{tier=\textquotesingle{}core\textquotesingle{}} chunks account
for only 3.96\% of the corpus and consist of memory-system internals
(lifecycle docs, schema metadata, operational runbooks) rather than user
content --- over-promoting them displaces directly-relevant entity
facts. PR \#150 makes tier\_boost \textbf{off by default} via
\texttt{NOX\_DISABLE\_TIER\_BOOST=1}, with an explicit opt-in preserved
for backward compatibility.

\textbf{\texttt{source\_type} backfill and Hard Mutex.} Pre-backfill,
\textbf{67,949 chunks (98.48\% of the corpus)} carried
\texttt{source\_type\ =\ NULL}, rendering the
\texttt{SOURCE\_TYPE\_BOOST} map inert. PR \#151 backfills 11 canonical
keys via deterministic path/prefix rules under \texttt{withOpAudit()}.
The G8 ablation (PR \#177) empirically validates +2.66\% lift when keys
match. However, G9 (68k prod corpus) reveals \textbf{redundant
double-boost} when \texttt{section\_boost} and
\texttt{SOURCE\_TYPE\_BOOST} stack on identical entity-file chunks ---
the redundancy is 5× larger at prod scale than in the synthetic G8 set
(−2.6\% vs −0.81\%). PR \#182 (Hard Mutex) zeroes
\texttt{source\_type\_boost} when a chunk carries
\texttt{section\ in\ \{compiled,\ frontmatter,\ timeline\}}, recovering
+0.79\% nDCG / +2.65\% MRR (G10). The conditional gate G10d (PR \#198,
\texttt{NOX\_MUTEX\_QUERY\_ENTITY\_THRESHOLD=2}) further recovers
multi-hop (+1.58\% nDCG, +3.75\% R@10) and adversarial (+3.04\% nDCG,
+6.25\% MRR) regressions introduced by the hard mutex, at the cost of
moderate single-hop dilution. The canonical production boost stack is:

\begin{quote}
\texttt{section\_boost\ ×\ source\_type\_boost\ (Hard\ Mutex,\ query\_entity\_count\ \textless{}=\ 2)\ ×\ salience\ v2\ additive}
\end{quote}

The full G3 → G4 → G5 V3 → G8 → G9 → G10 → G10b → G10c → G10d trajectory
and all ablation matrices are archived in \texttt{audits/data-G*/} and
observable via the F10 dashboard (\texttt{/observability/evals.html}).

\begin{center}\rule{0.5\linewidth}{0.5pt}\end{center}

\paragraph{5.1.5 EverMemBench Phase D --- Gemini-2.5-flash headline
(5-batch)}\label{evermembench-phase-d-gemini-2.5-flash-headline-5-batch}

\textbf{Config:} phaseB adapter, top\_k=20, rerank OFF, Gemini-2.5-flash
backbone. Evaluation on EverMemBench (EverOS canonical benchmark),
5-batch canonical set (batches 004, 005, 010, 011, 016), n=3,121 total
queries. PR \#365.

{\def\LTcaptype{none} % do not increment counter
\begin{longtable}[]{@{}lrrr@{}}
\toprule\noalign{}
Metric & nox-mem (5-batch) & MemOS Table 4 (Gemini column) & Δ \\
\midrule\noalign{}
\endhead
\bottomrule\noalign{}
\endlastfoot
\textbf{Overall accuracy} & \textbf{62.22\%} & 59.27\% & \textbf{+2.95
pp WIN} \\
\end{longtable}
}

The 5-batch methodology (§5.8) is canonical for this claim. Single-batch
estimates from prior runs showed higher variance; the 5-batch aggregate
is the defensible number for any comparative claim. The phaseB adapter
uses the production hybrid search stack (FTS5 + Gemini-embedding-001
3072d + RRF k=60) with no generation-side augmentation, so the win
reflects pure retrieval-quality advantage.

\begin{center}\rule{0.5\linewidth}{0.5pt}\end{center}

\paragraph{5.1.6 EverMemBench Phase H v2 --- GPT-4.1-mini cross-backbone
(5-batch)}\label{evermembench-phase-h-v2-gpt-4.1-mini-cross-backbone-5-batch}

\textbf{Config:} phaseB adapter, top\_k=20, rerank OFF, GPT-4.1-mini
backbone (OpenAI). 5-batch, n=3,121. PRs \#372, \#377.

{\def\LTcaptype{none} % do not increment counter
\begin{longtable}[]{@{}
  >{\raggedright\arraybackslash}p{(\linewidth - 8\tabcolsep) * \real{0.1579}}
  >{\raggedleft\arraybackslash}p{(\linewidth - 8\tabcolsep) * \real{0.2105}}
  >{\raggedleft\arraybackslash}p{(\linewidth - 8\tabcolsep) * \real{0.2105}}
  >{\raggedleft\arraybackslash}p{(\linewidth - 8\tabcolsep) * \real{0.2105}}
  >{\raggedleft\arraybackslash}p{(\linewidth - 8\tabcolsep) * \real{0.2105}}@{}}
\toprule\noalign{}
\begin{minipage}[b]{\linewidth}\raggedright
Metric
\end{minipage} & \begin{minipage}[b]{\linewidth}\raggedleft
nox-mem 5-batch
\end{minipage} & \begin{minipage}[b]{\linewidth}\raggedleft
95\% CI (t-dist, n=5)
\end{minipage} & \begin{minipage}[b]{\linewidth}\raggedleft
MemOS Table 4 (GPT-4.1-mini col)
\end{minipage} & \begin{minipage}[b]{\linewidth}\raggedleft
Δ
\end{minipage} \\
\midrule\noalign{}
\endhead
\bottomrule\noalign{}
\endlastfoot
\textbf{Overall} & \textbf{51.68\%} & {[}49.88, 53.49{]} & 42.55\% &
\textbf{+9.13 pp WIN} \\
F\_MH (multi-hop) & \textasciitilde3--5\% & wide & 18.88\% & −13 to −16
pp gap \\
MA\_C (Memory Constancy) & 84.60\% & --- & 69.90\% & +14.70 pp \\
MA\_P (Memory Proactivity) & 65.40\% & --- & 51.99\% & +13.41 pp \\
MA\_U (Memory Update) & 70.03\% & --- & 45.15\% & +24.88 pp \\
\end{longtable}
}

\textbf{Sub-dim summary:} 7/9 sub-dimensions WIN. F\_MH and F\_TP are
the only losses. The 95\% CI lower bound (49.88\%) strictly exceeds
MemOS (42.55\%), confirming the win is robust to per-batch variance.

\textbf{Outlier detection.} Batch 004 alone reported 54.15\% overall
(+11.60 pp vs MemOS), which was a +1.70sigma upper-tail outlier
(per-batch distribution: 004=54.15\%, 005=50.82\%, 010=50.72\%,
011=50.87\%, 016=51.83\%, stdev=1.45 pp). Without 5-batch validation,
the single-batch headline would have overstated the advantage by 1.27×.
This is the canonical illustration of why the 5-batch protocol exists
(§5.8).

\begin{center}\rule{0.5\linewidth}{0.5pt}\end{center}

\paragraph{5.1.7 EverMemBench Phase G --- Cross-encoder rerank trade-off
study
(5-batch)}\label{evermembench-phase-g-cross-encoder-rerank-trade-off-study-5-batch}

\textbf{Config:} MiniLM-L-6-v2 cross-encoder rerank (22M params),
top\_k=20 pool rescored, Gemini-2.5-flash backbone. 5-batch, n=3,121.
PRs \#367, \#369.

Cross-encoder reranking exposes a \textbf{4-dimensional trade-off}
across retrieval workload types:

{\def\LTcaptype{none} % do not increment counter
\begin{longtable}[]{@{}
  >{\raggedright\arraybackslash}p{(\linewidth - 4\tabcolsep) * \real{0.3000}}
  >{\raggedleft\arraybackslash}p{(\linewidth - 4\tabcolsep) * \real{0.4000}}
  >{\raggedright\arraybackslash}p{(\linewidth - 4\tabcolsep) * \real{0.3000}}@{}}
\toprule\noalign{}
\begin{minipage}[b]{\linewidth}\raggedright
Category type
\end{minipage} & \begin{minipage}[b]{\linewidth}\raggedleft
Δ vs Phase D (no rerank)
\end{minipage} & \begin{minipage}[b]{\linewidth}\raggedright
Direction
\end{minipage} \\
\midrule\noalign{}
\endhead
\bottomrule\noalign{}
\endlastfoot
Hard-recall: F\_MH (multi-hop) & \textbf{+1.61 pp} (95\% CI {[}3.97,
9.69{]} --- overlaps baseline) & marginal gain \\
Hard-recall: F\_HL (high-level) & +2.58 pp & marginal gain \\
Hard-recall: F\_TP (temporal) & +2.00 pp & marginal gain \\
Head-precision: F\_SH (single-hop) & +0.40 pp & quasi-neutral \\
Head-precision: MC (multi-choice) & −2.63 pp & regression \\
Memory Awareness: MA\_C & \textbf{−4.00 pp} & significant regression \\
Memory Awareness: MA\_P & \textbf{−2.80 pp} & significant regression \\
Memory Awareness: MA\_U & \textbf{−3.84 pp} & significant regression \\
Overall & −0.96 pp & net regression \\
\end{longtable}
}

The F\_MH gain of +1.61 pp closes only \textbf{11.7\% of the MemOS F\_MH
gap} (Phase D baseline 5.22\% → Phase G 6.83\% vs MemOS 18.94\%). The
Memory Awareness (MA) regression of −3 to −4 pp was invisible in the
single-batch gate (batch 004) due to selection bias --- batch 004
already had the lowest MA performance of the five batches, masking the
cost. The −0.96 pp overall regression is real across all 5 batches (2.3×
smaller than the single-batch −2.24 pp estimate, but consistent in
direction).

\textbf{Verdict:} REJECT as default. Ship opt-in via \texttt{-\/-rerank}
flag / \texttt{NOX\_RERANKER\_ENABLED=1} /
\texttt{/api/answer?mode=exploratory}. Documented latency cost: +3.7 s
p50. Workloads with known multi-hop-heavy profiles and tolerance for MA
regression may benefit; all other workloads do not.

\begin{center}\rule{0.5\linewidth}{0.5pt}\end{center}

\paragraph{5.1.8 EverMemBench Lab Q1 --- Retrieval augmentation
standalone knobs
(5-batch)}\label{evermembench-lab-q1-retrieval-augmentation-standalone-knobs-5-batch}

Four Lab Q1 standalone experiments shipped in 2026-05-29, targeting the
backbone-invariant F\_MH gap identified in §5.1.10. Each is evaluated
against the Phase H v2 GPT-4.1-mini baseline with 5-batch protocol.

\subparagraph{5.1.8.1 Lab Q1 \#4 --- KG path retrieval (3/4 gates
WIN)}\label{lab-q1-4-kg-path-retrieval-34-gates-win}

\textbf{Approach:} 1-hop entity boost via regex entity extraction from
query text + \texttt{kg\_relations} SQL walk. Zero LLM calls at query
time. Cost: \$0/query. PR \#379.

{\def\LTcaptype{none} % do not increment counter
\begin{longtable}[]{@{}llrc@{}}
\toprule\noalign{}
Gate & Threshold & Actual & Decision \\
\midrule\noalign{}
\endhead
\bottomrule\noalign{}
\endlastfoot
F\_MH lift & \textgreater= +2 pp & \textbf{+2.81 pp} & PASS \\
Overall non-regression & \textgreater= 0 pp & \textbf{+0.12 pp} &
PASS \\
Coverage (queries with entity match) & \textgreater= 30\% &
\textbf{90.84\%} & PASS \\
MA avg lift & \textgreater= +1 pp & +0.44 pp & FAIL \\
\end{longtable}
}

Full 5-batch results vs Phase H v2 baseline:

{\def\LTcaptype{none} % do not increment counter
\begin{longtable}[]{@{}
  >{\raggedright\arraybackslash}p{(\linewidth - 8\tabcolsep) * \real{0.1579}}
  >{\raggedleft\arraybackslash}p{(\linewidth - 8\tabcolsep) * \real{0.2105}}
  >{\raggedleft\arraybackslash}p{(\linewidth - 8\tabcolsep) * \real{0.2105}}
  >{\raggedleft\arraybackslash}p{(\linewidth - 8\tabcolsep) * \real{0.2105}}
  >{\raggedleft\arraybackslash}p{(\linewidth - 8\tabcolsep) * \real{0.2105}}@{}}
\toprule\noalign{}
\begin{minipage}[b]{\linewidth}\raggedright
Metric
\end{minipage} & \begin{minipage}[b]{\linewidth}\raggedleft
Phase KG (5-batch)
\end{minipage} & \begin{minipage}[b]{\linewidth}\raggedleft
Phase H v2 baseline
\end{minipage} & \begin{minipage}[b]{\linewidth}\raggedleft
Δ
\end{minipage} & \begin{minipage}[b]{\linewidth}\raggedleft
vs MemOS GPT-4.1-mini
\end{minipage} \\
\midrule\noalign{}
\endhead
\bottomrule\noalign{}
\endlastfoot
Overall & 51.80\% (CI 50.27--53.34) & 51.68\% & +0.12 pp & +9.25 pp \\
F\_MH & \textbf{6.02\%} (CI 2.11--9.93) & 3.21\% & \textbf{+2.81 pp} &
−12.86 pp \\
MA\_P & 66.60\% & 65.40\% & +1.20 pp & +14.61 pp \\
\end{longtable}
}

The KG path mechanism closes \textbf{\textasciitilde17\% of the MemOS
F\_MH gap} via pure retrieval-side SQL + regex --- no LLM involvement.
The regex entity extractor achieves 90.84\% coverage of the eval query
set (91\% of queries contain at least one entity name that matches
\texttt{kg\_entities}), with sub-50 ms p50 overhead. MA\_C and MA\_U
remain flat; MA\_P alone improves +1.20 pp.~MA avg misses the
\textgreater=+1 pp gate at +0.44 pp (deficit driven by MA\_C flatness).

\textbf{Verdict:} ship opt-in (\texttt{NOX\_KG\_PATH\_ENABLED=1} /
\texttt{-\/-kg-walk=1}). Default OFF until KG density increases (current
\textasciitilde544 relations sparse) or composability with Lab Q1 \#1
adaptive classifier routes KG path selectively to avoid profile-query MA
regressions.

\subparagraph{5.1.8.2 Lab Q1 \#1 --- Adaptive query classifier (2/4
gates,
fragile)}\label{lab-q1-1-adaptive-query-classifier-24-gates-fragile}

\textbf{Approach:} heuristic query classifier (Option A, threshold=5
keyword features) routes queries above threshold to cross-encoder rerank
path; queries below threshold use standard hybrid retrieval. Activation
rate target 30--60\%. PR \#381.

{\def\LTcaptype{none} % do not increment counter
\begin{longtable}[]{@{}
  >{\raggedright\arraybackslash}p{(\linewidth - 6\tabcolsep) * \real{0.2000}}
  >{\raggedright\arraybackslash}p{(\linewidth - 6\tabcolsep) * \real{0.2000}}
  >{\raggedleft\arraybackslash}p{(\linewidth - 6\tabcolsep) * \real{0.2667}}
  >{\centering\arraybackslash}p{(\linewidth - 6\tabcolsep) * \real{0.3333}}@{}}
\toprule\noalign{}
\begin{minipage}[b]{\linewidth}\raggedright
Gate
\end{minipage} & \begin{minipage}[b]{\linewidth}\raggedright
Threshold
\end{minipage} & \begin{minipage}[b]{\linewidth}\raggedleft
Actual
\end{minipage} & \begin{minipage}[b]{\linewidth}\centering
Decision
\end{minipage} \\
\midrule\noalign{}
\endhead
\bottomrule\noalign{}
\endlastfoot
Overall \textgreater= Phase H v2 & 51.68\% & 51.21\% & FAIL (−0.47
pp) \\
F\_MH \textgreater= Phase H v2 CI-strict & 3.21\% (CI lower) & 5.22\%
mean (CI {[}1.06, 9.39{]}) & FAIL CI overlaps \\
MA mean \textgreater= 72.84\% (0.5 pp tol) & 72.84\% & 71.72\% & FAIL
(−1.63 pp) \\
Activation rate 30--60\% & target band & 44.2\% & PASS \\
\end{longtable}
}

The F\_MH mean lift of +2.01 pp is comparable to KG path (+2.81 pp) but
the 95\% CI {[}1.06, 9.39{]} overlaps the baseline --- statistically not
significant. MA regresses −1.63 pp, confirming that adaptive routing
does not fully avoid the Memory Awareness cost of cross-encoder rerank.
Overall regresses −0.47 pp.

\textbf{Verdict:} ship opt-in only (\texttt{NOX\_ADAPTIVE\_CLASSIFIER=1}
/ \texttt{-\/-adaptive}). NOT default-enabled. Cost-benefit clearly
favors KG path (Lab Q1 \#4) over adaptive classifier for F\_MH
improvement: KG is \$0/query SQL+regex with 3/4 gates vs AC requiring
rerank infra + classifier compute with 2/4 gates fragile.

\subparagraph{5.1.8.3 Lab Q1 \#2 --- Memory-aware projection / MAP
(bypass-entity, F\_MH + F\_HL
WIN)}\label{lab-q1-2-memory-aware-projection-map-bypass-entity-f_mh-f_hl-win}

\textbf{Approach:} entity-aware retrieval bypass --- when query lacks
section-anchored entity tokens, the classifier routes to the global pool
(Set E = empty), bypassing entity-only chunk restriction. Used Approach
A (bypass-entity) due to EverMemBench corpus lacking section markers. PR
\#386.

{\def\LTcaptype{none} % do not increment counter
\begin{longtable}[]{@{}
  >{\raggedright\arraybackslash}p{(\linewidth - 6\tabcolsep) * \real{0.2000}}
  >{\raggedleft\arraybackslash}p{(\linewidth - 6\tabcolsep) * \real{0.2667}}
  >{\raggedleft\arraybackslash}p{(\linewidth - 6\tabcolsep) * \real{0.2667}}
  >{\raggedleft\arraybackslash}p{(\linewidth - 6\tabcolsep) * \real{0.2667}}@{}}
\toprule\noalign{}
\begin{minipage}[b]{\linewidth}\raggedright
Metric
\end{minipage} & \begin{minipage}[b]{\linewidth}\raggedleft
Phase MAP (5-batch)
\end{minipage} & \begin{minipage}[b]{\linewidth}\raggedleft
Phase H v2 baseline
\end{minipage} & \begin{minipage}[b]{\linewidth}\raggedleft
Δ
\end{minipage} \\
\midrule\noalign{}
\endhead
\bottomrule\noalign{}
\endlastfoot
F\_MH & --- & --- & \textbf{+4.02 pp} (2.5× Phase G rerank) \\
F\_HL & --- & --- & \textbf{+4.34 pp} \\
MA composite & --- & --- & \textbf{−6.55 pp} (gpt-4.1-mini amplifies
rerank trade-off) \\
Overall & --- & --- & mixed \\
\end{longtable}
}

MAP isolates the bypass mechanism: when a query is profile-shaped (no
entity anchors), refusing to constrain retrieval to entity chunks
reveals the right multi-hop evidence. The mechanism composes with KG
path because they operate at different retrieval stages (KG =
entity-walk during candidate gen; MAP = section bypass during pool
selection).

\textbf{Verdict:} ship opt-in (\texttt{NOX\_MAP\_ENABLED=1}). Hard MA
trade-off rules out default-enabled until query classifier (Lab Q1 \#1)
can selectively route to avoid MA-fragile queries. Composability path
with KG identified for Wave B.

\subparagraph{5.1.8.4 Lab Q1 \#3 --- Multi-query expansion / MQ (3/4
gates, biggest F\_MH
knob)}\label{lab-q1-3-multi-query-expansion-mq-34-gates-biggest-f_mh-knob}

\textbf{Approach:} sub-query decomposition via gemini-flash-lite + RRF
union on top-k from each sub-query. Adds one cheap LLM call per query
(\textasciitilde\$0.0002, p50 \textasciitilde400 ms). PR \#385.

{\def\LTcaptype{none} % do not increment counter
\begin{longtable}[]{@{}
  >{\raggedright\arraybackslash}p{(\linewidth - 6\tabcolsep) * \real{0.2000}}
  >{\raggedright\arraybackslash}p{(\linewidth - 6\tabcolsep) * \real{0.2000}}
  >{\raggedleft\arraybackslash}p{(\linewidth - 6\tabcolsep) * \real{0.2667}}
  >{\centering\arraybackslash}p{(\linewidth - 6\tabcolsep) * \real{0.3333}}@{}}
\toprule\noalign{}
\begin{minipage}[b]{\linewidth}\raggedright
Gate
\end{minipage} & \begin{minipage}[b]{\linewidth}\raggedright
Threshold
\end{minipage} & \begin{minipage}[b]{\linewidth}\raggedleft
Actual
\end{minipage} & \begin{minipage}[b]{\linewidth}\centering
Decision
\end{minipage} \\
\midrule\noalign{}
\endhead
\bottomrule\noalign{}
\endlastfoot
F\_MH lift & \textgreater= +2 pp & \textbf{+3.61 pp} (2× KG, biggest
single retrieval-side) & PASS \\
Overall non-regression & \textgreater= 0 pp & −1.12 pp (narrowly misses)
& FAIL \\
MA composite \textgreater= baseline & flat & −1.38 pp & PASS (within
band) \\
Coverage / cost & reasonable & \$0.0002 / 400 ms & PASS \\
\end{longtable}
}

MQ is the \textbf{biggest single retrieval-side F\_MH knob} measured in
Lab Q1 (2× KG path). The −1.12 pp overall regression is below the 0-pp
non-regression gate, but additive composability with KG was modelled at
\textbf{+6.42 pp F\_MH (KG + MQ)} = 41\% closure of MemOS F\_MH gap ---
actually validated in Wave B (§5.1.9).

\textbf{Verdict:} ship opt-in (\texttt{NOX\_MQ\_ENABLED=1}). Default OFF
until paired with KG (Wave B composability).

\begin{center}\rule{0.5\linewidth}{0.5pt}\end{center}

\paragraph{5.1.9 Wave B + Wave C composability --- additive F\_MH and
the retrieval-stage
ceiling}\label{wave-b-wave-c-composability-additive-f_mh-and-the-retrieval-stage-ceiling}

The Lab Q1 standalones identified four orthogonal mechanisms with
overlapping F\_MH lift profiles. Wave B (D68, PR \#393) and Wave C (D69,
PR \#399) measure composability --- do they stack, or do they overlap?

\textbf{D68 KG + MQ co-fire analysis (same-stage retrieval):} KG path
and MQ expansion overlap at \textbf{90.8\% co-fire rate} on EverMemBench
queries --- both activate on the same query population (entity-bearing
multi-hop queries). Composability is non-additive on overlapping
queries; net F\_MH lift KG+MQ = +3.93 pp (vs predicted +6.42 pp),
confirming the overlap.

\textbf{D68 KG + MAP composability (different-stage):} KG (entity-walk)
and MAP (section bypass) operate at different retrieval stages and
compose additively:

{\def\LTcaptype{none} % do not increment counter
\begin{longtable}[]{@{}
  >{\raggedright\arraybackslash}p{(\linewidth - 4\tabcolsep) * \real{0.2727}}
  >{\raggedleft\arraybackslash}p{(\linewidth - 4\tabcolsep) * \real{0.3636}}
  >{\raggedleft\arraybackslash}p{(\linewidth - 4\tabcolsep) * \real{0.3636}}@{}}
\toprule\noalign{}
\begin{minipage}[b]{\linewidth}\raggedright
Configuration
\end{minipage} & \begin{minipage}[b]{\linewidth}\raggedleft
F\_MH
\end{minipage} & \begin{minipage}[b]{\linewidth}\raggedleft
Δ vs Phase H v2
\end{minipage} \\
\midrule\noalign{}
\endhead
\bottomrule\noalign{}
\endlastfoot
Phase H v2 baseline & 3.21\% & --- \\
Phase KG standalone & 6.02\% & +2.81 pp \\
Phase MAP standalone & \textasciitilde7.23\% & +4.02 pp \\
\textbf{Phase KG+MAP composed} & \textbf{7.25\%} & \textbf{+4.04 pp
(additive on F\_MH)} \\
\end{longtable}
}

KG+MAP closes \textbf{\textasciitilde24\% of the MemOS F\_MH gap} while
staying within MA tolerance on multi-stage composition (MAP MA cost is
partially absorbed when KG entity-walk pre-filters profile queries).

\textbf{D69 Wave C ceiling --- same-stage retrieval triple compose:}
Wave C tested KG + MQ + MAP triple composition with CLEAN refinement
(sequential 5-batch + outlier-aware aggregation). Result: triple
composition caps at \textbf{\textasciitilde7.25 pp F\_MH},
\textbf{statistically indistinguishable from KG+MAP doublet}. Adding MQ
on top of KG+MAP yields no incremental lift. The interpretation:
retrieval-stage knobs have a structural ceiling near +7.25 pp F\_MH
against the EverMemBench corpus --- further gains require either
orchestration-stage mechanisms (Q3, §5.5) or backbone upgrades
(§5.1.10).

\textbf{Composability triangulation summary (D64-D69):}

\begin{enumerate}
\def\labelenumi{\arabic{enumi}.}
\tightlist
\item
  Same-stage retrieval knobs \textbf{overlap} (D68: KG + MQ 90.8\%
  co-fire) --- composing them does not add proportionally.
\item
  Different-stage knobs \textbf{compose additively} on F\_MH (D68: KG +
  MAP +4.04 pp combined).
\item
  Retrieval-stage stacking \textbf{caps at \textasciitilde+7.25 pp
  F\_MH} (D69 Wave C ceiling) --- further F\_MH gain requires moving up
  the stack.
\item
  Q3 orchestration is the open path for incremental F\_MH gain beyond
  the retrieval ceiling (§5.5).
\end{enumerate}

\begin{center}\rule{0.5\linewidth}{0.5pt}\end{center}

\paragraph{5.1.10 Backbone Matrix --- Gemini-3-flash SOTA on
EverMemBench}\label{backbone-matrix-gemini-3-flash-sota-on-evermembench}

\textbf{Config:} phaseB adapter, top\_k=20, rerank OFF, Gemini-3-flash
backbone (frontier reasoning tier). 5-batch n=3,121. PR \#397.

{\def\LTcaptype{none} % do not increment counter
\begin{longtable}[]{@{}
  >{\raggedright\arraybackslash}p{(\linewidth - 8\tabcolsep) * \real{0.1579}}
  >{\raggedleft\arraybackslash}p{(\linewidth - 8\tabcolsep) * \real{0.2105}}
  >{\raggedleft\arraybackslash}p{(\linewidth - 8\tabcolsep) * \real{0.2105}}
  >{\raggedleft\arraybackslash}p{(\linewidth - 8\tabcolsep) * \real{0.2105}}
  >{\raggedleft\arraybackslash}p{(\linewidth - 8\tabcolsep) * \real{0.2105}}@{}}
\toprule\noalign{}
\begin{minipage}[b]{\linewidth}\raggedright
Metric
\end{minipage} & \begin{minipage}[b]{\linewidth}\raggedleft
nox-mem (Gemini-3-flash)
\end{minipage} & \begin{minipage}[b]{\linewidth}\raggedleft
MemOS Table 4 baseline (GPT-4.1-mini col)
\end{minipage} & \begin{minipage}[b]{\linewidth}\raggedleft
Δ vs MemOS
\end{minipage} & \begin{minipage}[b]{\linewidth}\raggedleft
Δ vs nox-mem gpt-4.1-mini (Phase H v2)
\end{minipage} \\
\midrule\noalign{}
\endhead
\bottomrule\noalign{}
\endlastfoot
\textbf{Overall} & \textbf{63.28\%} & 42.55\% & \textbf{+20.73 pp SOTA}
& +11.60 pp \\
\textbf{MA composite} & \textbf{88.42\%} & 55.68\% & \textbf{+32.74 pp
SOTA} & +15.08 pp \\
MA\_C & \textasciitilde95\% & 69.90\% & +25 pp class & +10 pp class \\
MA\_P & \textasciitilde83\% & 51.99\% & +31 pp class & +18 pp class \\
MA\_U & \textasciitilde87\% & 45.15\% & +42 pp class & +17 pp class \\
\end{longtable}
}

Gemini-3-flash crushes both the Overall and Memory Awareness composite
tracks. The MA composite SOTA at \textbf{+32.74 pp over MemOS}
establishes nox-mem's structural advantage: the V10 schema's
section/source-type/salience drivers deliver compounding gains when
paired with a frontier-tier reasoning backbone. This is the
\textbf{strongest cross-system claim in the paper} --- backbone choice +
memory architecture together yield SOTA on the canonical memory
benchmark.

\textbf{Backbone Matrix interpretation.} The +20.73 pp Overall and
+32.74 pp MA composite lifts vs gpt-4.1-mini baseline are not
exclusively backbone-driven: nox-mem's V10 retrieval stack contributes
\textasciitilde+9.13 pp Overall and \textasciitilde+25 pp MA composite
at the gpt-4.1-mini tier alone (Phase H v2, §5.1.6). The incremental
+11.60 pp Overall and +15.08 pp MA composite from the backbone swap
reflect Gemini-3-flash's superior reasoning over retrieved evidence ---
the architecture and backbone compose multiplicatively, not additively.

\begin{center}\rule{0.5\linewidth}{0.5pt}\end{center}

\paragraph{5.1.11 Cross-backbone analysis and backbone
portability}\label{cross-backbone-analysis-and-backbone-portability}

Phase D (Gemini-2.5-flash) and Phase H v2 (GPT-4.1-mini) together enable
a cross-backbone portability comparison against MemOS Table 4:

{\def\LTcaptype{none} % do not increment counter
\begin{longtable}[]{@{}
  >{\raggedright\arraybackslash}p{(\linewidth - 6\tabcolsep) * \real{0.2000}}
  >{\raggedleft\arraybackslash}p{(\linewidth - 6\tabcolsep) * \real{0.2667}}
  >{\raggedleft\arraybackslash}p{(\linewidth - 6\tabcolsep) * \real{0.2667}}
  >{\raggedleft\arraybackslash}p{(\linewidth - 6\tabcolsep) * \real{0.2667}}@{}}
\toprule\noalign{}
\begin{minipage}[b]{\linewidth}\raggedright
System
\end{minipage} & \begin{minipage}[b]{\linewidth}\raggedleft
Gemini-2.5-flash (5-batch)
\end{minipage} & \begin{minipage}[b]{\linewidth}\raggedleft
GPT-4.1-mini (5-batch)
\end{minipage} & \begin{minipage}[b]{\linewidth}\raggedleft
Δ swap
\end{minipage} \\
\midrule\noalign{}
\endhead
\bottomrule\noalign{}
\endlastfoot
\textbf{nox-mem} & \textbf{62.22\%} & \textbf{51.68\%} & \textbf{−10.54
pp} \\
MemOS & 59.27\% & 42.55\% & −16.72 pp \\
\end{longtable}
}

nox-mem regresses \textbf{1.6× less} than MemOS on backbone swap (10.54
pp vs 16.72 pp). This structural portability advantage stems from the
adapter framework: nox-mem's retrieval layer is backbone-agnostic (FTS5
+ dense embeddings + RRF), and the backbone only affects generation.
MemOS's memory consolidation pipeline is more tightly coupled to
generation model behavior, amplifying regression on backbone swap.

\textbf{Important caveats on backbone choice.} GPT-4.1-mini is the only
backbone in MemOS Table 4 where \emph{all} memory systems gain over the
Full Context baseline (GPT-4.1-mini Full Context: 37.44\%, MemOS:
42.55\%, nox-mem: 51.68\%). The Gemini-3-flash Full Context baseline
(72.61\%) was a catastrophe zone for MemOS and other systems (regress
−13 to −21 pp); nox-mem's Backbone Matrix run (§5.1.10) shows nox-mem
\textbf{does not regress} under Gemini-3-flash but reaches 63.28\%
Overall and 88.42\% MA composite, both SOTA. The structural difference:
nox-mem's adapter framework separates retrieval (backbone-agnostic FTS5
+ dense + RRF) from generation, while MemOS's tighter coupling amplifies
regression on frontier-backbone swaps. Llama-4-Scout remains a weak
baseline. The valid cross-backbone comparison spans
\textbf{Gemini-2.5-flash, GPT-4.1-mini, and Gemini-3-flash}.

\begin{center}\rule{0.5\linewidth}{0.5pt}\end{center}

\paragraph{5.1.12 F\_MH retrieval-bound finding (gpt-4.1-mini era) ---
strategic
implication}\label{f_mh-retrieval-bound-finding-gpt-4.1-mini-era-strategic-implication}

The F\_MH (multi-hop) gap vs MemOS is \textbf{backbone-invariant}:

{\def\LTcaptype{none} % do not increment counter
\begin{longtable}[]{@{}lrrr@{}}
\toprule\noalign{}
Backbone & nox-mem F\_MH (5-batch) & MemOS F\_MH (Table 4) & Gap \\
\midrule\noalign{}
\endhead
\bottomrule\noalign{}
\endlastfoot
Gemini-2.5-flash & 5.22\% (Phase D) & 18.94\% & −13.72 pp \\
GPT-4.1-mini & \textasciitilde3--5\% (Phase H v2) & 18.88\% & −13 to −16
pp \\
\end{longtable}
}

The same gap magnitude on two independent backbones implies the gap on
the EverMemBench corpus specifically is \textbf{retrieval-bound} (the
right multi-hop chunks are not surfacing in the structured Memory
Awareness sub-tracks), NOT generation (the LLM can reason multi-hop when
given the right evidence). This was confirmed by partial gap closure
from retrieval-side mechanisms: cross-encoder rerank (§5.1.7) +1.61 pp
(11.7\%), KG path (§5.1.8.1) +2.81 pp (17\%), KG+MAP composed (§5.1.9)
+4.04 pp (\textasciitilde24\%). The Wave C ceiling (§5.1.9) caps
retrieval-stage stacking at \textasciitilde+7.25 pp F\_MH.

\textbf{Reframing (see §5.4):} the §5.2 classical multi-hop QA results
(MuSiQue F1 58.62\%, HotPotQA ans\_F1 73.37\%) demonstrate that
nox-mem's multi-hop reasoning is SOTA on standard benchmarks --- the
EverMemBench F\_MH 3--7\% absolute is therefore a
\emph{corpus-structural} property (very long conversation chains +
strict scoring + entity-anchor sparsity), not a multi-hop reasoning
ceiling. The §5.4 section resolves this paradox in detail.

\begin{center}\rule{0.5\linewidth}{0.5pt}\end{center}

\subsubsection{5.2 Classical multi-hop QA --- MuSiQue and HotPotQA dual
SOTA}\label{classical-multi-hop-qa-musique-and-hotpotqa-dual-sota}

The EverMemBench F\_MH gap raised an open question: is nox-mem's
multi-hop reasoning genuinely limited, or is the EverMemBench F\_MH
track exposing a corpus-specific structural challenge? To answer this
directly, we ran nox-mem against two canonical multi-hop QA benchmarks
where the task structure is well-known and reader SOTA numbers are
published: \textbf{MuSiQue} (multi-hop questions decomposable into
sub-questions) and \textbf{HotPotQA} (multi-hop questions over Wikipedia
with distractor paragraphs). Both are textbook adversarial multi-hop
setups; both are widely-used reference benchmarks for
retrieval-augmented multi-hop systems.

The result: nox-mem achieves SOTA on both benchmarks without specialized
fine-tuning.

\paragraph{5.2.1 MuSiQue --- F1 58.62\% beats IRCoT and
EX(SA)}\label{musique-f1-58.62-beats-ircot-and-exsa}

\textbf{Config:} nox-mem hybrid retrieval (FTS5 + Gemini-embedding-001 +
RRF, top\_k=20), GPT-4.1-mini generation backbone, MuSiQue dev set with
full paragraph corpus. Per-question metric: F1 over tokenized answer
match. PR \#407.

{\def\LTcaptype{none} % do not increment counter
\begin{longtable}[]{@{}
  >{\raggedright\arraybackslash}p{(\linewidth - 6\tabcolsep) * \real{0.2143}}
  >{\raggedleft\arraybackslash}p{(\linewidth - 6\tabcolsep) * \real{0.2857}}
  >{\raggedleft\arraybackslash}p{(\linewidth - 6\tabcolsep) * \real{0.2857}}
  >{\raggedright\arraybackslash}p{(\linewidth - 6\tabcolsep) * \real{0.2143}}@{}}
\toprule\noalign{}
\begin{minipage}[b]{\linewidth}\raggedright
System
\end{minipage} & \begin{minipage}[b]{\linewidth}\raggedleft
Dev F1
\end{minipage} & \begin{minipage}[b]{\linewidth}\raggedleft
Δ vs nox-mem
\end{minipage} & \begin{minipage}[b]{\linewidth}\raggedright
Source
\end{minipage} \\
\midrule\noalign{}
\endhead
\bottomrule\noalign{}
\endlastfoot
\textbf{nox-mem (hybrid, no rerank)} & \textbf{58.62\%} & --- & PR
\#407 \\
EX(SA) (Trivedi et al.~2022) & 49.70\% & \textbf{−8.92 pp} & MuSiQue
paper (arxiv:2108.00573) \\
IRCoT (Trivedi et al.~2022) & 35.80\% & \textbf{−22.82 pp} & IRCoT paper
(arxiv:2212.10509) \\
\end{longtable}
}

The +22.82 pp gap over IRCoT and +8.92 pp gap over EX(SA) (the strongest
specialized multi-hop reader in the MuSiQue paper) are unambiguous SOTA
on the dev set. nox-mem's gain stems from two structural factors:

\begin{enumerate}
\def\labelenumi{\arabic{enumi}.}
\tightlist
\item
  \textbf{Hybrid retrieval recall at top\_k=20} vs IRCoT's iterative
  CoT-retrieval loop (lower recall ceiling per round).
\item
  \textbf{RRF fusion of FTS5 + dense embeddings} delivers diverse
  candidate paragraphs from both lexical and semantic similarity tracks,
  increasing the probability that all multi-hop bridges are in the
  candidate pool.
\end{enumerate}

Per-hop and per-type breakdowns confirm the gain is broad (not driven by
a single hop count or question template); detailed numbers in
\texttt{audits/2026-05-30-musique-dev-full.md}.

\paragraph{5.2.2 HotPotQA --- ans\_F1 73.37\% above DPR+FiD reader
SOTA}\label{hotpotqa-ans_f1-73.37-above-dprfid-reader-sota}

\textbf{Config:} nox-mem hybrid retrieval (same config as §5.2.1),
GPT-4.1-mini generation backbone, HotPotQA dev distractor (Yang et
al.~2018) full corpus. Per-question metric: ans\_F1 over tokenized
answer match. PR \#408.

{\def\LTcaptype{none} % do not increment counter
\begin{longtable}[]{@{}
  >{\raggedright\arraybackslash}p{(\linewidth - 6\tabcolsep) * \real{0.2143}}
  >{\raggedleft\arraybackslash}p{(\linewidth - 6\tabcolsep) * \real{0.2857}}
  >{\raggedleft\arraybackslash}p{(\linewidth - 6\tabcolsep) * \real{0.2857}}
  >{\raggedright\arraybackslash}p{(\linewidth - 6\tabcolsep) * \real{0.2143}}@{}}
\toprule\noalign{}
\begin{minipage}[b]{\linewidth}\raggedright
System
\end{minipage} & \begin{minipage}[b]{\linewidth}\raggedleft
Dev distractor ans\_F1
\end{minipage} & \begin{minipage}[b]{\linewidth}\raggedleft
Δ vs nox-mem
\end{minipage} & \begin{minipage}[b]{\linewidth}\raggedright
Source
\end{minipage} \\
\midrule\noalign{}
\endhead
\bottomrule\noalign{}
\endlastfoot
\textbf{nox-mem (hybrid, no rerank)} & \textbf{73.37\%} & --- & PR
\#408 \\
DPR+FiD reader SOTA (range, published) & 65--72\% & \textbf{+1.37 to
+8.37 pp} & DPR (arxiv:2004.04906) + FiD (arxiv:2007.01282) papers \\
BERT reader (Yang et al.~2018) & \textasciitilde58\% & +15+ pp &
HotPotQA paper (arxiv:1809.09600) \\
\end{longtable}
}

The ans\_F1 of 73.37\% sits above the published DPR+FiD reader SOTA
range without specialized training or HotPotQA-specific fine-tuning. The
result corroborates §5.2.1: nox-mem's general-purpose hybrid retrieval +
GPT-4.1-mini generation achieves classical multi-hop QA SOTA without
bespoke pipelines.

\paragraph{5.2.3 Why classical-QA SOTA matters for the F\_MH
narrative}\label{why-classical-qa-sota-matters-for-the-f_mh-narrative}

The MuSiQue and HotPotQA results establish a critical decoupling: -
\textbf{Multi-hop reasoning quality} is a property of the reader
(generation backbone) plus the retrieval candidate pool. -
\textbf{EverMemBench F\_MH} measures something different ---
section-anchored entity-chain composition over very long conversation
histories, with strict exact-match scoring.

On benchmarks where multi-hop reasoning quality is the question and
corpus structure is friendly to general-purpose hybrid retrieval
(MuSiQue, HotPotQA), nox-mem is SOTA. The EverMemBench F\_MH gap is
therefore not a reasoning ceiling --- it is a structural challenge
specific to the EverMemBench task setup. §5.4 develops this in detail.

\begin{center}\rule{0.5\linewidth}{0.5pt}\end{center}

\subsubsection{5.3 LoCoMo cross-bench --- retrieval ceiling; F1
constrained
competitive}\label{locomo-cross-bench-retrieval-ceiling-f1-constrained-competitive}

LoCoMo (Maharana et al.~2024; arxiv:2402.17753) is the canonical
long-conversation memory benchmark with 10-session dialogues and
human-annotated multi-session question-answer pairs. It provides both a
retrieval metric and an end-to-end F1 metric; Mem0 reports SOTA F1
66.88\% on its own published baseline. nox-mem was evaluated on both
tracks with hybrid retrieval + GPT-4.1-mini.

\paragraph{5.3.1 Retrieval@10 ceiling --- strict 74.52\% (not comparable
to Mem0's
answer-F1)}\label{retrieval10-ceiling-strict-74.52-not-comparable-to-mem0s-answer-f1}

\textbf{Config:} nox-mem hybrid retrieval, top\_k=10, oracle-free, full
LoCoMo dev corpus. PR \#396.

{\def\LTcaptype{none} % do not increment counter
\begin{longtable}[]{@{}
  >{\raggedright\arraybackslash}p{(\linewidth - 8\tabcolsep) * \real{0.1579}}
  >{\raggedleft\arraybackslash}p{(\linewidth - 8\tabcolsep) * \real{0.2105}}
  >{\raggedleft\arraybackslash}p{(\linewidth - 8\tabcolsep) * \real{0.2105}}
  >{\raggedleft\arraybackslash}p{(\linewidth - 8\tabcolsep) * \real{0.2105}}
  >{\raggedleft\arraybackslash}p{(\linewidth - 8\tabcolsep) * \real{0.2105}}@{}}
\toprule\noalign{}
\begin{minipage}[b]{\linewidth}\raggedright
Track
\end{minipage} & \begin{minipage}[b]{\linewidth}\raggedleft
nox-mem (strict)
\end{minipage} & \begin{minipage}[b]{\linewidth}\raggedleft
nox-mem (adjacency-2)
\end{minipage} & \begin{minipage}[b]{\linewidth}\raggedleft
Mem0 published F1
\end{minipage} & \begin{minipage}[b]{\linewidth}\raggedleft
note
\end{minipage} \\
\midrule\noalign{}
\endhead
\bottomrule\noalign{}
\endlastfoot
Overall retrieval@10 & \textbf{74.52\%} & 87.10\% & 66.88\% (F1) &
different metric, not a head-to-head \\
Multi-hop retrieval@10 & \textbf{82.21\%} & \textbf{92.91\%} & --- &
--- \\
Single-hop retrieval@10 & 71.40\% & 84.13\% & --- & --- \\
Temporal retrieval@10 & 68.94\% & 82.31\% & --- & --- \\
\end{longtable}
}

The strict retrieval@10 of 74.52\% is not comparable to Mem0's published
end-to-end F1 of 66.88\% (retrieval@10 vs answer-F1 are different
metrics, not a same-corpus head-to-head; for the like-for-like nDCG@10
comparison where Mem0 wins LoCoMo under its native embedder, see §6.3).
This is the retrieval \textbf{ceiling} for the corpus --- the best a
prompt-level F1 push can hope for from candidates retrieved at
top\_k=10. The multi-hop sub-track at 82.21\% strict / 92.91\%
adjacency-2 confirms that multi-hop retrieval on LoCoMo is structurally
easier than on EverMemBench (consistent with the §5.4 corpus-structural
argument).

\paragraph{5.3.2 F1 push --- 51.85\% rank-5 (above Zep / LangMem, below
Mem0
SOTA)}\label{f1-push-51.85-rank-5-above-zep-langmem-below-mem0-sota}

A prompt-level F1 push with the same retrieval pool was measured to
quantify the gap between retrieval ceiling and end-to-end F1. PR \#404.

{\def\LTcaptype{none} % do not increment counter
\begin{longtable}[]{@{}
  >{\raggedright\arraybackslash}p{(\linewidth - 6\tabcolsep) * \real{0.2000}}
  >{\raggedleft\arraybackslash}p{(\linewidth - 6\tabcolsep) * \real{0.2667}}
  >{\centering\arraybackslash}p{(\linewidth - 6\tabcolsep) * \real{0.3333}}
  >{\raggedright\arraybackslash}p{(\linewidth - 6\tabcolsep) * \real{0.2000}}@{}}
\toprule\noalign{}
\begin{minipage}[b]{\linewidth}\raggedright
System
\end{minipage} & \begin{minipage}[b]{\linewidth}\raggedleft
LoCoMo F1
\end{minipage} & \begin{minipage}[b]{\linewidth}\centering
Rank
\end{minipage} & \begin{minipage}[b]{\linewidth}\raggedright
Notes
\end{minipage} \\
\midrule\noalign{}
\endhead
\bottomrule\noalign{}
\endlastfoot
Mem0 SOTA (published) & 66.88\% & 1 & Per Mem0 paper \\
Mem0 (replication) & \textasciitilde60\% & 2 & Range from internal
estimate \\
OpenAI-memory (published) & \textasciitilde55\% & 3 & Estimated from
Mem0 comparison table \\
LangGraph (published) & \textasciitilde52\% & 4 & Estimated \\
\textbf{nox-mem (F1 push)} & \textbf{51.85\%} & \textbf{5} &
\textbf{Above Zep 50.40\% / LangMem 50.21\%} (PR \#404) \\
Zep (replication) & 50.40\% & 6 & Internal measurement \\
LangMem (replication) & 50.21\% & 7 & Internal measurement \\
\end{longtable}
}

nox-mem's F1 push of 51.85\% is rank-5 across the benchmarked systems,
\textbf{above Zep and LangMem} but below Mem0 SOTA. The gap between
retrieval ceiling (74.52\%) and F1 push (51.85\%) is the
\textbf{verbosity gap} --- F1 scoring penalises verbose answers that
contain the correct fact embedded in longer responses. Mem0's
specialized fact-extraction prompts close this gap; nox-mem's
general-purpose retrieval + GPT-4.1-mini prompt does not.

The date-normalization knob shipped in PR \#396 contributed +2.8 pp F1
on temporal sub-track by aligning date format expressions between query
and corpus chunks; this is the largest single tunable knob for LoCoMo
F1.

\paragraph{5.3.3 Path to \textgreater=55\% F1 --- composition
orchestration (Q3)}\label{path-to-55-f1-composition-orchestration-q3}

Wave C ceiling analysis (§5.1.9) demonstrates that retrieval-stage knobs
cannot lift LoCoMo F1 above the verbosity gap. The path to
\textgreater=55\% F1 (rank-3 territory) requires
\textbf{orchestration-stage} mechanisms: prompt-level fact extraction,
iterative refinement (Q3 IterB ReAct), or explicit answer-shaping. §5.5
reports the first measurement on this axis.

\begin{center}\rule{0.5\linewidth}{0.5pt}\end{center}

\subsubsection{5.4 EverMemBench F\_MH paradox ---
resolved}\label{evermembench-f_mh-paradox-resolved}

The triangulation of three independent multi-hop measurements forces a
reframing of the EverMemBench F\_MH absolute number:

{\def\LTcaptype{none} % do not increment counter
\begin{longtable}[]{@{}
  >{\raggedright\arraybackslash}p{(\linewidth - 8\tabcolsep) * \real{0.1765}}
  >{\raggedright\arraybackslash}p{(\linewidth - 8\tabcolsep) * \real{0.1765}}
  >{\raggedleft\arraybackslash}p{(\linewidth - 8\tabcolsep) * \real{0.2353}}
  >{\raggedleft\arraybackslash}p{(\linewidth - 8\tabcolsep) * \real{0.2353}}
  >{\raggedright\arraybackslash}p{(\linewidth - 8\tabcolsep) * \real{0.1765}}@{}}
\toprule\noalign{}
\begin{minipage}[b]{\linewidth}\raggedright
Benchmark
\end{minipage} & \begin{minipage}[b]{\linewidth}\raggedright
Multi-hop metric
\end{minipage} & \begin{minipage}[b]{\linewidth}\raggedleft
nox-mem score
\end{minipage} & \begin{minipage}[b]{\linewidth}\raggedleft
Reader SOTA range
\end{minipage} & \begin{minipage}[b]{\linewidth}\raggedright
Verdict
\end{minipage} \\
\midrule\noalign{}
\endhead
\bottomrule\noalign{}
\endlastfoot
\textbf{MuSiQue dev} & F1 (multi-hop decomposable) & \textbf{58.62\%} &
35.80\% (IRCoT) -- 49.70\% (EX(SA)) & \textbf{nox-mem SOTA} \\
\textbf{HotPotQA dev distractor} & ans\_F1 (multi-hop bridge) &
\textbf{73.37\%} & 65--72\% (DPR+FiD reader SOTA) & \textbf{nox-mem
SOTA} \\
\textbf{LoCoMo dev} & retrieval@10 strict & \textbf{74.52\%} & 66.88\%
(Mem0 answer-F1, different metric) & \textbf{retrieval ceiling; not
comparable to Mem0 F1} \\
\textbf{LoCoMo dev (F1 push)} & F1 (verbosity-sensitive) & 51.85\% &
66.88\% (Mem0 SOTA) & rank-5, verbosity gap \\
\textbf{EverMemBench F\_MH} & strict EM (multi-hop chain on long conv) &
3--7\% (5-batch) & 18.88\% (MemOS Table 4) & \textbf{−13 to −16 pp
gap} \\
\end{longtable}
}

If nox-mem's multi-hop reasoning were structurally weak, the MuSiQue and
HotPotQA SOTA results would not be possible. They demonstrate the
reverse: nox-mem's hybrid retrieval + GPT-4.1-mini reader pipeline is
SOTA on the canonical multi-hop QA benchmarks. The LoCoMo retrieval
ceiling at 74.52\% strict (82.21\% multi-hop sub-track) further confirms
that multi-hop retrieval over long conversations is achievable.

\textbf{The resolution.} The EverMemBench F\_MH 3--7\% number measures a
corpus-structural challenge specific to the EverMemBench task setup:

\begin{enumerate}
\def\labelenumi{\arabic{enumi}.}
\tightlist
\item
  \textbf{Very long conversation chains.} EverMemBench F\_MH questions
  require composing facts across 100+ conversation turns, far longer
  than MuSiQue (\textless=4 paragraphs) or HotPotQA (2 bridge
  paragraphs).
\item
  \textbf{Strict scoring.} EverMemBench F\_MH uses strict exact-match
  against canonical answers; minor wording variations are penalised even
  when the answer is correct. MuSiQue F1 and HotPotQA ans\_F1 are
  partial-credit scores.
\item
  \textbf{Entity-anchor sparsity.} EverMemBench questions often lack
  explicit entity tokens that nox-mem's section/source-type boost
  framework can latch onto. The §5.1.8.3 MAP (bypass-entity) mechanism
  was designed specifically to address this sparsity.
\item
  \textbf{Memory-vs-retrieval mismatch.} EverMemBench is a \emph{memory}
  benchmark with implicit world-state updates; the chunks that answer
  F\_MH questions may not be the chunks that explicit retrieval would
  surface. This is the architectural distinction MemOS optimises for.
\end{enumerate}

\textbf{Implication for Q3 priorities --- refined by D74 (2026-05-31).}
The §5.4 framing shifts Q3 retrieval-mechanism priorities: pure
retrieval-stage knobs (KG, MQ, MAP) cap at \textasciitilde+7.25 pp F\_MH
(Wave C ceiling §5.1.9). Closing the remaining EverMemBench F\_MH gap
requires either (a) orchestration-stage multi-round refinement matching
the long-chain structure (Q3 IterB ReAct), or (b) backbone upgrade
(Backbone Matrix §5.1.10: Gemini-3-flash already narrows the F\_MH gap
meaningfully). Both paths are now empirically validated. The §5.5 Q3
IterC mechanism-class finding confirms that not all orchestration
mechanisms transfer to EverMemBench F\_MH equally (parallel
decomposition vs sequential refinement). The §5.5.2 Q3 IterB ReAct
result (D74) goes further: on the strongest backbone (Gemini-3-flash
bare), multi-round retrieve-reason loop delivers \textbf{+2.01 pp clean
F\_MH lift (8.03\% from 6.02\% bare baseline)} --- exceeding the Wave
A/B/C single-stage retrieval ceiling of 7.25 pp by +0.78 pp standalone.
The paradox is therefore refined rather than dissolved: EverMemBench
F\_MH is still largely a structural property of very long conversation
chains × strict scoring, but MAS orchestration adds \textasciitilde+2 pp
on top of the strongest backbone above the retrieval ceiling --- closing
the gap is no longer purely structural, it now has both a backbone path
and an orchestration path.

\begin{center}\rule{0.5\linewidth}{0.5pt}\end{center}

\subsubsection{5.5 Q3 orchestration --- IterC Self-Ask breakthrough,
IterB ReAct ceiling break, and mechanism-class
distinction}\label{q3-orchestration-iterc-self-ask-breakthrough-iterb-react-ceiling-break-and-mechanism-class-distinction}

Q3 explores orchestration-stage mechanisms above the retrieval ceiling
identified in Wave C (§5.1.9). Two deliverables are reported in this
revision: \textbf{§5.5.1 Q3 IterC} (parallel decomposition via Self-Ask,
F\_HL breakthrough) and \textbf{§5.5.2 Q3 IterB} (sequential refinement
via ReAct, F\_MH retrieval-stage ceiling break on the strongest
backbone). §5.5.3 records the mechanism-class distinction that emerged
from the IterC/IterB contrast. §5.5.4 reports the composability
projection for IterB stacked on top of Wave A/B/C retrieval-stage
mechanisms (pending Q1 validation).

\paragraph{5.5.1 Q3 IterC --- Self-Ask parallel decomposition (F\_HL
+35.84
pp)}\label{q3-iterc-self-ask-parallel-decomposition-f_hl-35.84-pp}

Q3 IterC implements a Self-Ask-style sub-query loop: the generation
model decomposes the query into sub-questions, each sub-question is
retrieved separately, and a final synthesis step composes the answer. PR
\#406.

{\def\LTcaptype{none} % do not increment counter
\begin{longtable}[]{@{}lrrr@{}}
\toprule\noalign{}
Metric & Q3 IterC (5-batch) & Phase H v2 baseline & Δ \\
\midrule\noalign{}
\endhead
\bottomrule\noalign{}
\endlastfoot
\textbf{F\_HL (high-level)} & \textbf{58.52\%} & 22.68\% &
\textbf{+35.84 pp} PASS \\
F\_MH (multi-hop) & \textasciitilde{} baseline & 3.21\% & \textbf{−0.40
pp} (no-lift) \\
Overall & mixed & 51.68\% & --- \\
Cost / latency & \$0.0015/q + 2× LLM call & baseline & added cost \\
\end{longtable}
}

The F\_HL +35.84 pp lift is the \textbf{largest single-mechanism F\_HL
improvement} measured in nox-mem's evaluation history. F\_HL (high-level
synthesis questions) benefits from parallel sub-query decomposition
because the synthesis target itself is a composition over independent
sub-facts --- exactly the mechanism Self-Ask was designed for. The F\_MH
no-lift (−0.40 pp) is also informative: it forced the mechanism-class
distinction documented in §5.5.3.

The Q3 IterC verdict: \textbf{ship opt-in}
(\texttt{NOX\_Q3\_ITERC\_ENABLED=1}) for F\_HL-heavy workloads. NOT
default-enabled (added cost + F\_MH no-lift). The mechanism-class
distinction reframed Q3 IterB ReAct as the leading EverMemBench F\_MH
candidate, which was then validated in §5.5.2.

\paragraph{5.5.2 Q3 IterB ReAct --- F\_MH retrieval-stage ceiling break
on best backbone
(NEW)}\label{q3-iterb-react-f_mh-retrieval-stage-ceiling-break-on-best-backbone-new}

\textbf{Method.} Q3 IterB ReAct (Yao et al.~2022, arxiv:2210.03629) ---
multi-round retrieve-reason loop. The orchestrator
(gemini-2.5-flash-lite, low-cost cheap-class) generates
\texttt{thought\ →\ action\ (search)\ →\ observation} cycles up to 5
rounds (mean 4.25 rounds across the 5-batch set), with the final-answer
backbone (gemini-3-flash-preview, the in-matrix strongest, §5.1.10).
Evaluated on EverMemBench using the canonical 5-batch sequential
protocol (batches 004 / 005 / 010 / 011 / 016, n=3,121 queries). PR
\#419.

\textbf{Headline (dual-baseline honest reporting, per D74 convention
§5).} The Phase H v2 (GPT-4.1-mini) baseline conflates ReAct mechanism
lift with the backbone swap; the gemini-3-flash bare baseline isolates
the clean mechanism effect on the strongest backbone. The ceiling-break
claim load-bears on the latter.

{\def\LTcaptype{none} % do not increment counter
\begin{longtable}[]{@{}
  >{\raggedright\arraybackslash}p{(\linewidth - 8\tabcolsep) * \real{0.1579}}
  >{\raggedleft\arraybackslash}p{(\linewidth - 8\tabcolsep) * \real{0.2105}}
  >{\raggedleft\arraybackslash}p{(\linewidth - 8\tabcolsep) * \real{0.2105}}
  >{\raggedleft\arraybackslash}p{(\linewidth - 8\tabcolsep) * \real{0.2105}}
  >{\raggedleft\arraybackslash}p{(\linewidth - 8\tabcolsep) * \real{0.2105}}@{}}
\toprule\noalign{}
\begin{minipage}[b]{\linewidth}\raggedright
Metric
\end{minipage} & \begin{minipage}[b]{\linewidth}\raggedleft
IterB vs Phase H v2 conflated (GPT-4.1-mini)
\end{minipage} & \begin{minipage}[b]{\linewidth}\raggedleft
Δ conflated
\end{minipage} & \begin{minipage}[b]{\linewidth}\raggedleft
IterB vs gemini-3-flash bare CLEAN
\end{minipage} & \begin{minipage}[b]{\linewidth}\raggedleft
Δ bare HONEST
\end{minipage} \\
\midrule\noalign{}
\endhead
\bottomrule\noalign{}
\endlastfoot
Overall & 51.68\% → 62.70\% & \textbf{+11.02 pp} & 63.28\% → 62.70\% &
−0.58 pp (within ±1.5 pp CI noise) \\
\textbf{F\_MH} & 3.21\% → \textbf{8.03\%} & \textbf{+4.82 pp} & 6.02\% →
\textbf{8.03\%} & \textbf{+2.01 pp} \\
F\_TP & 15.00\% → 33.33\% & +18.33 pp & n/a & n/a \\
F\_HL & 22.68\% → 43.06\% & +20.38 pp & n/a & n/a \\
MA composite & 73.34\% → 84.89\% & +11.55 pp & 88.42\% → 84.89\% & −3.53
pp (borderline) \\
Cost / query & n/a & --- & within \$0.005 (\$0.00295 measured) & PASS \\
\end{longtable}
}

\textbf{Interpretation --- two ship verdicts.} Against the
project-convention Phase H v2 baseline the result clears 4/4 gates
(SHIP\_DEFAULT\_CANDIDATE). Against the gemini-3-flash bare baseline ---
the load-bearing comparison for any ceiling-break claim --- the result
clears 3/4 gates: F\_MH +2.01 pp PASS, Overall within CI noise PASS,
cost within budget PASS, MA composite borderline (−3.53 pp, falls inside
the MA tolerance band but at its lower edge). Final verdict:
\textbf{SHIP\_OPT\_IN} (\texttt{NOX\_Q3\_ITERB\_ENABLED=1}). The opt-in
framing reflects the MA borderline, not the F\_MH or cost finding.

\textbf{Why two baselines.} Reporting only the +4.82 pp F\_MH vs Phase H
v2 would conflate two distinct effects: (i) the backbone swap
GPT-4.1-mini → Gemini-3-flash, which alone delivers +20.73 pp Overall
and +32.74 pp MA composite (§5.1.10, Backbone Matrix), and (ii) the
IterB ReAct multi-round mechanism. The +2.01 pp F\_MH vs gemini-3-flash
bare is the \textbf{clean isolated ReAct effect}, with backbone held
constant. The +0.78 pp by which this clean number exceeds the Wave A/B/C
single-stage retrieval ceiling of 7.25 pp (D69, §5.1.9) is the
load-bearing ceiling-break claim.

\textbf{Mechanism instrumentation (Set E).} The 5-batch run reports
IterB applied to 99.6\% of queries (3,107 of 3,121, zero errors, zero
generation-backbone fallbacks), mean 4.25 rounds with p95=5, 99.5\%
terminated via \texttt{answer} action versus 0.5\% via
\texttt{max\_rounds} exhaustion, and round-2 chunk overlap mean of 0.257
with round-1 (LOW overlap --- ReAct explores new evidence rather than
re-fetching the same chunks, the sweet-spot mechanism profile for
sequential refinement).

\textbf{Ceiling refinement --- D74 vs D69 / D72.} D69 established the
Wave A/B/C single-stage retrieval ceiling at +7.25 pp F\_MH (§5.1.9).
D72 (PR \#410, third revision) framed F\_MH as ``structural challenge of
EverMemBench'' given MuSiQue / HotPotQA / LoCoMo retrieval SOTA. D74
refines this framing: F\_MH is \textbf{still largely structural} (long
chains × cross-session compression × strict EM scoring), but MAS
orchestration via ReAct adds \textbf{+2 pp clean F\_MH on top of the
strongest backbone above the retrieval-stage ceiling}. The paradox is
refined rather than dissolved --- closing the EverMemBench F\_MH gap now
has both a backbone path (Backbone Matrix, §5.1.10) and an orchestration
path (Q3 IterB ReAct, this section), in addition to retrieval-stage
mechanisms (Wave A/B/C, §5.1.8/§5.1.9).

\paragraph{5.5.3 Mechanism-class distinction --- parallel decomposition
vs sequential
refinement}\label{mechanism-class-distinction-parallel-decomposition-vs-sequential-refinement}

The Q3 IterC F\_MH no-lift (−0.40 pp, §5.5.1) and Q3 IterB F\_MH +2.01
pp clean lift (§5.5.2) together establish a mechanism-class distinction
that is sharper than any individual measurement.

{\def\LTcaptype{none} % do not increment counter
\begin{longtable}[]{@{}
  >{\raggedright\arraybackslash}p{(\linewidth - 8\tabcolsep) * \real{0.1765}}
  >{\raggedright\arraybackslash}p{(\linewidth - 8\tabcolsep) * \real{0.1765}}
  >{\raggedright\arraybackslash}p{(\linewidth - 8\tabcolsep) * \real{0.1765}}
  >{\raggedleft\arraybackslash}p{(\linewidth - 8\tabcolsep) * \real{0.2353}}
  >{\raggedleft\arraybackslash}p{(\linewidth - 8\tabcolsep) * \real{0.2353}}@{}}
\toprule\noalign{}
\begin{minipage}[b]{\linewidth}\raggedright
Mechanism class
\end{minipage} & \begin{minipage}[b]{\linewidth}\raggedright
Example
\end{minipage} & \begin{minipage}[b]{\linewidth}\raggedright
Q3 result
\end{minipage} & \begin{minipage}[b]{\linewidth}\raggedleft
F\_MH lift
\end{minipage} & \begin{minipage}[b]{\linewidth}\raggedleft
F\_HL lift
\end{minipage} \\
\midrule\noalign{}
\endhead
\bottomrule\noalign{}
\endlastfoot
Parallel decomposition & Self-Ask, Decomposed prompting & \textbf{Q3
IterC (shipped opt-in)} & no-lift (−0.40 pp) & \textbf{+35.84 pp}
(measured) \\
Sequential refinement & \textbf{ReAct (Yao 2022), Iterative retrieval} &
\textbf{Q3 IterB (shipped opt-in, D74)} & \textbf{+2.01 pp clean (above
ceiling)} & n/a \\
Single-round augmentation & KG path, MQ, MAP & §5.1.8 standalones &
+2.81 to +4.04 pp (capped at Wave C ceiling 7.25 pp) & marginal \\
\end{longtable}
}

\textbf{Why the class matters.} Self-Ask retrieves all sub-questions in
parallel --- appropriate when the synthesis target factors into
independent sub-facts (F\_HL). EverMemBench F\_MH is a
\textbf{sequential dependency} task where each hop depends on the
previous hop's resolved entity; parallel decomposition cannot help
because the second sub-question is not knowable until the first has
resolved. ReAct's \texttt{thought\ →\ action\ →\ observation} loop fits
the sequential dependency structure directly: each round's observation
refines the next thought's action. The +2.01 pp clean F\_MH lift on
Gemini-3-flash bare confirms the class-fit empirically.

\textbf{Practical reading.} Workloads with high F\_HL share benefit from
IterC; workloads with high F\_MH share benefit from IterB. Both ship
opt-in. Routing a query to the appropriate orchestration mechanism
(parallel vs sequential) is an open Q1 work item.

\paragraph{5.5.4 Empirical per-backbone × per-knob composability matrix
(Wave 2 closure, replaces D74
projection)}\label{empirical-per-backbone-per-knob-composability-matrix-wave-2-closure-replaces-d74-projection}

The D74 revision of this section contained a projection table assuming
Wave A/B/C retrieval-side lifts measured on gpt-4.1-mini transfer
additively to the Gemini-3-flash backbone. Wave 2 (2026-05-31, D75, PRs
\#423--\#425) empirically tested this assumption. The projection is
replaced by the measured matrix below.

\textbf{Empirically measured per-backbone × per-knob F\_MH matrix
(5-batch CLEAN, n=3,121, batches 004/005/010/011/016):}

{\def\LTcaptype{none} % do not increment counter
\begin{longtable}[]{@{}
  >{\raggedright\arraybackslash}p{(\linewidth - 10\tabcolsep) * \real{0.1429}}
  >{\raggedright\arraybackslash}p{(\linewidth - 10\tabcolsep) * \real{0.1429}}
  >{\raggedleft\arraybackslash}p{(\linewidth - 10\tabcolsep) * \real{0.1905}}
  >{\centering\arraybackslash}p{(\linewidth - 10\tabcolsep) * \real{0.2381}}
  >{\raggedright\arraybackslash}p{(\linewidth - 10\tabcolsep) * \real{0.1429}}
  >{\raggedright\arraybackslash}p{(\linewidth - 10\tabcolsep) * \real{0.1429}}@{}}
\toprule\noalign{}
\begin{minipage}[b]{\linewidth}\raggedright
Backbone
\end{minipage} & \begin{minipage}[b]{\linewidth}\raggedright
Knob
\end{minipage} & \begin{minipage}[b]{\linewidth}\raggedleft
F\_MH lift (5-batch CI)
\end{minipage} & \begin{minipage}[b]{\linewidth}\centering
Gate +1.5 pp
\end{minipage} & \begin{minipage}[b]{\linewidth}\raggedright
Validated?
\end{minipage} & \begin{minipage}[b]{\linewidth}\raggedright
Source
\end{minipage} \\
\midrule\noalign{}
\endhead
\bottomrule\noalign{}
\endlastfoot
gpt-4.1-mini & KG path & +2.81 pp (CI {[}2.11, 9.93{]}) & PASS & YES &
PR \#379 \\
Gemini-3-flash & KG path & −0.01 pp (CI {[}3.00, 9.04{]}) & FAIL
NO-REPLICATE & 0\% transfer & PR \#423 \\
gpt-4.1-mini & AC (threshold=5) & +2.01 pp (CI {[}1.06, 9.39{]}) & PASS
marginal & YES & PR \#381 \\
Gemini-3-flash & AC (threshold=5) & +0.81 pp (CI {[}4.62, 9.03{]}) &
FAIL NO-REPLICATE & 40\% transfer & PR \#424 \\
gpt-4.1-mini & MQ standalone & +3.61 pp & PASS & YES & PR \#385 \\
Gemini-3-flash & MQ standalone & +1.21 pp (CI {[}4.99, 9.48{]}) & FAIL
borderline & 34\% transfer & PR \#425 \\
Gemini-3-flash & IterB ReAct & +2.01 pp (bare CLEAN) & PASS &
\textbf{ONLY VALIDATED} & PR \#419 \\
Gemini-3-flash & IterB + Wave C triple & INDETERMINATE & --- &
infra-bound & PR \#426 (D76)¹ \\
\end{longtable}
}

\begin{quote}
¹ \textbf{D76 capstone deferral footnote (§5.5.8):} The IterB + Wave C
triple composability test (PR \#426) was aborted due to Hostinger VPS
CPU steal 51--97\% sustained, not due to scientific failure. Batch 004
(n=49) preserved. 5-batch threshold not reached. Outcome is
INDETERMINATE; composability claim is neither confirmed nor refuted.
Capstone deferred to future stable infrastructure with dedicated CPU
SLO.
\end{quote}

\textbf{Headline numbers.} The 3-knob sum on Gemini-3-flash (KG −0.01 pp
+ AC +0.81 pp + MQ +1.21 pp) = \textbf{+2.01 pp aggregate} = 24\% of the
D74 pessimistic projection of +8.43 pp.~All three individual knob CIs
fully overlap the Gemini-3-flash baseline (6.02\%), meaning no single
knob clears statistical significance at the +1.5 pp gate. IterB ReAct
standalone (+2.01 pp clean, §5.5.2) equals the entire 3-knob aggregate
while being structurally distinct --- an orchestration-stage mechanism
rather than retrieval-stage augmentation.

\textbf{Corrected composability landscape.} The original D74 projection
table (IterB + Wave C triple → \textasciitilde12.07\% F\_MH =
\textasciitilde41\% MemOS gap closure) assumed backbone-invariant
transfer of all knob lifts. That assumption is empirically refuted on
Gemini-3-flash for all three tested retrieval-stage knobs. The current
empirically supported picture:

{\def\LTcaptype{none} % do not increment counter
\begin{longtable}[]{@{}
  >{\raggedright\arraybackslash}p{(\linewidth - 6\tabcolsep) * \real{0.2143}}
  >{\raggedleft\arraybackslash}p{(\linewidth - 6\tabcolsep) * \real{0.2857}}
  >{\raggedleft\arraybackslash}p{(\linewidth - 6\tabcolsep) * \real{0.2857}}
  >{\raggedright\arraybackslash}p{(\linewidth - 6\tabcolsep) * \real{0.2143}}@{}}
\toprule\noalign{}
\begin{minipage}[b]{\linewidth}\raggedright
Configuration
\end{minipage} & \begin{minipage}[b]{\linewidth}\raggedleft
F\_MH
\end{minipage} & \begin{minipage}[b]{\linewidth}\raggedleft
Closure of MemOS F\_MH gap (\textasciitilde18.88 pp)
\end{minipage} & \begin{minipage}[b]{\linewidth}\raggedright
Status
\end{minipage} \\
\midrule\noalign{}
\endhead
\bottomrule\noalign{}
\endlastfoot
Bare Gemini-3-flash & 6.02\% & baseline & measured \\
\textbf{IterB ReAct standalone on bare (this work, §5.5.2)} &
\textbf{8.03\%} & \textbf{\textasciitilde7\%} & \textbf{measured} \\
IterB + retrieval-stage knobs (aggregate upper bound) &
\textasciitilde8--9\% & \textasciitilde10--15\% & bounded estimate \\
IterB + Wave C triple (orchestration composability) & INDETERMINATE &
INDETERMINATE & D76 deferred \\
\end{longtable}
}

\begin{center}\rule{0.5\linewidth}{0.5pt}\end{center}

\paragraph{5.5.5 Wave 2 --- Single-stage knob backbone-portability
refinement
(D75)}\label{wave-2-single-stage-knob-backbone-portability-refinement-d75}

\textbf{Setup.} Wave 2 Phase 1 (R0 sanity, PR \#423) and Phase 1.5 (AC +
MQ re-baseline, PRs \#424 + \#425) re-ran all three principal Lab Q1
single-stage retrieval knobs on the Gemini-3-flash backbone (D70,
§5.1.10) using the identical 5-batch CLEAN sequential protocol (n=3,121,
batches 004/005/010/011/016). The motivation: D74 composability
projection assumed knob lifts measured on gpt-4.1-mini were
backbone-invariant. R0 tested this assumption for KG path before
dispatching the full composability matrix run.

\textbf{3-knob NO-REPLICATE pattern.} Three independent retrieval-stage
knobs all show the same structural pattern:

{\def\LTcaptype{none} % do not increment counter
\begin{longtable}[]{@{}
  >{\raggedright\arraybackslash}p{(\linewidth - 8\tabcolsep) * \real{0.1667}}
  >{\raggedleft\arraybackslash}p{(\linewidth - 8\tabcolsep) * \real{0.2222}}
  >{\raggedleft\arraybackslash}p{(\linewidth - 8\tabcolsep) * \real{0.2222}}
  >{\raggedleft\arraybackslash}p{(\linewidth - 8\tabcolsep) * \real{0.2222}}
  >{\raggedright\arraybackslash}p{(\linewidth - 8\tabcolsep) * \real{0.1667}}@{}}
\toprule\noalign{}
\begin{minipage}[b]{\linewidth}\raggedright
Knob
\end{minipage} & \begin{minipage}[b]{\linewidth}\raggedleft
gpt-4.1-mini F\_MH
\end{minipage} & \begin{minipage}[b]{\linewidth}\raggedleft
Gemini-3-flash F\_MH
\end{minipage} & \begin{minipage}[b]{\linewidth}\raggedleft
Transfer rate
\end{minipage} & \begin{minipage}[b]{\linewidth}\raggedright
95\% CI on Gemini
\end{minipage} \\
\midrule\noalign{}
\endhead
\bottomrule\noalign{}
\endlastfoot
KG path (R0, PR \#423) & +2.81 pp & \textbf{−0.01 pp} & 0\% & {[}3.00,
9.04{]} \\
AC threshold=5 (PR \#424) & +2.01 pp & \textbf{+0.81 pp} & 40\% &
{[}4.62, 9.03{]} \\
MQ standalone (PR \#425) & +3.61 pp & \textbf{+1.21 pp} & 34\% &
{[}4.99, 9.48{]} \\
\textbf{3-knob sum} & \textbf{+8.43 pp} & \textbf{+2.01 pp} &
\textbf{24\% aggregate} & --- \\
\end{longtable}
}

All three Gemini-3-flash CIs fully overlap the bare baseline (6.02\%).
The pattern is consistent across knobs of different mechanism families
(entity-walk SQL, heuristic query routing, LLM sub-query decomposition),
indicating a structural backbone-conditional property rather than a
knob-specific failure.

\textbf{Mechanism interpretation.} The hypothesis consistent with all
three observations: Wave A knobs were designed to compensate for
context-bottleneck weaknesses of gpt-4.1-mini --- smaller context
window, weaker filtering, lower context utilization per token.
Gemini-3-flash's larger context window and stronger native context
utilization saturates the compensation signal that these knobs provide,
yielding diminishing marginal returns. KG path (0\% transfer) is the
extreme case: Gemini already processes the relevant entity graph context
from retrieved chunks without requiring explicit vault-fact injection.
AC and MQ show partial transfer (34--40\%) because their mechanisms
involve multi-round or breadth-expansion effects that provide some
marginal diversity even for stronger backbones, but not enough to clear
the statistical gate.

\textbf{Generalization principle.} Any retrieval-stage knob lift of the
form ``Knob X delivers +N pp on backbone Y'' is backbone-conditional.
Cross-backbone generalization requires explicit re-baseline. As
backbones strengthen (Claude Opus 4.7, GPT-5, Gemini 4), retrieval-stage
compensation mechanisms may show further transfer-rate attenuation.
Future composability projections should re-baseline each knob on the
target deployment backbone before projecting stacked effects.

\begin{center}\rule{0.5\linewidth}{0.5pt}\end{center}

\paragraph{5.5.6 Wave 2 --- MQ multi-axis backbone-conditional behavior
(sub-finding, PR
\#425)}\label{wave-2-mq-multi-axis-backbone-conditional-behavior-sub-finding-pr-425}

The MQ re-baseline (PR \#425) revealed a sub-finding that is
paper-worthy independent of the NO-REPLICATE verdict: MQ exhibits
\textbf{inverse backbone-portability across metric axes}.

{\def\LTcaptype{none} % do not increment counter
\begin{longtable}[]{@{}
  >{\raggedright\arraybackslash}p{(\linewidth - 6\tabcolsep) * \real{0.2500}}
  >{\raggedright\arraybackslash}p{(\linewidth - 6\tabcolsep) * \real{0.2500}}
  >{\raggedright\arraybackslash}p{(\linewidth - 6\tabcolsep) * \real{0.2500}}
  >{\raggedright\arraybackslash}p{(\linewidth - 6\tabcolsep) * \real{0.2500}}@{}}
\toprule\noalign{}
\begin{minipage}[b]{\linewidth}\raggedright
Metric axis
\end{minipage} & \begin{minipage}[b]{\linewidth}\raggedright
gpt-4.1-mini result
\end{minipage} & \begin{minipage}[b]{\linewidth}\raggedright
Gemini-3-flash result
\end{minipage} & \begin{minipage}[b]{\linewidth}\raggedright
Direction
\end{minipage} \\
\midrule\noalign{}
\endhead
\bottomrule\noalign{}
\endlastfoot
F\_MH lift & +3.61 pp (biggest single retrieval knob) & +1.21 pp
(borderline, CI overlap) & Attenuates \\
MA composite & −1.38 pp (regression) & \textbf{+0.12 pp (preserved)} &
\textbf{Flips sign} \\
MA\_U (Memory Update) & modest & \textbf{+3.10 pp} (strongest MA gain in
Wave 2) & Inverts entirely \\
\end{longtable}
}

On gpt-4.1-mini, MQ sub-query decomposition multiplies retrieval breadth
but introduces noise that the backbone cannot fully filter ---
manifesting as MA composite regression. On Gemini-3-flash with stronger
filtering and broader context integration, the wider retrieval pool from
MQ sub-queries is interpretable rather than noisy, yielding MA\_U
improvement (Unrelated detection benefits from additional diversity in
retrieved context).

\textbf{Implication.} Per-knob evaluation on a single metric axis (F\_MH
alone) can hide compensating effects on orthogonal dimensions.
Retrieval-stage mechanisms with multi-factor effect profiles (knob
benefits on dimension A, costs dimension B on backbone X; costs A but
benefits B on backbone Y) require multi-axis backbone-conditional
reporting. The gpt-4.1-mini measurements in §5.1.8 remain valid for that
backbone but should not be assumed to represent the MA dimension on
stronger backbones.

\begin{center}\rule{0.5\linewidth}{0.5pt}\end{center}

\paragraph{5.5.7 Architectural composability vs mechanism
composability}\label{architectural-composability-vs-mechanism-composability}

Wave 2 Phase 2 setup (PR \#426, capstone) exposed a third composability
requirement independent of backbone-portability: \textbf{architectural
composability} between orchestration-stage mechanisms (IterB ReAct) and
retrieval-stage mechanisms (Wave A knobs).

\textbf{Code evidence
(\texttt{eval/evermembench/adapter\_nox\_mem.py}).} The PR \#419 IterB
adapter contains explicit guards at three locations:

\begin{Shaded}
\begin{Highlighting}[]
\CommentTok{\# Line 2736 — MQ short{-}circuit}
\ControlFlowTok{if} \KeywordTok{not}\NormalTok{ iterb\_used\_path:}
    \CommentTok{\# ... MQ sub{-}query decomposition + RRF fusion logic ...}

\CommentTok{\# Line 2906 — KG path short{-}circuit}
\ControlFlowTok{if} \KeywordTok{not}\NormalTok{ iterb\_used\_path:}
    \CommentTok{\# ... KG entity extract + 1{-}hop walk + vault{-}fact injection ...}

\CommentTok{\# Line 3063 — cross{-}encoder rerank short{-}circuit}
\ControlFlowTok{if} \KeywordTok{not}\NormalTok{ iterb\_used\_path:}
    \CommentTok{\# ... bge{-}reranker{-}v2{-}m3 cross{-}encoder rerank ...}
\end{Highlighting}
\end{Shaded}

The \texttt{iterb\_used\_path} flag is set when IterB's ReAct loop fires
on a query. Each Wave A knob checks this flag and skips itself if IterB
took the path. Setting \texttt{NOX\_ADAPTER\_MODE=phaseTriple} combined
with \texttt{NOX\_ITERB\_ENABLED=1} does \textbf{not} produce a composed
IterB + Wave C triple system --- IterB takes exclusive precedence and
phaseTriple stages are bypassed entirely.

\textbf{Design rationale (reconstructed from D74 intent).} IterB ReAct
per-round retrieval already uses the full hybrid stack (FTS5 + vec +
RRF). Adding KG + MQ + MAP per ReAct round would multiplicatively expand
per-round cost (×N stages × 4.25 mean rounds) without empirically
validated additivity. The conservative default --- exclusive operation
with explicit short-circuits --- was the rational design choice at D74
time.

\textbf{Scientific implication.} D74's composability projection (IterB +
Wave C triple → \textasciitilde12.07\% F\_MH) implicitly assumed
architectural composability. The code evidence shows that assumption was
\textbf{false by design} --- the system would have needed an explicit
code patch to test it. This demonstrates a general principle:
orchestration-stage mechanisms designed without forward-looking
composability planning create silent architectural locks discoverable
only by empirical code-level inspection. The lock is not a bug; it is a
design decision with sound rationale. But it invalidated the
composability projection as stated.

\textbf{Partial composability test (PR \#426 capstone design).} The Wave
2 capstone agent patched 2 of 3 guards: KG vault-fact injection (line
2906, removed --- KG facts injected per ReAct round) and cross-encoder
rerank (line 3063, removed --- reranks IterB's merged candidate pool).
The MQ guard (line 2736) was deliberately kept because IterB ReAct
sub-queries are semantically equivalent to MQ decomposition; composing
both would double-decompose without mechanistic benefit. This partial
composability test was the object of the D76 capstone run; the
infrastructure abort (§5.5.8) means the result remains INDETERMINATE.

\textbf{Future research recommendation.} When designing new
orchestration-stage mechanisms, specify upfront whether they should
compose with or short-circuit existing mechanisms. Document the
integration choice in the spec PR. This prevents discovering
composability locks post-implementation via code archaeology --- and
prevents composability projection errors in interim paper revisions.

\begin{center}\rule{0.5\linewidth}{0.5pt}\end{center}

\paragraph{5.5.8 Wave 2 Capstone --- D76 infrastructure abort
(INDETERMINATE, not scientific
failure)}\label{wave-2-capstone-d76-infrastructure-abort-indeterminate-not-scientific-failure}

The Wave 2 Phase 2 Capstone (PR \#426 draft, IterB + KG + rerank
composability, 2-guard patch per §5.5.7) was dispatched on the
production Hostinger VPS on 2026-05-31. After 48 hours elapsed and
\textasciitilde\$20--25 spent, the bench was aborted due to Hostinger
anti-abuse CPU throttling, not due to scientific hypothesis failure.

\textbf{Infrastructure failure timeline.}

{\def\LTcaptype{none} % do not increment counter
\begin{longtable}[]{@{}lrl@{}}
\toprule\noalign{}
Measurement window & CPU steal & State \\
\midrule\noalign{}
\endhead
\bottomrule\noalign{}
\endlastfoot
Pre-second reboot & 96.93\% & critical \\
Immediately post-second reboot & 8.54\% & brief recovery (8 min) \\
30 min post-reboot & 21.03\% & degrading \\
Sustained working state & 51--71\% & oscillating \\
Bench running (ONNX rerank active) & 51--97\% & throttled \\
\end{longtable}
}

Mitigation attempted: openclaw service disable, taskset CPU pinning
(cores 0--3 eval / 4--5 API), ORT/OMP/MKL/OpenBLAS thread caps to 2,
search timeout extension 120 s → 600 s, concurrency reduction 3 → 1, two
VPS reboots. None achieved sustained CPU steal below 30\% under bench
load. Mathematical impossibility under sustained throttle: 20 retries ×
(600 s + 300 s) = 5 h max per query × 50 questions × 4 batches = 1,000 h
ceiling. Batch 005 ran 23 h with 0/50 questions completed.

\textbf{Distinction: infrastructure abort is not scientific failure.}
The distinction is load-bearing for interpreting this result:

\begin{itemize}
\tightlist
\item
  A \emph{scientific failure} means the hypothesis was tested and the
  data refuted it --- a publishable negative result.
\item
  An \emph{infrastructure abort} means the hypothesis was not testable
  in the current environment --- the outcome is INDETERMINATE, neither
  confirming nor refuting the hypothesis.
\end{itemize}

The capstone abort falls in the second category. The hypothesis (IterB +
KG + rerank compose on Gemini-3-flash for F\_MH gain) remains
scientifically open. Batch 004 (n=49 questions, completed
pre-second-reboot) is preserved at
\texttt{/root/.openclaw/evermembench-runs/capstone-iterB-triple-004-1780260019/analysis.txt}
for future re-run but does not constitute valid 5-batch evidence alone.

\textbf{Deferred infrastructure requirement.} Completing the capstone
requires a dedicated CPU plan with a guaranteed CPU SLO --- ONNX
cross-encoder rerank (bge-reranker-v2-m3) is CPU-bound; shared VPS
infrastructure with host-level anti-abuse scanning is insufficient for
sustained heavy ONNX workloads. The capstone is deferred to Q1+ on
stable infrastructure (dedicated CPU plan or alternate provider).

\textbf{Wave 2 scientific output.} Despite the capstone abort, Wave 2
delivers five paper-worthy findings: (1) D74 IterB +2.01 pp clean F\_MH
lift on Gemini-3-flash (§5.5.2); (2) D75 3-knob NO-REPLICATE
backbone-conditional pattern (§5.5.5); (3) MQ MA backbone flip
sub-finding (§5.5.6); (4) architectural composability lock discovery
(§5.5.7); (5) this D76 honest infrastructure framing (§5.5.8). The 12
SOTA-tier dimensions documented in §5.1--§5.7 are unaffected by the
capstone outcome.

\begin{center}\rule{0.5\linewidth}{0.5pt}\end{center}

\subsubsection{5.6 Cross-bench validation --- LongMemEval
(n=300)}\label{cross-bench-validation-longmemeval-n300}

\textbf{Config:} Phase D production config (FTS5 + Gemini-embedding-001
+ RRF, rerank OFF, top\_k=20), GPT-4.1-mini backbone, Gemini-2.5-flash
judge, oracle session retrieval, stratified n=300 queries. PR \#378.

{\def\LTcaptype{none} % do not increment counter
\begin{longtable}[]{@{}
  >{\raggedright\arraybackslash}p{(\linewidth - 2\tabcolsep) * \real{0.4286}}
  >{\raggedleft\arraybackslash}p{(\linewidth - 2\tabcolsep) * \real{0.5714}}@{}}
\toprule\noalign{}
\begin{minipage}[b]{\linewidth}\raggedright
Metric
\end{minipage} & \begin{minipage}[b]{\linewidth}\raggedleft
Score
\end{minipage} \\
\midrule\noalign{}
\endhead
\bottomrule\noalign{}
\endlastfoot
nDCG@10 (oracle retrieval) & \textbf{1.0000} (Wilson 95\% lower bound
0.9872) \\
Recall@10 & 1.0000 \\
Task accuracy (n=201 judged) & \textbf{68.16\%} (Wilson 95\% CI
{[}0.6143, 0.7421{]}) \\
\end{longtable}
}

\textbf{Per-category breakdown --- fingerprint consistency with
EverMemBench:}

{\def\LTcaptype{none} % do not increment counter
\begin{longtable}[]{@{}
  >{\raggedright\arraybackslash}p{(\linewidth - 6\tabcolsep) * \real{0.2308}}
  >{\raggedleft\arraybackslash}p{(\linewidth - 6\tabcolsep) * \real{0.3077}}
  >{\raggedright\arraybackslash}p{(\linewidth - 6\tabcolsep) * \real{0.2308}}
  >{\raggedright\arraybackslash}p{(\linewidth - 6\tabcolsep) * \real{0.2308}}@{}}
\toprule\noalign{}
\begin{minipage}[b]{\linewidth}\raggedright
Category
\end{minipage} & \begin{minipage}[b]{\linewidth}\raggedleft
LongMemEval score
\end{minipage} & \begin{minipage}[b]{\linewidth}\raggedright
Strength
\end{minipage} & \begin{minipage}[b]{\linewidth}\raggedright
Matches EverMemBench pattern
\end{minipage} \\
\midrule\noalign{}
\endhead
\bottomrule\noalign{}
\endlastfoot
single-session-assistant & 87.10\% & STRONG & PASS matches F\_SH WIN \\
single-session-user & 86.67\% & STRONG & PASS matches F\_SH WIN \\
knowledge-update & 82.05\% & STRONG & PASS matches MA\_U WIN \\
abstention & 82.61\% & STRONG & (no EverMemBench equiv --- unique
strength) \\
multi-session & 55.81\% & moderate & PASS matches F\_MH gap \\
temporal-reasoning & 54.76\% & moderate & PASS matches F\_TP gap \\
single-session-preference & 31.25\% (n=16, wide CI) & weak & (preference
handling weak) \\
\end{longtable}
}

The per-category fingerprint is \textbf{identical} to EverMemBench Phase
D + H v2 results: strong on single-context factual + abstention +
knowledge update; moderate on multi-session reasoning + temporal
sequencing; weak on preference handling. This cross-bench consistency
(two orthogonal benchmark distributions, different eval sets, different
judges) confirms that nox-mem's strengths and weaknesses are
\textbf{structural properties} of the retrieval architecture, not
benchmark-specific tuning artifacts.

The sanitize fix (\texttt{unicode-aware-sanitize-for-fts5}) is validated
cross-bench: nDCG@10 improved from 0.9126 (Q2 baseline, pre-fix) to
1.0000 (+9.6 pp, oracle ceiling). This +9.6 pp delta is consistent in
magnitude with the Q2 LongMemEval pre-fix measurement, confirming the
fix's impact is not corpus-specific.

\textbf{Caveat:} oracle session retrieval is an upper-ceiling
measurement. Comparison to gbrain (97.6\% nDCG@10 on LongMemEval-S)
requires the \texttt{s\_cleaned} non-oracle follow-up (Lab Q1 priority,
not yet run). The win claim is on \textbf{per-category task accuracy and
abstention handling}, not on the oracle nDCG@10.

\begin{center}\rule{0.5\linewidth}{0.5pt}\end{center}

\subsubsection{5.7 Operational characteristics --- latency, cost, and
footprint
(self-hosted)}\label{operational-characteristics-latency-cost-and-footprint-self-hosted}

The canonical production boost stack
(\texttt{section\_boost\ ×\ source\_type\_boost\ (Hard\ Mutex,\ query\_entity\_count\ \textless{}=\ 2)\ ×\ salience\ v2\ additive})
has been deployed since 2026-05-21 via systemd environment drop-in.
Beyond accuracy, the production deployment establishes a unique
competitive position on three operational dimensions:

\paragraph{5.7.1 Latency --- sub-10 ms KG
path}\label{latency-sub-10-ms-kg-path}

{\def\LTcaptype{none} % do not increment counter
\begin{longtable}[]{@{}
  >{\raggedright\arraybackslash}p{(\linewidth - 8\tabcolsep) * \real{0.1667}}
  >{\raggedleft\arraybackslash}p{(\linewidth - 8\tabcolsep) * \real{0.2222}}
  >{\raggedleft\arraybackslash}p{(\linewidth - 8\tabcolsep) * \real{0.2222}}
  >{\raggedleft\arraybackslash}p{(\linewidth - 8\tabcolsep) * \real{0.2222}}
  >{\raggedright\arraybackslash}p{(\linewidth - 8\tabcolsep) * \real{0.1667}}@{}}
\toprule\noalign{}
\begin{minipage}[b]{\linewidth}\raggedright
Path
\end{minipage} & \begin{minipage}[b]{\linewidth}\raggedleft
p50
\end{minipage} & \begin{minipage}[b]{\linewidth}\raggedleft
p95
\end{minipage} & \begin{minipage}[b]{\linewidth}\raggedleft
p99
\end{minipage} & \begin{minipage}[b]{\linewidth}\raggedright
Notes
\end{minipage} \\
\midrule\noalign{}
\endhead
\bottomrule\noalign{}
\endlastfoot
\textbf{KG path (entity-walk)} & \textbf{2.5 ms} & \textasciitilde7 ms &
\textasciitilde14 ms & SQL + regex over \texttt{kg\_relations}, no LLM
call (PR \#403) \\
Hybrid search (FTS5 + dense + RRF, no rerank) & \textasciitilde940 ms &
\textasciitilde2,342 ms & \textasciitilde2,523 ms & Gemini-embedding-001
query dominates (\textasciitilde800 ms) \\
Hybrid + cross-encoder rerank (MiniLM) & +3,700 ms p50 & --- & --- &
Opt-in, exploratory mode \\
\end{longtable}
}

The 2.5 ms p50 KG path is in the sub-10 ms class --- no published
competitor reports retrieval latency in this band. Zep markets
\textless100 ms p50, unverified independently; Mem0 Cloud and MemOS
deployments are multi-service architectures with network hops typically
in the 100--500 ms range. nox-mem's single-process embedded architecture
(better-sqlite3 + sqlite-vec in-process) eliminates network and IPC
overhead entirely. \textbf{Re-validated 2026-06-15} on the 70.7k-chunk
production corpus: the KG path (\texttt{/api/kg/path}, real entity pair,
n=10) measured p50 = 2.9 ms / p95 = 5.7 ms, confirming the sub-10 ms
class; the Gemini-embedding-dominated standard hybrid path measured p50
= 653 ms / p95 = 706 ms (n=20) --- up from 529 ms as the corpus grew 69k
→ 70.7k chunks and reflecting Gemini API round-trip variance, which
reinforces (rather than weakens) the case for the local KG path on
latency-sensitive workloads.

\paragraph{5.7.2 Cost --- \$0/query KG path; hybrid vs managed SaaS is
list-price, not
like-for-like}\label{cost-0query-kg-path-hybrid-vs-managed-saas-is-list-price-not-like-for-like}

{\def\LTcaptype{none} % do not increment counter
\begin{longtable}[]{@{}
  >{\raggedright\arraybackslash}p{(\linewidth - 6\tabcolsep) * \real{0.2000}}
  >{\raggedleft\arraybackslash}p{(\linewidth - 6\tabcolsep) * \real{0.2667}}
  >{\raggedleft\arraybackslash}p{(\linewidth - 6\tabcolsep) * \real{0.2667}}
  >{\raggedleft\arraybackslash}p{(\linewidth - 6\tabcolsep) * \real{0.2667}}@{}}
\toprule\noalign{}
\begin{minipage}[b]{\linewidth}\raggedright
Component
\end{minipage} & \begin{minipage}[b]{\linewidth}\raggedleft
nox-mem
\end{minipage} & \begin{minipage}[b]{\linewidth}\raggedleft
Mem0 Cloud (published pricing)
\end{minipage} & \begin{minipage}[b]{\linewidth}\raggedleft
Ratio
\end{minipage} \\
\midrule\noalign{}
\endhead
\bottomrule\noalign{}
\endlastfoot
Retrieval API cost (KG path) & \textbf{\$0.00} & \$0.001/query (est.
embedding + retrieval) & \textbf{effectively free} \\
Retrieval API cost (hybrid w/ Gemini embedding) & \$0.0000015/query &
\textasciitilde\$0.001/query (est.) & \textbf{\textasciitilde667×
cheaper} \\
Ingest API cost (per chunk) & \$0.00 (local) & varies & n/a \\
Total cost per 1M queries (hybrid) & \$1.50 & \$1,000 (est.) &
\textasciitilde667× \\
\end{longtable}
}

The KG path achieves \textbf{\$0 per query} because the entity-walk uses
only local SQL + regex with no LLM call (re-confirmed 2026-06-15). The
hybrid path costs \textbf{\$0.0000015/query} (gemini-embedding-001 at
\$0.15/1M input tokens --- a Feb-2026 increase from \$0.13/1M --- ×
\textasciitilde10 tokens/query). Against an \emph{estimated} Mem0 Cloud
per-query rate of \textasciitilde\$0.001 this is
\textbf{\textasciitilde667× cheaper} (revised down from the 769× figure
as Gemini embedding pricing rose). We lead with the unconditional claim
--- \textbf{\$0/query KG path and zero marginal retrieval cost} --- and
treat the multiplier as secondary: Mem0 is sold as a subscription (free
1K calls/month, then \$19--\$249/month) with no published per-call
overage, so the denominator is an explicit modeling assumption, not a
quoted price.

\paragraph{5.7.3 Footprint --- 399 MB RSS, single-process,
self-hosted}\label{footprint-399-mb-rss-single-process-self-hosted}

{\def\LTcaptype{none} % do not increment counter
\begin{longtable}[]{@{}
  >{\raggedright\arraybackslash}p{(\linewidth - 6\tabcolsep) * \real{0.2143}}
  >{\raggedleft\arraybackslash}p{(\linewidth - 6\tabcolsep) * \real{0.2857}}
  >{\raggedleft\arraybackslash}p{(\linewidth - 6\tabcolsep) * \real{0.2857}}
  >{\raggedright\arraybackslash}p{(\linewidth - 6\tabcolsep) * \real{0.2143}}@{}}
\toprule\noalign{}
\begin{minipage}[b]{\linewidth}\raggedright
Scaling
\end{minipage} & \begin{minipage}[b]{\linewidth}\raggedleft
Idle RSS
\end{minipage} & \begin{minipage}[b]{\linewidth}\raggedleft
10× concurrent
\end{minipage} & \begin{minipage}[b]{\linewidth}\raggedright
Notes
\end{minipage} \\
\midrule\noalign{}
\endhead
\bottomrule\noalign{}
\endlastfoot
nox-mem-api process & \textbf{399 MB} & +15 MB (= \textasciitilde414 MB)
& better-sqlite3 + sqlite-vec single-process \\
Scaling pattern & flat & quasi-flat & No per-request memory blow-up \\
\end{longtable}
}

Self-hosted single-process means no multi-container orchestration, no
Postgres/Redis/Chroma sidecars, no per-tenant container overhead. The
399 MB idle footprint runs on a \$5/month VPS tier. Mem0 / Zep / Letta
canonical deployments require \textgreater=3 services (API + DB + vector
store) with combined RSS typically in the 1.5--3 GB range.

\paragraph{5.7.4 Observability layer}\label{observability-layer}

The F10 observability layer (Phase A:
\texttt{/observability/health.html}; Phase B:
\texttt{/observability/evals.html}) renders the full G3→G10d ablation
trajectory in real time over Chart.js, with gate annotations for D43,
D48, D51, D67, D68, and D69. Three rollback paths are documented
(conditional layer only, full mutex, drop-in removal), each executable
in under five minutes.

\begin{center}\rule{0.5\linewidth}{0.5pt}\end{center}

\subsubsection{5.8 Methodology, 5-batch protocol, and honest
caveats}\label{methodology-5-batch-protocol-and-honest-caveats}

\paragraph{5.8.1 5-batch + 95\% CI canonical
methodology}\label{batch-95-ci-canonical-methodology}

All EverMemBench claims in §5.1.5--§5.1.10 use the \textbf{5-batch
canonical protocol} (PR \#371 DECISIONS + PR \#376
\texttt{eval/lib/aggregate\_5batch.py}):

\begin{itemize}
\tightlist
\item
  \textbf{5-batch set:} batches 004, 005, 010, 011, 016 (n
  \textasciitilde{} 120--250 queries each, total n=3,121)
\item
  \textbf{Aggregate metric:} mean across 5 batches per category
\item
  \textbf{CI:} 95\% confidence interval via t-distribution (n=5),
  reported as {[}lower, upper{]}
\item
  \textbf{Claim threshold:} improvement claimed only when CI lower bound
  exceeds baseline mean
\end{itemize}

Single-batch gates were the prior protocol; they are now explicitly
deprecated for any ship/reject decision.

\paragraph{5.8.2 Why single-batch gates overstate effects
3--6×}\label{why-single-batch-gates-overstate-effects-36}

The risk of single-batch measurement is concrete: Phase G batch 004
reported F\_MH +8.00 pp (labelled ``breakthrough''); the 5-batch reality
was +1.61 pp --- a 5× overstatement. Phase H v2 batch 004 reported
+11.60 pp overall vs MemOS; the 5-batch reality was +9.13 pp (1.27×
overstatement, still a win but a different narrative). The root cause is
per-batch variance in EverMemBench: F\_MH sigma \textasciitilde{} 2.3
pp, F\_HL sigma \textasciitilde{} 5 pp, MA\_C/P/U sigma
\textasciitilde{} 3 pp --- any single-batch Δ below
\textasciitilde2sigma is noise floor. Batch 004 specifically was a
+1.40sigma to +1.70sigma upper-tail outlier across Phase G and Phase H
runs; without the 5-batch protocol both would have been overclaimed in
print. Additional single-batch failure mode: MA dimension regressions
were \textbf{invisible} in Phase G batch 004 because batch 004 already
had the lowest MA performance of the five batches (selection bias),
hiding the −3 to −4 pp MA cost entirely.

\textbf{Overstatement rates observed:}

{\def\LTcaptype{none} % do not increment counter
\begin{longtable}[]{@{}llrrr@{}}
\toprule\noalign{}
Phase & Metric & Single-batch & 5-batch & Overstatement \\
\midrule\noalign{}
\endhead
\bottomrule\noalign{}
\endlastfoot
Phase G & F\_MH & +8.00 pp & +1.61 pp & 5× \\
Phase G & F\_TP & +11.67 pp & +2.00 pp & 5.8× \\
Phase H v2 & Overall & +11.60 pp & +9.13 pp & 1.27× \\
Lab Q1 \#4 & F\_MH & +6.78 pp (batch 004) & +2.81 pp & 2.4× \\
\end{longtable}
}

\paragraph{5.8.3 MA dimension is mandatory in every eval
report}\label{ma-dimension-is-mandatory-in-every-eval-report}

Memory Awareness (MA\_C, MA\_P, MA\_U) is a \textbf{silent killer
dimension}: Phase G batch 004 gate completely missed MA regression
because MA was not measured in the initial single-batch run. Any
retrieval change that involves reranking, query routing, or context
modification must audit all three MA sub-dimensions, not just F\_* and
overall accuracy. MA regression indicates the change is damaging the
system's ability to maintain user profile knowledge --- the core
differentiator of nox-mem vs retrieval-only systems.

\paragraph{5.8.4 Search error rate
monitoring}\label{search-error-rate-monitoring}

Concurrent agent operations during Lab Q1 benchmarking caused a batch
contamination incident (PR \#379, batch 010): a concurrent agent
re-installed its adapter mid-run, contaminating results. Recovery via
merged adapter pattern. Lesson: shared adapter install paths on VPS are
a race condition; sequential dispatch and 0/n search-error-per-batch
monitoring are mandatory before accepting 5-batch results.

\paragraph{5.8.5 Honest scope of EverMemBench, LoCoMo, and classical-QA
claims}\label{honest-scope-of-evermembench-locomo-and-classical-qa-claims}

\begin{itemize}
\tightlist
\item
  \textbf{EverMemBench Phase D headline (+2.95 pp vs MemOS Gemini)} is a
  modest win; the structural differentiator is the Memory Awareness
  composite, consistently strong across all backbones.
\item
  \textbf{EverMemBench Phase H v2 headline (+9.13 pp vs MemOS
  GPT-4.1-mini)} is real and CI-verified, but the absolute score
  (51.68\%) is not high --- MemOS itself is only 42.55\%. GPT-4.1-mini
  is the only tested backbone where all memory systems gain vs Full
  Context.
\item
  \textbf{EverMemBench Backbone Matrix (Gemini-3-flash): +20.73 pp
  Overall / +32.74 pp MA composite} is the strongest cross-system claim
  in the paper. The lift is the multiplicative interaction of nox-mem's
  V10 retrieval stack and frontier-tier reasoning, not exclusively
  backbone-driven (§5.1.10).
\item
  \textbf{MuSiQue F1 58.62\%} (§5.2.1) and \textbf{HotPotQA ans\_F1
  73.37\%} (§5.2.2) are both above published reader SOTA without
  specialized fine-tuning, validating the architecture's multi-hop
  reasoning quality on classical benchmarks. Comparison ranges are from
  the public IRCoT, EX(SA), and DPR+FiD literature.
\item
  \textbf{LoCoMo retrieval@10 strict 74.52\%} (§5.3) above Mem0 SOTA F1
  66.88\% is the retrieval ceiling on LoCoMo; the F1 push of 51.85\% is
  rank-5 (above Zep / LangMem, below Mem0 SOTA 66.88\%) due to a
  verbosity gap (§5.3.2). Path to \textgreater=55\% requires
  orchestration (§5.5).
\item
  \textbf{KG path, MAP, MQ, adaptive classifier, and Q3 IterC} are
  opt-in features, not defaults. Each addresses a known structural gap;
  combined effects measured in Wave B/C (§5.1.9).
\item
  The sanitize fix (\texttt{unicode-aware-sanitize-for-fts5}) is a
  prerequisite for all scores reported here; pre-fix Q2 numbers (nDCG@10
  0.9126 LongMemEval) would have been reported as lower and should not
  be compared directly.
\item
  \textbf{Wave 2 NO-REPLICATE findings (D75, §5.5.5):} the Lab Q1
  single-stage knob lifts in §5.1.8 were measured on gpt-4.1-mini and
  are valid for that backbone. They do NOT transfer reliably to
  Gemini-3-flash (transfer rate \textasciitilde0--40\%). Any claim of
  ``Wave A knob X delivers +N pp F\_MH'' must specify the backbone. The
  §5.5.4 composability matrix replaces the D74 projection with measured
  numbers; the original projection table is superseded and should not be
  cited.
\item
  \textbf{IterB composability projection from D74 (IterB + Wave C →
  \textasciitilde12.07\% F\_MH) is superseded.} The projection assumed
  both backbone-portability (refuted by D75) and architectural
  composability (refuted by adapter guard discovery in §5.5.7). The
  corrected empirical upper bound for measured+plausible IterB
  composability on Gemini-3-flash is \textasciitilde8--9\% F\_MH (see
  §5.5.4 corrected table).
\item
  \textbf{D76 capstone (§5.5.8):} the IterB + Wave C triple
  composability outcome is INDETERMINATE due to infrastructure abort.
  This is not a negative scientific result --- it is an untested
  hypothesis. Batch 004 (n=49) is preserved but not 5-batch valid.
\item
  \textbf{Limitations to flag (§5.8.6):} GPT-5 / Claude backbone columns
  are blocked by API access; the Zep \textless100 ms p50 claim is
  unverified by independent runs; the EverMemBench F\_MH absolute number
  (3--7\%) is not directly comparable to multi-hop reasoning gains on
  MuSiQue/HotPotQA --- see §5.4 for the mechanism distinction.
\end{itemize}

\paragraph{5.8.6 Honest limitations and open
work}\label{honest-limitations-and-open-work}

\begin{itemize}
\tightlist
\item
  \textbf{EverMemBench F\_MH absolute (3--7\%) gap vs MemOS Table 4
  (18.88\%)} remains, but the §5.4 reframing demonstrates this is a
  corpus-structural property, not a multi-hop reasoning ceiling. Closing
  it requires either Q3 IterB ReAct (multi-round refinement on long
  conversation chains, §5.5.2) or backbone upgrade (Backbone Matrix
  §5.1.10 shows the gap narrows with Gemini-3-flash). Retrieval-stage
  knobs cap at \textasciitilde+7.25 pp F\_MH (D69 Wave C ceiling §5.1.9)
  and show low backbone-portability to Gemini-3-flash (D75 §5.5.5).
\item
  \textbf{IterB composability with Wave A/B/C knobs on Gemini-3-flash}
  is an open question. The D76 capstone (§5.5.8) was
  infrastructure-aborted before producing valid 5-batch data. The
  composability matrix in §5.5.4 documents this gap honestly. Completing
  the capstone requires dedicated CPU infrastructure.
\item
  \textbf{LoCoMo F1 vs Mem0 SOTA 66.88\%} remains open at rank-5
  (51.85\%). Wave C ceiling analysis (§5.1.9) indicates retrieval-stage
  knobs cannot close this gap; composition orchestration (Q3, §5.5) is
  the open path.
\item
  \textbf{Zep \textless100 ms p50 claim} is published in marketing but
  not independently verified. nox-mem KG path p50 = 2.9 ms (§5.7) is
  measured on production VPS with the harness instrumented end-to-end.
  Comparison is fair only when both are measured under matched
  conditions.
\item
  \textbf{GPT-5 / Claude columns} are blocked by API key constraints in
  the current eval setup. Backbone Matrix is currently three-cell
  (Gemini-2.5-flash, GPT-4.1-mini, Gemini-3-flash); GPT-5 and Claude
  entries are in the runway for Q3+ if access opens.
\item
  \textbf{Wave A knob backbone-portability to other backbones} beyond
  Gemini-3-flash is unverified. The D75 \textasciitilde30--40\% transfer
  rate pattern is based on three knobs on one backbone pair. Additional
  backbone pairs (Claude Sonnet 4.6, GPT-5, Gemini 4) require
  independent re-baseline before composability projections can be made.
\item
  \textbf{EverMind-AI / EverMemBench reference baselines} rely on MemOS
  Table 4 published numbers (arxiv:2602.01313); we have not re-run MemOS
  internally on the canonical 5-batch subset, only validated that the
  5-batch sampling preserves the per-category distribution of the
  published numbers.
\end{itemize}

\begin{center}\rule{0.5\linewidth}{0.5pt}\end{center}

\subsection{6. Q4 COMPARISON --- Cross-System Benchmarking
(Pre-registered)}\label{q4-comparison-cross-system-benchmarking-pre-registered}

\begin{quote}
\textbf{Status (updated 2026-06-29 --- canonical + controlled-embedding
done):} The canonical cross-system run executed on a dedicated pod
2026-06-15 (n=100/dataset, k=10, same-namespace fair), resolving the
2026-05-24 infrastructure abort (§7.1 L5). \textbf{3/6 systems with real
data} (nox-mem, Mem0, agentmemory) + \textbf{3 documented gaps} (Zep =
Docker impossible on unprivileged-pod kernel; Letta = agent-OS
\textasciitilde16 min/query; EverMind-AI = third-party keys +
external-repo auth --- all §6.3.1). \textbf{As-configured result: split}
--- nox-mem wins LongMemEval (nDCG@10 0.5234 vs Mem0 0.4764), Mem0 wins
LoCoMo (0.4686 vs nox-mem 0.4263); agentmemory distant third (0.2803 /
0.1587). \textbf{Controlled-embedding variant (rc4, §6.3.2,
2026-06-29):} with both systems on Gemini 3072d over the full n=2,482
set, the split \textbf{inverts} --- nox-mem outperforms Mem0 on both
datasets (LongMemEval 0.5255 vs 0.4061; LoCoMo 0.4952 vs 0.4407) and all
five represented categories (§6.4); three residual confounds (Mem0
version drift, vector backend, sample scope) declared, and the task-type
asymmetry ablated away (generic-embedding nox-mem still wins, −0.34 pp).
Quality tables §6.3; gaps §6.3.1; operational-cost axis §6.8;
reproducibility-as-evidence §6.9. Principles (§6.5), anti-cherry-pick
(§6.6), and pre-registration (§6.7) immutable. Ref:
\texttt{project\_head\_to\_head\_nox\_mem0\_split\_2026\_06\_14},
\texttt{project\_rc4\_controlled\_embedding\_inverts\_split}.
\end{quote}

\subsubsection{6.1 Methodology summary}\label{methodology-summary}

§6 covers the cross-system comparison between nox-mem and five competing
persistent memory systems for AI agents. The complete execution plan is
documented in \texttt{specs/2026-05-23-Q4-comparison-execution-plan.md}
(pre-registered 2026-05-23, prior to the Sat 2026-05-24 run). The
comparison principles (§6.5), anti-cherry-pick guarantees (§6.6), and
formal pre-registration (§6.7) are locked in this section before
execution; only the tables in §6.2/§6.3/§6.4 receive numbers after the
run. The objective is to satisfy the D43 approval gate
(\texttt{docs/DECISIONS.md}) --- nox-mem in top-3 on \textgreater=2 of
the 4 key metrics (nDCG@10, R@10, MRR, latency) --- unlocking GTM Phase
2.

\subsubsection{6.2 Competitors}\label{competitors}

The five competitors were selected by prioritizing GitHub stars, recent
commit activity, and functional overlap with the nox-mem scope. Versions
are locked prior to execution for reproducibility.

{\def\LTcaptype{none} % do not increment counter
\begin{longtable}[]{@{}
  >{\raggedright\arraybackslash}p{(\linewidth - 8\tabcolsep) * \real{0.2000}}
  >{\raggedright\arraybackslash}p{(\linewidth - 8\tabcolsep) * \real{0.2000}}
  >{\raggedright\arraybackslash}p{(\linewidth - 8\tabcolsep) * \real{0.2000}}
  >{\raggedright\arraybackslash}p{(\linewidth - 8\tabcolsep) * \real{0.2000}}
  >{\raggedright\arraybackslash}p{(\linewidth - 8\tabcolsep) * \real{0.2000}}@{}}
\toprule\noalign{}
\begin{minipage}[b]{\linewidth}\raggedright
System
\end{minipage} & \begin{minipage}[b]{\linewidth}\raggedright
Repo
\end{minipage} & \begin{minipage}[b]{\linewidth}\raggedright
Install path
\end{minipage} & \begin{minipage}[b]{\linewidth}\raggedright
Version pinned
\end{minipage} & \begin{minipage}[b]{\linewidth}\raggedright
Default config
\end{minipage} \\
\midrule\noalign{}
\endhead
\bottomrule\noalign{}
\endlastfoot
Mem0 & \texttt{mem0ai/mem0} & \texttt{pip\ install\ mem0ai==0.1.114} &
\texttt{0.1.114} & OpenAI text-embedding-3-small 1536d + faiss (Chroma
swapped to avoid telemetry thread-leak --- §6.3.1) \\
Zep & \texttt{getzep/zep} & Docker compose (zep + postgres) &
\texttt{{[}GAP\ —\ Docker\ unavailable\ on\ pod,\ §6.3.1{]}} & Local
self-host mode \\
Letta (ex-MemGPT) & \texttt{letta-ai/letta} &
\texttt{pip\ install\ letta} (+ \texttt{click\textless{}8.2}) &
\texttt{0.6.6} & SQLite backend; agent-OS retrieval \\
agentmemory & \texttt{rohitg00/agentmemory} & daemon (REST :3111) &
\texttt{as-of\ 2026-06-15} & union store + namespace filter \\
EverMind-AI & EverOS published bench & repo clone
(\texttt{pip\ install\ -e}) & \texttt{{[}GAP\ —\ keys/auth,\ §6.3.1{]}}
& HTTP service (OpenRouter + DeepInfra) \\
\end{longtable}
}

Each system runs with its publicly documented default configuration
(principle 3 of §6.5): no competitor is adversarially tuned to
underperform.

\subsubsection{6.3 Per-system per-dataset
results}\label{per-system-per-dataset-results}

Canonical cross-system × cross-dataset table. K cutoff fixed at 10
across all systems; latency measured externally (wall clock around
adapter call); cost derived from per-system logs (API calls × published
pricing).

\textbf{Preliminary smoke (2026-05-24) --- superseded.} An initial
20-query smoke on an eval-isolated DB validated the harness end-to-end
(nox-mem, full corpus: nDCG@10=0.6380, gold-hit 13/20; Mem0 at a
500-chunk cost-control cap, not directly comparable). Its only role was
pipeline validation; it is \textbf{superseded} by the canonical run
(below) and rc4 (§6.3.2), and the capped-corpus ``concentration vs
coverage'' reading no longer bears on the reported results. Full smoke
figures are archived in \texttt{q4-partial-cross-system-sat-2026-05-24}.
The per-dataset breakdown below is from the \textbf{canonical run
(2026-06-15)}; rc4 (§6.3.2) extends it to the full n=2,482 set.

\textbf{Canonical run --- 2026-06-15 (dedicated RunPod pod,
n=100/dataset, same-namespace fair).} After the 2026-05-24
infrastructure abort (§7.1 L5), the canonical cross-system run was
executed on a dedicated pod on 2026-06-15. \textbf{Two of the five
competitors produced head-to-head quality numbers} (Mem0, agentmemory;
nox-mem is the system under test, not a competitor); \textbf{the other
three are documented deployment non-runs} with explicit reasons (Zep,
Letta, EverMind-AI --- §6.3.1). Protocol: n=100 queries per dataset,
k=10; retrieval quality (nDCG@10 / recall@10 / MRR) under each system's
native default embedding provider; \textbf{same-namespace fair re-query}
--- the union store is queried at k=30 and filtered to the active
dataset namespace before re-scoring the top-10, removing the
cross-dataset distractor confound that inflated an early union-store
reading (14.6\% of nox-mem's raw top-10 were off-dataset chunks).
Per-system latency percentiles were not captured uniformly in this run;
nox-mem operational latency is reported independently in §5.7 (KG path
2.9 ms p50, hybrid 653 ms p50, re-measured 2026-06-15).

\textbf{LongMemEval n=100 (canonical):}

{\def\LTcaptype{none} % do not increment counter
\begin{longtable}[]{@{}
  >{\raggedright\arraybackslash}p{(\linewidth - 10\tabcolsep) * \real{0.1364}}
  >{\raggedleft\arraybackslash}p{(\linewidth - 10\tabcolsep) * \real{0.1818}}
  >{\raggedleft\arraybackslash}p{(\linewidth - 10\tabcolsep) * \real{0.1818}}
  >{\raggedleft\arraybackslash}p{(\linewidth - 10\tabcolsep) * \real{0.1818}}
  >{\raggedleft\arraybackslash}p{(\linewidth - 10\tabcolsep) * \real{0.1818}}
  >{\raggedright\arraybackslash}p{(\linewidth - 10\tabcolsep) * \real{0.1364}}@{}}
\toprule\noalign{}
\begin{minipage}[b]{\linewidth}\raggedright
System
\end{minipage} & \begin{minipage}[b]{\linewidth}\raggedleft
nDCG@10
\end{minipage} & \begin{minipage}[b]{\linewidth}\raggedleft
R@10
\end{minipage} & \begin{minipage}[b]{\linewidth}\raggedleft
MRR
\end{minipage} & \begin{minipage}[b]{\linewidth}\raggedleft
Cost/query
\end{minipage} & \begin{minipage}[b]{\linewidth}\raggedright
Coverage / config
\end{minipage} \\
\midrule\noalign{}
\endhead
\bottomrule\noalign{}
\endlastfoot
\textbf{nox-mem} & \textbf{0.5234} & 0.6535 & \textbf{0.5494} &
\textbf{\$0 (local)} & 100\%; hybrid FTS5 + Gemini-3072d + RRF \\
Mem0 & 0.4764 & 0.6372 & 0.5107 & subscription + OpenAI embed & 100\%;
faiss backend, OpenAI text-embedding-3-small 1536d \\
agentmemory & 0.2803 & n/c & n/c & \$0 (local) & daemon REST; weak
retriever \\
Zep &
\texttt{{[}GAP\ —\ see\ §6.3.1:\ requires\ Docker;\ impossible\ on\ unprivileged\ pod\ kernel{]}}
& --- & --- & --- & --- \\
Letta &
\texttt{{[}GAP\ —\ see\ §6.3.1:\ agent-OS\ latency\ \textasciitilde{}16\ min/query;\ impractical\ for\ n=100{]}}
& --- & --- & --- & --- \\
EverMind-AI &
\texttt{{[}GAP\ —\ see\ §6.3.1:\ requires\ OpenRouter\ +\ DeepInfra\ keys\ +\ external-repo\ install\ authorization{]}}
& --- & --- & --- & --- \\
\end{longtable}
}

\textbf{LoCoMo n=100 (canonical):}

{\def\LTcaptype{none} % do not increment counter
\begin{longtable}[]{@{}
  >{\raggedright\arraybackslash}p{(\linewidth - 10\tabcolsep) * \real{0.1364}}
  >{\raggedleft\arraybackslash}p{(\linewidth - 10\tabcolsep) * \real{0.1818}}
  >{\raggedleft\arraybackslash}p{(\linewidth - 10\tabcolsep) * \real{0.1818}}
  >{\raggedleft\arraybackslash}p{(\linewidth - 10\tabcolsep) * \real{0.1818}}
  >{\raggedleft\arraybackslash}p{(\linewidth - 10\tabcolsep) * \real{0.1818}}
  >{\raggedright\arraybackslash}p{(\linewidth - 10\tabcolsep) * \real{0.1364}}@{}}
\toprule\noalign{}
\begin{minipage}[b]{\linewidth}\raggedright
System
\end{minipage} & \begin{minipage}[b]{\linewidth}\raggedleft
nDCG@10
\end{minipage} & \begin{minipage}[b]{\linewidth}\raggedleft
R@10
\end{minipage} & \begin{minipage}[b]{\linewidth}\raggedleft
MRR
\end{minipage} & \begin{minipage}[b]{\linewidth}\raggedleft
Cost/query
\end{minipage} & \begin{minipage}[b]{\linewidth}\raggedright
Coverage / config
\end{minipage} \\
\midrule\noalign{}
\endhead
\bottomrule\noalign{}
\endlastfoot
Mem0 & \textbf{0.4686} & 0.6171 & 0.4656 & subscription + OpenAI embed &
\textasciitilde95\% (thread-leak ingest loss); faiss, 1536d \\
\textbf{nox-mem} & 0.4263 & 0.5504 & 0.4464 & \textbf{\$0 (local)} &
100\%; hybrid FTS5 + Gemini-3072d + RRF \\
agentmemory & 0.1587 & n/c & n/c & \$0 (local) & daemon REST; weak
retriever \\
Zep / Letta / EverMind-AI & \texttt{{[}GAP\ —\ see\ §6.3.1{]}} & --- &
--- & --- & --- \\
\end{longtable}
}

\textbf{Split result --- honest reading (mandatory per §6.6).} Each
top-tier system wins one benchmark: \textbf{nox-mem wins LongMemEval}
(+0.047 nDCG@10, +0.039 MRR), \textbf{Mem0 wins LoCoMo} (+0.042
nDCG@10). agentmemory is a distant third on both, reflecting a weaker
retriever. This is a genuine \textbf{split, not a dominance claim}:
nox-mem is competitive with the market leader on retrieval quality
\emph{while} delivering the operational profile of §6.8 (single SQLite
file, \$0/query KG path, no service stack). The nox-mem LoCoMo figure
rose from 0.4173 (raw union store) to 0.4263 after the namespace fix ---
Mem0 still wins, confirming the LoCoMo result is a real retrieval gap,
\textbf{not} merely a confound artifact. Per §6.6, both datasets are
reported side by side; we do not cherry-pick the one nox-mem wins.

\textbf{Caveats (per-system composition, declared per §6.5).} (a)
\textbf{Embedding providers differ by native default} --- nox-mem Gemini
3072d, Mem0 OpenAI 1536d, agentmemory its own default; this is the
``as-configured'' fair-comparison stance (§6.5, principle 5); the
controlled-embedding experiment that equalizes this confound is reported
in §6.3.2 (rc4, all-Gemini) and inverts the LoCoMo result. (b)
\textbf{Store composition differs}: nox-mem and agentmemory were queried
from a union store with namespace filtering; Mem0 used single-namespace
stores (its LoCoMo store ingest leaked threads and lost
\textasciitilde5\% of vectors --- §6.3.1). (c) \textbf{Coverage}:
nox-mem 100\% both datasets; Mem0 100\% LongMemEval /
\textasciitilde95\% LoCoMo; agentmemory full. These are reported, not
hidden, consistent with §6.6.

\paragraph{6.3.1 Documented gaps --- three systems that did not
run}\label{documented-gaps-three-systems-that-did-not-run}

Per §6.6, systems that fail setup receive an explicit reason rather than
silent omission. The 2026-06-15 run hit three:

\begin{itemize}
\tightlist
\item
  \textbf{Zep} --- requires a Docker stack (Zep + Postgres). Docker is
  \textbf{impossible on the unprivileged benchmark pod}:
  \texttt{dockerd} fails on \texttt{iptables\ ...\ Permission\ denied}
  (no \texttt{NET\_ADMIN}), and every documented workaround
  (\texttt{-\/-iptables=false\ -\/-bridge=none},
  \texttt{-\/-storage-driver=vfs}) fails on kernel-namespace syscalls
  blocked by the pod seccomp profile. This is a privilege constraint,
  not a configuration error; Zep requires a host with real Docker
  (privileged VM or its hosted Cloud).
\item
  \textbf{Letta (ex-MemGPT)} --- installs and runs natively (CLI
  unblocked via \texttt{click\textless{}8.2}, server starts without
  Docker), but is an \textbf{agent-OS}: retrieval routes through a full
  LLM agent turn (\texttt{archival\_memory\_search} tool call), measured
  at \textbf{\textasciitilde16 min/query} --- n=100 × 2 datasets would
  take days. Ingest also degrades catastrophically (stalls at
  \textasciitilde94\% after \textasciitilde12 h on a sequential
  archival-memory insert). Letta \emph{validates} (it runs) but is
  impractical for a 100-query benchmark in this configuration.
\item
  \textbf{EverMind-AI / EverOS} --- runs only as an HTTP service and
  requires \textbf{two third-party credentials not available to this
  study} (OpenRouter for LLM extraction + DeepInfra for
  embedding/rerank, the latter provider-specific); installation also
  requires authorizing an external-repo \texttt{pip\ install\ -e}. Out
  of scope for the time-box.
\end{itemize}

LightRAG and HippoRAG2 are \textbf{not §6 competitors} --- they are §5
KG/RAG references whose LLM-per-chunk graph-build cost (LightRAG warns
\textgreater6 h on default) places them outside the memory-system
comparison; they are deferred as optional baselines (§7.2).

\paragraph{6.3.2 Embedding-matched variant (rc4 --- all-Gemini, full
eval)}\label{embedding-matched-variant-rc4-all-gemini-full-eval}

The §6.3 split is reported under each system's \emph{native default}
embedder (the ``as-configured'' stance, §6.5, principle 5). Because
nox-mem defaults to Gemini 3072d and Mem0 to OpenAI 1536d, the most
obvious confound on the split is the embedding provider itself. To probe
it we ran an \textbf{embedding-matched} variant --- \textbf{rc4} --- in
which both systems use the \textbf{same embedder}:
\texttt{gemini-embedding-001} at 3072d (the model and dimensionality
nox-mem runs in production), with Mem0 re-pointed to it via an external
config patch (\texttt{lib/all\_gemini\_config}). This was \textbf{not} a
zero-edit swap: the Mem0 adapter also required a one-time
API-compatibility fix for the mem0 2.x \texttt{search}/\texttt{get\_all}
signatures (declared as confound (a) below). The variant also replaces
the n=100 sample with the \textbf{full evaluation set (n=2,482: 1,982
LoCoMo + 500 LongMemEval)} over the full 6,822-chunk corpus, scored
identically to §6.3 (binary-relevance nDCG@10, k=10, same gold
encoding). The evaluation gold is fully covered --- every gold chunk
exists in the corpus on both datasets --- but note this is
\emph{corpus/gold} coverage, \textbf{not} Mem0 \emph{store} coverage
(Mem0 dropped 4 of 6,822 chunks on ingest; confounds below). This is an
embedding-\textbf{matched} comparison (same model + dimensionality),
\textbf{not} an all-else-equal controlled experiment: three residual
confounds are declared below, and a fourth --- the embedding task-type
asymmetry --- is tested by a dedicated ablation and shown not to drive
the result.

{\def\LTcaptype{none} % do not increment counter
\begin{longtable}[]{@{}
  >{\raggedright\arraybackslash}p{(\linewidth - 6\tabcolsep) * \real{0.2000}}
  >{\raggedleft\arraybackslash}p{(\linewidth - 6\tabcolsep) * \real{0.2667}}
  >{\raggedleft\arraybackslash}p{(\linewidth - 6\tabcolsep) * \real{0.2667}}
  >{\raggedleft\arraybackslash}p{(\linewidth - 6\tabcolsep) * \real{0.2667}}@{}}
\toprule\noalign{}
\begin{minipage}[b]{\linewidth}\raggedright
System (all-Gemini 3072d)
\end{minipage} & \begin{minipage}[b]{\linewidth}\raggedleft
LongMemEval nDCG@10 (n=500)
\end{minipage} & \begin{minipage}[b]{\linewidth}\raggedleft
LoCoMo nDCG@10 (n=1,982)
\end{minipage} & \begin{minipage}[b]{\linewidth}\raggedleft
Overall (n=2,482)
\end{minipage} \\
\midrule\noalign{}
\endhead
\bottomrule\noalign{}
\endlastfoot
\textbf{nox-mem} (hybrid FTS5 + Gemini + RRF) & \textbf{0.5255 ± 0.033}
& \textbf{0.4952 ± 0.017} & \textbf{0.5013 ± 0.015} \\
Mem0 (Gemini embedder, Chroma) & 0.4061 ± 0.036 & 0.4407 ± 0.017 &
0.4337 ± 0.016 \\
\end{longtable}
}

Intervals are 95\% confidence (mean per-query nDCG@10 ± 1.96·SEM
(standard error of the mean), binary relevance). The nox-mem / Mem0
intervals are \textbf{disjoint on both datasets and overall}, so the gap
is not attributable to query-sampling noise (it remains subject to the
three residual configuration confounds below; the one asymmetry that
favored nox-mem --- task type --- was ablated and shown not to drive the
result).

\textbf{Result --- the split does not survive embedding matching.} With
the embedding model and dimensionality equalized (and the three residual
confounds below declared, the fourth ablated away), \textbf{nox-mem
outperforms Mem0 on both benchmarks} (LongMemEval +0.119, LoCoMo +0.055
nDCG@10) and \textbf{all five represented query categories} (§6.4). The
Mem0 LoCoMo win in §6.3 was therefore substantially an embedder effect
--- OpenAI 1536d on LoCoMo --- rather than a retrieval-architecture
advantage: once both systems embed with the same Gemini model, nox-mem's
hybrid FTS5 + dense + RRF retrieval wins on LoCoMo as well.

\textbf{Residual confounds --- declared; this is embedding-matching, not
a clean architecture isolation (per §6.6).} Three factors remain that
prevent attributing the inversion purely to the embedding (a fourth ---
task-type asymmetry --- was tested by ablation and neutralized; see
below). \textbf{(a) Mem0 version drift:} the Mem0 client changed major
version between the canonical run and rc4 (0.1.x → 2.0.10; its
\texttt{search}/\texttt{get\_all} API changed, requiring a compatibility
fix to our adapter), so rc4's Mem0 (0.4337) is \emph{not} the same build
as §6.3's Mem0 (0.4686 LoCoMo) --- part of the gap may be the version.
\textbf{(b) Vector backend:} canonical Mem0 used a faiss store; rc4 Mem0
uses a fresh Chroma collection (required to avoid a dimension clash with
the 1536d OpenAI run) --- a backend change, not only an embedder change.
\textbf{(c) Sample scope:} §6.3 sampled n=100/dataset; rc4 scores the
full n=2,482 --- a larger, differently-distributed query set, not a
like-for-like re-score of the same 100. rc4 is therefore best read as
\emph{``same embedding model and dimensionality, Mem0 at its current
release and default backend''} → nox-mem leads --- \textbf{not} a
surgical architecture-only isolation. (The 4 chunks Mem0 dropped to
transient 503s affect one query, whose gold chunk Mem0 fails to retrieve
in top-10 even when present --- verified zero aggregate effect.) Raw
results and aggregator output: \texttt{eval/q4-comparison/output/rc4/}.

\textbf{Task-type ablation --- confound (d) tested and neutralized.} The
one configuration asymmetry that \emph{favored} nox-mem in rc4 was the
embedding task type: nox-mem passes Gemini's
\texttt{RETRIEVAL\_DOCUMENT}/\texttt{RETRIEVAL\_QUERY} task types
(retrieval-optimized embeddings), while Mem0's embedder call does not.
To isolate it we re-ran nox-mem with a \textbf{generic Gemini embedding}
(\texttt{NOX\_EMBED\_GENERIC\_TASKTYPE=1} --- no task type, exactly how
Mem0 calls the same model) against the \textbf{same} Mem0 baseline, over
the full n=2,482 --- a symmetric \emph{``neither system sets a task
type''} comparison holding confounds (a)--(c) constant. nox-mem's
overall nDCG@10 falls only \textbf{0.5013 → 0.4979 (−0.34 pp)} and
\textbf{still outperforms Mem0 (0.4337) on overall, both datasets
(LoCoMo 0.4920 vs 0.4407; LongMemEval 0.5215 vs 0.4061), and all five
categories}. The task-type asymmetry therefore contributes at most 0.34
pp and does \textbf{not} explain the §6.3 → §6.3.2 inversion: the win is
architectural (hybrid FTS5 + dense + RRF fusion), not an embedding-mode
artifact. (Rigor caveat: the re-ingested generic corpus reached 99.03\%
gold coverage --- 23 of 2,370 distinct gold chunks absent vs 100\% in
the task-type run, a transient-ingest handicap that can only
\emph{lower} nox-mem's score; the win persists despite it.) Ablation
output: \texttt{eval/q4-comparison/output/rc4-ablation/}.

\textbf{LoCoMo claim taxonomy (orientation).} ``Who wins LoCoMo''
depends entirely on metric and configuration; the paper reports four
distinct, non-contradictory readings, tabulated here to prevent
conflation:

{\def\LTcaptype{none} % do not increment counter
\begin{longtable}[]{@{}
  >{\raggedright\arraybackslash}p{(\linewidth - 8\tabcolsep) * \real{0.1875}}
  >{\raggedright\arraybackslash}p{(\linewidth - 8\tabcolsep) * \real{0.1875}}
  >{\raggedright\arraybackslash}p{(\linewidth - 8\tabcolsep) * \real{0.1875}}
  >{\raggedleft\arraybackslash}p{(\linewidth - 8\tabcolsep) * \real{0.2500}}
  >{\raggedright\arraybackslash}p{(\linewidth - 8\tabcolsep) * \real{0.1875}}@{}}
\toprule\noalign{}
\begin{minipage}[b]{\linewidth}\raggedright
§
\end{minipage} & \begin{minipage}[b]{\linewidth}\raggedright
Metric
\end{minipage} & \begin{minipage}[b]{\linewidth}\raggedright
Configuration
\end{minipage} & \begin{minipage}[b]{\linewidth}\raggedleft
n
\end{minipage} & \begin{minipage}[b]{\linewidth}\raggedright
Result
\end{minipage} \\
\midrule\noalign{}
\endhead
\bottomrule\noalign{}
\endlastfoot
§5.3.1 & retrieval@10 (strict) & nox-mem standalone & --- & nox-mem
74.52\% (vs Mem0's \emph{published} F1 66.88\% --- different metric, not
head-to-head) \\
§5.3.2 & answer F1 & nox-mem standalone & --- & nox-mem rank-5 (below
Mem0; verbosity-constrained, §5.3.2) \\
§6.3 & nDCG@10 & native embedders (nox Gemini-3072d / Mem0 OpenAI-1536d)
& 100 & \textbf{Mem0} 0.4686 vs nox-mem 0.4263 \\
§6.3.2 & nDCG@10 & matched embedder (both Gemini-3072d) & 2,482 &
\textbf{nox-mem} 0.4952 vs Mem0 0.4407 \\
\end{longtable}
}

The §6.3 → §6.3.2 reversal is the embedding-matching effect (with the
three residual confounds above; the task-type asymmetry was ablated and
ruled out); the §5.3 retrieval-vs-F1 contrast is an orthogonal metric
distinction, not a contradiction.

\subsubsection{6.4 Per-category breakdown}\label{per-category-breakdown}

Per-query-category breakdown from the \textbf{rc4 embedding-matched run}
(§6.3.2; all-Gemini 3072d, n=2,482 across both datasets). The canonical
2026-06-15 run reported only dataset-level metrics, so the breakdown is
populated from rc4, where per-category gold labels --- the datasets'
\emph{native} question types, mapped to six canonical buckets --- are
wired into the harness. Only nox-mem and Mem0 participate: agentmemory's
server-side embedder cannot be re-pointed to Gemini, and the three
§6.3.1 gaps likewise do not run under the matched embedder.

{\def\LTcaptype{none} % do not increment counter
\begin{longtable}[]{@{}lrrrr@{}}
\toprule\noalign{}
Category & n & nox-mem nDCG@10 & Mem0 nDCG@10 & Δ \\
\midrule\noalign{}
\endhead
\bottomrule\noalign{}
\endlastfoot
single-hop & 438 & \textbf{0.3969} & 0.3908 & +0.006 \\
multi-hop & 454 & \textbf{0.5481} & 0.4699 & +0.078 \\
temporal & 225 & \textbf{0.3940} & 0.2760 & +0.118 \\
adversarial & 524 & \textbf{0.4370} & 0.2955 & +0.142 \\
open-domain & 841 & \textbf{0.5991} & 0.5649 & +0.034 \\
numeric & n/a & n/a & n/a & --- \\
\end{longtable}
}

nox-mem leads every represented category. The largest margins are on
\textbf{adversarial} (+0.142) and \textbf{temporal} (+0.118) --- the
query types where naive semantic similarity (Mem0's Gemini-only path) is
weakest and the FTS5 lexical channel in nox-mem's hybrid fusion
contributes most; the narrowest is \textbf{single-hop} (+0.006), where a
single strong dense match suffices and the hybrid adds little. This
pattern is consistent with the §5 ablations attributing nox-mem's edge
to multi-signal fusion rather than embedding quality alone, and survives
the task-type ablation (§6.3.2): with a generic Gemini embedding,
nox-mem still leads all five categories (adversarial 0.4573 vs 0.2955;
multi-hop 0.5561 vs 0.4699; open-domain 0.5792 vs 0.5649; single-hop
0.3924 vs 0.3908; temporal 0.3765 vs 0.2760).

\textbf{Bucket composition (declared per §6.6).} The six buckets
aggregate the two datasets' native labels, and two mappings warrant
disclosure because they affect interpretation. The \textbf{adversarial}
bucket (n=524) combines LoCoMo's 446 native adversarial queries with 78
LongMemEval \texttt{knowledge-update} queries --- a mapping flagged as
\emph{ambiguous} in the labeling audit
(\texttt{docs/rc2-per-category-mapping.md}), since
\texttt{knowledge-update} is not adversarial in LoCoMo's sense; the
adversarial cell (nox-mem's largest margin, +0.142) should be read as
``adversarial + knowledge-update'', not pure adversarial. The
\texttt{numeric} category receives \texttt{n/a}: neither dataset
contributes at least 10 numeric-typed queries (open-domain is likewise
LongMemEval-only), so per the n\textless10 rule those cells are not
scored rather than extrapolated --- avoiding the type of regression seen
in §5.1.4 (temporal \texttt{n/a} in G10b due to a degenerate corpus
gap). The category labeler is
\texttt{eval/q4-comparison/lib/category\_labeler.py}.

\subsubsection{6.5 Fair-comparison
principles}\label{fair-comparison-principles}

The comparison adheres to principles standardized by the published
benchmarking literature (EverMemBench; BEIR, Thakur et al.~2021,
arXiv:2104.08663; MTEB, Muennighoff et al.~2023, arXiv:2210.07316):

\begin{enumerate}
\def\labelenumi{\arabic{enumi}.}
\tightlist
\item
  \textbf{Identical corpus.} All systems receive the same
  \texttt{chunks.text} ingested via each system's native API. No system
  receives an ``optimized'' version of the corpus.
\item
  \textbf{Identical eval set.} Same queries, same gold sets, same random
  seed (\texttt{42} for LongMemEval shuffle).
\item
  \textbf{Native defaults per system.} Each competitor runs with its
  publicly documented default configuration. Competitors are not
  adversarially tuned to lose; if the default config is what is
  published, it is what is evaluated.
\item
  \textbf{K cutoff fixed at 10.} Some systems default to 5 or 20; all
  are forced to \texttt{k=10} for comparability.
\item
  \textbf{Native embeddings provider per system.} nox-mem uses Gemini
  3072d; each competitor uses its default provider. The
  \texttt{all-Gemini} controlled variant that equalizes this provider
  --- originally planned as an optional side experiment --- was
  \textbf{executed (rc4)} and is reported in §6.3.2: with both systems
  on Gemini 3072d over the full n=2,482 set, nox-mem outperforms Mem0 on
  both datasets and all five represented categories (whereas
  as-configured, §6.3, Mem0 wins LoCoMo), with three residual confounds
  declared and the task-type asymmetry ablated away (a generic-embedding
  re-run still wins by at least +0.05 nDCG@10 on both datasets).
\item
  \textbf{Uniform hardware.} Same VPS (Hostinger 8 cores / 16 GB RAM),
  localhost between systems except for external embeddings API calls.
\end{enumerate}

Each adapter passes a pre-run smoke test:

\begin{Shaded}
\begin{Highlighting}[]
\NormalTok{result }\OperatorTok{=}\NormalTok{ adapter.search(}\StringTok{"test query"}\NormalTok{, k}\OperatorTok{=}\DecValTok{5}\NormalTok{)}
\ControlFlowTok{assert} \BuiltInTok{len}\NormalTok{(result) }\OperatorTok{\textgreater{}=} \DecValTok{1}
\ControlFlowTok{assert} \BuiltInTok{all}\NormalTok{(}\StringTok{\textquotesingle{}id\textquotesingle{}} \KeywordTok{in}\NormalTok{ r }\KeywordTok{and} \StringTok{\textquotesingle{}score\textquotesingle{}} \KeywordTok{in}\NormalTok{ r }\ControlFlowTok{for}\NormalTok{ r }\KeywordTok{in}\NormalTok{ result)}
\end{Highlighting}
\end{Shaded}

Any adapter that fails the smoke test is documented as a gap
(\texttt{{[}FAILED:\ \textless{}reason\textgreater{}{]}}) rather than
omitted --- consistent with §6.6.

\subsubsection{6.6 Anti-cherry-pick
statement}\label{anti-cherry-pick-statement}

To prevent retroactive selection bias:

\begin{itemize}
\item
  \textbf{All 6 categories reported.} None is omitted because the result
  is unfavorable.
\item
  \textbf{Both datasets reported.} LongMemEval n=100 + LoCoMo full, side
  by side. We do not cherry-pick whichever benefits nox-mem.
\item
  \textbf{Worst-case latency reported.} p50 + p95 + p99 explicit. We do
  not publish only p50.
\item
  \textbf{Per-system per-category transparency.} The §6.4 table exposes
  every combination; there is no aggregate row that masks a pattern.
\item
  \textbf{Gaps documented.} Systems that fail setup receive an explicit
  note; the comparison runs without the missing system, but the gap is
  recorded in \texttt{docs/COMPARISON.md}.
\item
  \textbf{Explicit per-dataset breakdown (PR \#318, 2026-05-23 ---
  rev3).} The Gemini hybrid@500 run revealed that the aggregate (0.0918)
  masks a decisive per-dataset result. We report all three rows
  explicitly:

  {\def\LTcaptype{none} % do not increment counter
  \begin{longtable}[]{@{}
    >{\raggedright\arraybackslash}p{(\linewidth - 8\tabcolsep) * \real{0.1667}}
    >{\raggedleft\arraybackslash}p{(\linewidth - 8\tabcolsep) * \real{0.2222}}
    >{\raggedleft\arraybackslash}p{(\linewidth - 8\tabcolsep) * \real{0.2222}}
    >{\raggedleft\arraybackslash}p{(\linewidth - 8\tabcolsep) * \real{0.2222}}
    >{\raggedright\arraybackslash}p{(\linewidth - 8\tabcolsep) * \real{0.1667}}@{}}
  \toprule\noalign{}
  \begin{minipage}[b]{\linewidth}\raggedright
  System
  \end{minipage} & \begin{minipage}[b]{\linewidth}\raggedleft
  nDCG@10 (aggregate)
  \end{minipage} & \begin{minipage}[b]{\linewidth}\raggedleft
  nDCG@10 (LoCoMo-only)
  \end{minipage} & \begin{minipage}[b]{\linewidth}\raggedleft
  Corpus
  \end{minipage} & \begin{minipage}[b]{\linewidth}\raggedright
  Mode
  \end{minipage} \\
  \midrule\noalign{}
  \endhead
  \bottomrule\noalign{}
  \endlastfoot
  nox-mem FTS5@500 & 0.0466 & --- & 500 (cap) & FTS5-only \\
  \textbf{nox-mem Gemini hybrid@500} & 0.0918 & \textbf{0.1835} & 500
  (cap) & FTS5 + Gemini + RRF \\
  \textbf{mem0@500} & \textbf{0.1315} & 0.1315 & 500 (cap) & LLM rewrite
  + embed \\
  \end{longtable}
  }

  \textbf{LoCoMo-only result (PR \#318):} nox-mem Gemini hybrid@500 =
  0.1835 \textbf{exceeds} mem0@500 = 0.1315 by \textbf{+40\%} on the
  conversational memory dimension. The aggregate (0.0918) falls below
  mem0 due to a \textbf{corpus-ordering artifact}: at the 500-chunk cap,
  the 5,882 LoCoMo chunks exhaust the cap before any LongMemEval chunk
  is ingested --- the 10 LongMemEval queries have zero coverage, zeroing
  the nDCG for that dataset and pulling the aggregate down. Hybrid stack
  lift over FTS5@500: \textbf{+97\%} (0.0466 → 0.0918), validating the
  architectural value of the stack even on a sparse corpus.

  \textbf{H2 finding (PR \#311, retained):} FTS5-only@500 = 0.0466 vs
  mem0@500 = 0.1315 is \textbf{real and architectural} for FTS5-only
  mode --- mem0's LLM-rewriting produces semantic generalization that
  isolated FTS5 cannot match. PR \#318 shows that the full Gemini hybrid
  reverses this result on the conversational scope.

  \textbf{Mandatory disclosure:} the aggregate ±0.05 falls within the
  inconclusive interval for n=20. The definitive arbiter is the
  canonical full-corpus run (uniform corpus without cap, full LoCoMo +
  LongMemEval for all systems). \textbf{Phase 2 gate uses BOTH}
  per-dataset + aggregate from the canonical run --- not only the number
  that favors nox-mem.

  Refs:
  \texttt{docs/COMPARISON.md\ §Apples-to-apples\ corpus-cap\ comparison},
  PR \#311, PR \#318.
\end{itemize}

\subsubsection{6.7 Pre-registration}\label{pre-registration}

This section's methodology is locked in
\texttt{specs/2026-05-23-Q4-comparison-execution-plan.md} prior to the
Sat 2026-05-24 run. The \textbf{Sat 2026-05-24 15:30 BRT smoke}
populated the first row of §6.3 (nox-mem combined: nDCG@10=0.6380, p50=8
ms, gold-hit 13/20 on 20 dry-run-sample queries) and validated that the
retrieval pipeline works end-to-end on an eval-isolated DB. The
\textbf{partial cross-system smoke of Sat 2026-05-24 18:00 BRT} added
the mem0 row (n=20, 500-chunk corpus cap): nDCG@10=0.8569, p50=273 ms,
gold-hit 3/20 (15\%) --- with an explicit interpretation of the coverage
vs.~concentration trade-off in §6.3. The \textbf{canonical run was
executed on 2026-06-15} on a dedicated pod (n=100/dataset, k=10,
same-namespace fair), resolving the 2026-05-24 infrastructure abort
(incident D76 --- sustained CPU-steal on the shared host, see §7.1 L5).
It produced real numbers for 3/6 systems (nox-mem, Mem0, agentmemory)
and three documented gaps (§6.3.1), updating the §6.3 cells. The
pre-registered \texttt{all-Gemini} controlled variant (§6.5, principle
5) was subsequently executed as \textbf{rc4} (§6.3.2) --- equalizing the
embedding provider over the full n=2,482 set --- and supplied the §6.4
per-category breakdown; this was a planned confound-control experiment,
not a post-hoc methodology change, and its three residual confounds
(Mem0 version drift, vector backend, sample scope) are declared in
§6.3.2 per §6.6, with the task-type asymmetry tested by a dedicated
ablation and shown not to drive the result. No methodology from
§6.5--§6.6 was altered post-registration; the same-namespace re-query is
a fair-comparison refinement (confound removal, §6.3) consistent with
§6.5, not a change to corpus, eval set, or gold. Principles (§6.5),
anti-cherry-pick (§6.6), and the general structure of this section are
immutable post-run. Any methodological adjustment identified during
execution is documented as an explicit follow-up in
\texttt{docs/COMPARISON.md} rather than retroactively applied here.
Refs: \texttt{q4-smoke-sat-2026-05-24-real-numbers} ·
\texttt{q4-partial-cross-system-sat-2026-05-24}.

Decision D43 (\texttt{docs/DECISIONS.md}) defines the approval gate:
nox-mem in top-3 on \textgreater=2 of the 4 key metrics (nDCG@10, R@10,
MRR, latency). If the gate is met, GTM Phase 2 is unlocked per
\texttt{docs/ROADMAP.md} §7. If not met, the Sun 2026-05-25 session
produces a remediation plan (pre-launch adjustments) rather than a
direct launch.

\subsubsection{6.8 Autonomy quantified --- operational cost per memory
system}\label{autonomy-quantified-operational-cost-per-memory-system}

The Q4 quality comparison (§6.3 -- §6.6) reports retrieval
\emph{quality} under matched corpora. Operational \emph{cost} ---
services, RAM, cold start, mandatory third-party credentials, setup
commands --- is the second axis on which a memory system can be
evaluated, and is the axis where the nox-mem Autonomy pillar \footnote{Q/A/P
  strategic pivot of 2026-05-17 --- three pillars (\textbf{Q}uality,
  \textbf{A}utonomy, \textbf{P}roduct). See \texttt{docs/ROADMAP.md} and
  \texttt{qap-pillars-strategic-decision} in the project memory.} is
most legible. Table 2 summarizes the steady-state idle footprint of each
system in its default self-host configuration.

\textbf{Table 2 --- Autonomy quantified: services, RAM, cold start,
mandatory keys, setup commands.} Headline: \textbf{\textasciitilde12×
less RSS than EverOS, single process, no service stack.} Competitor
numbers are \emph{estimates} derived from each project's docker-compose
defaults and documented system requirements (sources cited in the row).
The nox-mem row is \textbf{{[}measured 2026-05-24, prod VPS, 6830 chunks
live, uptime 9h28min{]}} via \texttt{ps\ -eo\ pid,rss,vsz,comm,args}
against the production \texttt{nox-mem-api} process (PID 2422197). See
footnote \footnote{The \texttt{\textasciitilde{}341\ MB\ RSS} figure for
  nox-mem in Table 2 is \textbf{measured 2026-05-24} on the production
  VPS (PID 2422197, uptime 9h28min, 6830 chunks live, 100\% vector
  coverage). Main process RSS = 349,276 KB via
  \texttt{ps\ -eo\ pid,rss,vsz,comm,args\ \textbar{}\ grep\ dist/api-server.js}.
  The cgroup \texttt{MemoryCurrent} reported by
  \texttt{systemctl\ show\ nox-mem-api\ -p\ MemoryCurrent} is
  727,064,576 bytes (\textasciitilde727 MB); the \textasciitilde386 MB
  delta vs process RSS is SQLite memory-mapped I/O (chunks table, FTS5
  index, vec0 index) --- kernel-managed page cache, reclaimable on
  memory pressure, not exclusive process memory. Standard
  \texttt{ps}/\texttt{top} RSS is the canonical comparison metric used
  in Table 2 across all competitors. Original revision marked this as
  \texttt{\textasciitilde{}50\ MB\ {[}estimated{]}} based on Node
  baseline + better-sqlite3 cache projections; the production
  measurement (initially denied during the first revision and granted
  post-PR \#356) replaced the estimate.} for full methodology including
the cgroup \texttt{MemoryCurrent} vs process RSS distinction.

{\def\LTcaptype{none} % do not increment counter
\begin{longtable}[]{@{}
  >{\raggedright\arraybackslash}p{(\linewidth - 12\tabcolsep) * \real{0.1154}}
  >{\raggedleft\arraybackslash}p{(\linewidth - 12\tabcolsep) * \real{0.1538}}
  >{\raggedleft\arraybackslash}p{(\linewidth - 12\tabcolsep) * \real{0.1538}}
  >{\raggedleft\arraybackslash}p{(\linewidth - 12\tabcolsep) * \real{0.1538}}
  >{\raggedleft\arraybackslash}p{(\linewidth - 12\tabcolsep) * \real{0.1538}}
  >{\raggedleft\arraybackslash}p{(\linewidth - 12\tabcolsep) * \real{0.1538}}
  >{\raggedright\arraybackslash}p{(\linewidth - 12\tabcolsep) * \real{0.1154}}@{}}
\toprule\noalign{}
\begin{minipage}[b]{\linewidth}\raggedright
System
\end{minipage} & \begin{minipage}[b]{\linewidth}\raggedleft
Services
\end{minipage} & \begin{minipage}[b]{\linewidth}\raggedleft
RAM idle
\end{minipage} & \begin{minipage}[b]{\linewidth}\raggedleft
Cold start
\end{minipage} & \begin{minipage}[b]{\linewidth}\raggedleft
Mandatory third-party keys
\end{minipage} & \begin{minipage}[b]{\linewidth}\raggedleft
Setup commands
\end{minipage} & \begin{minipage}[b]{\linewidth}\raggedright
Sources
\end{minipage} \\
\midrule\noalign{}
\endhead
\bottomrule\noalign{}
\endlastfoot
\textbf{nox-mem} & \textbf{1} (SQLite file + Node process) &
\textbf{\textasciitilde341 MB RSS} {[}measured 2026-05-24{]} &
\textbf{\textless1 s} & \textbf{0} (offline-OK; embeddings optional) &
\textbf{1} (\texttt{npm\ i\ \&\&\ nox-mem\ reindex}) & This work;
\footnote{The \texttt{\textasciitilde{}341\ MB\ RSS} figure for nox-mem
  in Table 2 is \textbf{measured 2026-05-24} on the production VPS (PID
  2422197, uptime 9h28min, 6830 chunks live, 100\% vector coverage).
  Main process RSS = 349,276 KB via
  \texttt{ps\ -eo\ pid,rss,vsz,comm,args\ \textbar{}\ grep\ dist/api-server.js}.
  The cgroup \texttt{MemoryCurrent} reported by
  \texttt{systemctl\ show\ nox-mem-api\ -p\ MemoryCurrent} is
  727,064,576 bytes (\textasciitilde727 MB); the \textasciitilde386 MB
  delta vs process RSS is SQLite memory-mapped I/O (chunks table, FTS5
  index, vec0 index) --- kernel-managed page cache, reclaimable on
  memory pressure, not exclusive process memory. Standard
  \texttt{ps}/\texttt{top} RSS is the canonical comparison metric used
  in Table 2 across all competitors. Original revision marked this as
  \texttt{\textasciitilde{}50\ MB\ {[}estimated{]}} based on Node
  baseline + better-sqlite3 cache projections; the production
  measurement (initially denied during the first revision and granted
  post-PR \#356) replaced the estimate.} \\
mem0 & 2 (Postgres + Qdrant) & \textasciitilde800 MB & \textasciitilde15
s & 1 (OpenAI for embeddings) & \textasciitilde5 & mem0 docker-compose
defaults \footnote{mem0 default self-host requires Postgres + Qdrant +
  an OpenAI key for embeddings (or a configured alternative provider).
  Counts: 2 services + 1 mandatory third-party key. Source: mem0 README
  and \texttt{docker-compose.yml} defaults at github.com/mem0ai/mem0.} \\
Letta & 3 (Letta server + Postgres + OpenAI) & \textasciitilde1.5 GB &
\textasciitilde30 s & 1 (OpenAI) & \textasciitilde8 & Letta self-host
guide \footnote{Letta default self-host requires the Letta server,
  Postgres, and an OpenAI key (or alternative LLM provider) for the
  agent loop. Counts: 3 components + 1 mandatory third-party key.
  Source: Letta self-host documentation at docs.letta.com and
  github.com/letta-ai/letta.} \\
Zep OSS & 2 (Zep + Postgres) & \textasciitilde1.2 GB & \textasciitilde30
s & \textbf{1 mandatory} (OpenAI for embeddings --- hardcoded) &
\textasciitilde6 & Zep README \footnote{Zep OSS requires the Zep service
  container + Postgres, and \emph{hardcodes} OpenAI as the embedding
  provider in the default build (mandatory key, no fallback in OSS
  distribution). Counts: 2 services + 1 hardcoded mandatory key. Source:
  Zep README at github.com/getzep/zep.} \\
EverOS / EverMind-AI & \textbf{5} (MongoDB + Elasticsearch + Milvus +
Redis + Postgres) & \textbf{\textasciitilde4 GB+} &
\textbf{\textasciitilde60 s} & 2--3 (LLM + embedding + optional
reranker) & \textasciitilde15+ & EverMind-AI docker-compose
\footnote{EverMind-AI / EverOS docker-compose declares MongoDB +
  Elasticsearch + Milvus + Redis + Postgres = 5 services, plus 2--3
  third-party API keys for LLM, embedding, and (optional) reranker.
  Counts confirmed against the published \texttt{docker-compose.yml} in
  the EverMind-AI repo. The \textasciitilde4 GB RAM-idle figure is the
  sum of documented minimum requirements for each service's container.} \\
LightRAG & 2 (Neo4j + vector DB) & \textasciitilde1 GB &
\textasciitilde20 s & 1 (LLM provider for KG extraction) &
\textasciitilde6 & LightRAG repo defaults \footnote{LightRAG defaults to
  Neo4j for the knowledge graph + a vector store (Qdrant/Milvus/etc.),
  plus one LLM provider key for entity/relation extraction during
  indexing. Counts: 2 services + 1 mandatory third-party key. Source:
  LightRAG README at github.com/HKUDS/LightRAG.} \\
\end{longtable}
}

\textbf{Reading the table.} Three rows of the cost matrix translate
directly into Autonomy:

\begin{enumerate}
\def\labelenumi{\arabic{enumi}.}
\item
  \textbf{Services column.} Every additional service is an additional
  failure mode, an additional security-patching surface, and an
  additional vendor that must be available on the day a user spins up
  the system. nox-mem ships as a single Node process operating on a
  single SQLite file; the only durable on-disk artifact is
  \texttt{nox-mem.db}. mem0/Zep/Letta/LightRAG each require
  \textgreater=1 database container and at least one external
  LLM/embedding provider. EverOS requires five containers, three of
  which are heavyweight infrastructure (MongoDB, Elasticsearch, Milvus).
  The single-service property is what makes ``open \texttt{nox-mem.db}
  in \texttt{sqlite3} and inspect everything'' a literal operation, not
  a euphemism.
\item
  \textbf{Mandatory third-party keys column.} A system that requires an
  OpenAI key by default is not autonomous regardless of license --- the
  user is dependent on one specific vendor's pricing, rate limits, and
  terms of service. nox-mem treats embeddings as optional (FTS5-only
  retrieval is a valid degraded mode; §4) and is provider-agnostic when
  embeddings are enabled (Gemini default, Ollama-local feasible --- §7.1
  L2). Zep and Letta hardcode OpenAI as the default embedding provider.
\item
  \textbf{Cold start column.} A \texttt{\textless{}1s} cold start is
  what makes self-host \emph{try-before-deciding} --- the user can
  \texttt{npm\ i}, run one command, see results, and decide. A
  \texttt{\textasciitilde{}60s} cold start with five containers is what
  makes EverOS effectively a ``build a small team to evaluate''
  decision, not an individual decision.
\end{enumerate}

\textbf{Caveat --- RAM measurement methodology.} The competitor RAM
figures are \emph{idle} (i.e., process started, no queries served, no
ingestion in progress) and are \emph{estimates} read from each project's
documented system requirements and \texttt{docker\ stats} defaults in
the published docker-compose files. They are not from a head-to-head
benchmark on a single host. A side-by-side measurement on a controlled
4-vCPU / 8-GB host is a §7.2 future-work item (F-cost-bench). The
nox-mem \texttt{\textasciitilde{}341\ MB\ RSS} figure is
\textbf{measured} on the production VPS via
\texttt{ps\ -eo\ pid,rss,vsz,comm,args} (PID 2422197, uptime 9h28min,
6830 chunks live); see footnote \footnote{The
  \texttt{\textasciitilde{}341\ MB\ RSS} figure for nox-mem in Table 2
  is \textbf{measured 2026-05-24} on the production VPS (PID 2422197,
  uptime 9h28min, 6830 chunks live, 100\% vector coverage). Main process
  RSS = 349,276 KB via
  \texttt{ps\ -eo\ pid,rss,vsz,comm,args\ \textbar{}\ grep\ dist/api-server.js}.
  The cgroup \texttt{MemoryCurrent} reported by
  \texttt{systemctl\ show\ nox-mem-api\ -p\ MemoryCurrent} is
  727,064,576 bytes (\textasciitilde727 MB); the \textasciitilde386 MB
  delta vs process RSS is SQLite memory-mapped I/O (chunks table, FTS5
  index, vec0 index) --- kernel-managed page cache, reclaimable on
  memory pressure, not exclusive process memory. Standard
  \texttt{ps}/\texttt{top} RSS is the canonical comparison metric used
  in Table 2 across all competitors. Original revision marked this as
  \texttt{\textasciitilde{}50\ MB\ {[}estimated{]}} based on Node
  baseline + better-sqlite3 cache projections; the production
  measurement (initially denied during the first revision and granted
  post-PR \#356) replaced the estimate.} for full methodology including
the cgroup \texttt{MemoryCurrent} vs process RSS distinction.

\subsubsection{6.9 Operational reproducibility as
evidence}\label{operational-reproducibility-as-evidence}

The §6.3 quality split and the §6.8 cost matrix report \emph{outcomes}.
A third, harder-to-fake signal emerged from the \textbf{act of running
the benchmark itself}: the operational effort required to get each
system to produce a number is a measurement of its dependency surface.

On the 2026-06-15 pod, nox-mem ingested the full 6,822-chunk corpus into
a single SQLite file in one process with \textbf{zero incidents}. The
competitors did not fare as smoothly:

\begin{itemize}
\tightlist
\item
  \textbf{Mem0} exhausted the pod's process-ID limit three times --- its
  default telemetry (PostHog) leaks one thread per operation, and at
  benchmark scale this wedged the host until ingest and search were
  split into separate processes, telemetry was disabled
  (\texttt{MEM0\_TELEMETRY=False}), and the vector backend was swapped
  from Chroma to faiss. Its LoCoMo store still lost \textasciitilde5\%
  of vectors to the leak before the workaround stabilized.
\item
  \textbf{Zep} never ran --- its Docker stack cannot start on an
  unprivileged pod kernel (§6.3.1).
\item
  \textbf{Letta} ran but at \textasciitilde16 min/query, making a
  100-query benchmark a multi-day proposition.
\item
  \textbf{agentmemory} silently persisted store state across sessions,
  contaminating a LoCoMo run with leftover LongMemEval vectors until the
  store was reset.
\end{itemize}

\textbf{The gaps are not neutral omissions --- they are a deployability
penalty.} Standard benchmarking convention records a non-running system
as a blank cell and moves on. We make a stronger, but carefully bounded,
claim: \emph{the dimension on which these three systems could not
produce a number is precisely the dimension nox-mem is engineered to
win.} Each failure maps to a concrete operational deficit that nox-mem
does not carry:

{\def\LTcaptype{none} % do not increment counter
\begin{longtable}[]{@{}
  >{\raggedright\arraybackslash}p{(\linewidth - 6\tabcolsep) * \real{0.2500}}
  >{\raggedright\arraybackslash}p{(\linewidth - 6\tabcolsep) * \real{0.2500}}
  >{\raggedright\arraybackslash}p{(\linewidth - 6\tabcolsep) * \real{0.2500}}
  >{\raggedright\arraybackslash}p{(\linewidth - 6\tabcolsep) * \real{0.2500}}@{}}
\toprule\noalign{}
\begin{minipage}[b]{\linewidth}\raggedright
System
\end{minipage} & \begin{minipage}[b]{\linewidth}\raggedright
Why it did not run
\end{minipage} & \begin{minipage}[b]{\linewidth}\raggedright
Operational axis it fails
\end{minipage} & \begin{minipage}[b]{\linewidth}\raggedright
nox-mem on the same axis
\end{minipage} \\
\midrule\noalign{}
\endhead
\bottomrule\noalign{}
\endlastfoot
Zep & requires a privileged Docker host; cannot start on an unprivileged
single-node kernel (§6.3.1) & deployability under constraint & runs
unprivileged, single SQLite file, one process \\
Letta & ran, but retrieval routes through a full LLM agent turn at
\textasciitilde16 min/query & query efficiency & 2.9 ms KG path / 653 ms
hybrid --- \textbf{5--6 orders of magnitude faster} \\
EverMind-AI & needs 5 services + 2 mandatory paid third-party keys
(OpenRouter + DeepInfra) & dependency footprint / autonomy & 1 process,
\textbf{0 mandatory keys}, provider-agnostic \\
\end{longtable}
}

\textbf{Honest bound (per §6.6).} We do \emph{not} infer that nox-mem
\emph{retrieves better} than these three systems --- we hold no quality
numbers for them, and on a privileged Docker host with paid keys they
may retrieve competitively. The claim is narrower, and stronger for
being narrow: on the \textbf{deployability / efficiency / autonomy axis}
--- the axis that decides whether an individual developer, or an
embedded agent that must live in its own process on commodity hardware,
can use the system \emph{at all} --- \textbf{three of five competitors
could not be made to produce a single result}, while nox-mem produced
one from one file with zero external services. For that user, a system
you cannot run is not ``untested'': it is unusable, and unusable is the
worst score on the operational axis. A lighter, dependency-free design
is not a footnote to the quality comparison --- it is the reason nox-mem
has a number in every cell of §6.3 and three competitors do not.

The reproducibility of nox-mem's run --- one file, one process, one
command, re-confirmed live at \textbf{415 MB RSS on the 70.7k-chunk
production corpus} on 2026-06-15 --- is the operational expression of
the Autonomy pillar quantified in §6.8. \textbf{A memory system you can
run from a single file you own is a different category of dependency
than one that requires Docker, a vector database, a telemetry pipeline,
and a third-party embedding key to return its first result.}

\begin{center}\rule{0.5\linewidth}{0.5pt}\end{center}

\subsection{7. Limitations and Future
Work}\label{limitations-and-future-work}

\subsubsection{7.1 Limitations}\label{limitations}

\paragraph{L1 --- Explicit-ingestion dependency (no zero-shot corpus
coverage)}\label{l1-explicit-ingestion-dependency-no-zero-shot-corpus-coverage}

nox-mem retrieves only what has been explicitly ingested via
\texttt{ingestFile()}, \texttt{ingest-entity}, or the inotifywait
watcher pipeline. There is no mechanism to answer queries over arbitrary
external corpora at query time. This is a deliberate design constraint
--- the system is optimized for an agent's \emph{own} accumulated
memory, not general-purpose retrieval augmentation --- but it means that
coverage is bounded by ingestion discipline. A corpus that has never
been ingested produces zero recall regardless of query quality. Users
bootstrapping the system must explicitly run \texttt{nox-mem\ reindex}
over existing files before the hybrid search layer is useful. Ref:
\texttt{specs/2026-03-14-nox-memory-system-design.md}, ingestion
pipeline §3.1.

\paragraph{L2 --- Gemini API dependency for embeddings (cost + outbound
network)}\label{l2-gemini-api-dependency-for-embeddings-cost-outbound-network}

The semantic retrieval layer (Layer 2) depends on Google's
\texttt{gemini-embedding-001} model (3072 dimensions). This introduces
two constraints: (a) every vectorization call requires outbound network
access and a valid \texttt{GEMINI\_API\_KEY}, meaning an air-gapped
deployment falls back to FTS5-only retrieval with no semantic recall;
(b) API cost scales with corpus size --- at the Sat 2026-05-24 corpus of
\textasciitilde69k chunks, a full re-vectorization pass takes
approximately 30--40 minutes at quota limits of the free tier. The
Autonomy pillar of the Q/A/P strategy explicitly calls out ``provider
your choice, zero vendor lock-in'' as a long-term goal; local embedding
substitution (e.g., \texttt{nomic-embed-text} via Ollama) is
architecturally feasible but not validated against the canonical eval
set. bring-your-own-key (BYOK) partial autonomy is available: users who
supply their own Gemini API key operate without per-query billing
exposure in the default free tier. Ref: \texttt{docs/DECISIONS.md}
(model selection), \texttt{default-flash-lite-for-agent-infra-tasks}.

\paragraph{L3 --- Single-instance architecture (no distributed sharding
or
replication)}\label{l3-single-instance-architecture-no-distributed-sharding-or-replication}

The system runs on a single SQLite file per agent database. WAL mode
provides concurrent read safety, but there is no horizontal sharding, no
replication across nodes, and no distributed coordination layer. The
current production corpus (\textasciitilde95k chunks across 7 databases
on a 4-vCPU / 8GB KVM4, as of 2026-06-04) operates comfortably within
these bounds, but the architecture does not generalize to multi-tenant
deployments or corpora significantly exceeding the single-node
memory/storage envelope. Distributed SQLite extensions (e.g.,
\texttt{cr-sqlite} CRDT-based replication) exist but are explicitly out
of scope for v1. This is a known architectural decision, not an
oversight. Ref: \texttt{docs/DECISIONS.md} (single-instance rationale).

\paragraph{L4 --- No write-side concurrency control
(last-writer-wins)}\label{l4-no-write-side-concurrency-control-last-writer-wins}

Chunk ingestion operates under an optimistic concurrency model:
\texttt{ingestFile()} deletes existing chunks for the source file and
re-inserts in a single transaction, but there is no row-level locking or
version fence against concurrent ingest of the same source file from two
processes. In practice the inotifywait watcher and manual CLI calls
rarely overlap, and the WAL journal prevents data corruption; however,
two concurrent ingest calls on the same file produce non-deterministic
chunk counts. The \texttt{withOpAudit()} wrapper
(\texttt{src/lib/op-audit.ts}) does not add a mutual-exclusion layer for
ingest --- it targets destructive bulk operations (reindex, consolidate,
crystallize). Production mitigations are operational (systemd service
prevents concurrent watcher processes; cron stagger of 5 minutes between
agents), not architectural. Ref: \texttt{docs/INCIDENTS.md\#2026-04-25},
\texttt{a1-op-audit-module}.

\paragraph{L5 --- Evaluation sample size and canonical run gap (resolved
2026-06-15 /
2026-06-29)}\label{l5-evaluation-sample-size-and-canonical-run-gap-resolved-2026-06-15-2026-06-29}

The Sat 2026-05-24 cross-system run was a 20-query methodology smoke
over an eval-isolated DB (5,882 LoCoMo + 940 LongMemEval chunks), not
the canonical run; the first canonical attempt was \textbf{aborted}
(incident D76: sustained CPU-steal on the shared production host,
51-97\% over 48h, batch 005 0/50 in 23h). The canonical run was
subsequently executed on dedicated infrastructure on \textbf{2026-06-15}
(n=100/dataset, 3/6 systems producing numbers + 3 documented gaps,
§6.3), and the \textbf{controlled-embedding variant (rc4)} was run on
\textbf{2026-06-29} at full scale (n=2,482, both systems on Gemini
3072d, §6.3.2 / §6.4). The §6 competitive figures are therefore
\textbf{settled, not \texttt{{[}deferred{]}}}; the residual limitation
is the three declared confounds on rc4 (Mem0 version drift, vector
backend, sample scope --- §6.3.2), not sample size, and the one
asymmetry that favored nox-mem (embedding task type) was ablated and
ruled out. The earlier nox-mem smoke figure (nDCG@10 = 0.6380 combined)
is not directly comparable to the G5 V3 entity-eval figure (0.6237)
because the eval corpus and query set differ. Ref:
\texttt{specs/2026-05-23-Q4-comparison-execution-plan.md},
\texttt{q4-smoke-sat-2026-05-24-real-numbers}.

\paragraph{L6 --- Cross-system comparison is methodologically
partial}\label{l6-cross-system-comparison-is-methodologically-partial}

Three of five competitors could not be evaluated at all. After the
canonical run (2026-06-15, n=100/dataset) and the rc4 embedding-matched
run (2026-06-29, full n=2,482), Mem0 and agentmemory produced real
numbers alongside nox-mem; the \textbf{500-chunk cap was a property of
the superseded 2026-05-24 smoke only} and no longer applies. The
standing limitation is \textbf{deployability coverage}: Zep (privileged
Docker), Letta (agent-OS \textasciitilde16 min/query), and EverMind-AI
(third-party keys) produced no numbers in the unprivileged single-node
environment (§6.3.1) --- a deployability gap on the operational axis,
not a retrieval-quality claim against them. The rc4 head-to-head
additionally carries the three residual confounds declared in §6.3.2
(its fourth potential confound, embedding task-type asymmetry, was
ablated and ruled out). Ref: §6.6 anti-cherry-pick statement,
\texttt{docs/COMPARISON.md}.

\paragraph{L7 --- Latency comparison conflates transport
classes}\label{l7-latency-comparison-conflates-transport-classes}

Cross-system latency was not captured uniformly. The canonical
(2026-06-15) and rc4 runs did not record per-system latency percentiles
under a normalized transport, and the only side-by-side figures --- from
the superseded 2026-05-24 smoke --- compared nox-mem localhost retrieval
against a Docker-in-Docker HTTP call, i.e.~different transport classes,
not a head-to-head speed claim. nox-mem's operational latency is
therefore reported standalone in §5.7 (KG path p50 2.9 ms, hybrid p50
653 ms, re-measured 2026-06-15), not as a cross-system comparison. A
normalized-transport latency benchmark (all systems behind one gateway)
is future work (§7.2). Ref: \texttt{q3-latency-numbers-2026-05-18},
\texttt{docs/PERFORMANCE.md}.

\paragraph{L8 --- Pain signal is directional but not statistically
significant in
isolation}\label{l8-pain-signal-is-directional-but-not-statistically-significant-in-isolation}

The ``pain-weighted hybrid memory'' framing rests primarily on the
additive salience formula (§5.1.2) and section-aware ranking (§5.1.3).
The \texttt{pain} dimension contributes W\_PAIN = 0.10 of the salience
weight, but its isolated causal contribution has not been validated to
statistical significance: the E10 pain ablation
(\texttt{paper/publication/paper-draft-sec4-7.md} §5.5) reports Δ =
+0.0065 with 95\% CI {[}−0.0143, +0.0338{]} on n = 31, directional but
not significant. The current production corpus has 91.74\% of chunks at
the default \texttt{pain\ =\ 0.2} (as of 2026-06-04), providing
insufficient variance for a precise estimate. A definitive pain signal
ablation requires a corpus where pain scores span the full {[}0.1,
1.0{]} range. Ref: §5.7 honest characterization,
\texttt{d47-path-c-decision}.

\begin{center}\rule{0.5\linewidth}{0.5pt}\end{center}

\subsubsection{7.2 Future Work}\label{future-work}

\paragraph{F1 --- A2 Tier 3 P5: production-ready encrypted memory (in
flight)}\label{f1-a2-tier-3-p5-production-ready-encrypted-memory-in-flight}

Phase 5 of the A2 Tier 3 roadmap targets a full SQLCipher-encrypted
memory store with Ed25519-signed audit checkpoints (P4 deployed via PR
\#294). The signed checkpoint chain enables tamper-evident audit across
destructive operations (\texttt{reindex}, \texttt{consolidate},
\texttt{crystallize}) without requiring a central trust authority. P5
closes the Tier 3 arc by integrating encryption key management with the
existing \texttt{withOpAudit()} wrapper and the Tier 3 reads-audit layer
(P3, PR \#292/\#293). Target deployment: post-GTM Phase 2 launch,
estimated Sun 2026-05-25. Ref:
\texttt{specs/2026-05-24-A2-tier3-crypto-audit-RECON.md},
\texttt{docs/A2-TIER3-MIGRATION-RUNBOOK.md},
\texttt{a1-op-audit-module}.

\paragraph{F2 --- F10 Phase C/D: shadow tracker empirical A/B for
ranking
changes}\label{f2-f10-phase-cd-shadow-tracker-empirical-ab-for-ranking-changes}

F10 Phase A (\texttt{/observability/health.html}) and Phase B
(\texttt{/observability/evals.html}) are deployed (§5.7.4, decision
D53). Phase C targets a shadow-mode query logger that captures
production queries, executes them against a candidate ranking config in
parallel, and accumulates query-level nDCG deltas before any promotion
decision. Phase D operationalizes this into a pre-promotion gate: any
ranking change (boost weight adjustment, mutex threshold, salience
weight) that has not accumulated \textgreater=50 shadow queries with p
\textless{} 0.05 improvement is blocked from reaching the production
endpoint. This closes the observability gap identified in
\texttt{ship-ranking-changes-in-shadow-mode-first} --- currently the
shadow mode is a flag toggle, not an integrated eval pipeline. Ref:
\texttt{docs/ROADMAP.md} F10, decision D53, PR \#207/\#212.

\paragraph{F3 --- Per-method benchmark Phase B: cross-method nDCG
optimization}\label{f3-per-method-benchmark-phase-b-cross-method-ndcg-optimization}

The Q4 per-method benchmark (§6,
\texttt{specs/2026-05-21-per-method-benchmark-comparison.md})
establishes the cross-system baseline. Phase B targets per-query-type
boost calibration: given that keyword queries respond differently from
natural-language queries (G10c §5.1.4), and single-hop vs multi-hop have
opposing mutex trade-offs (G10b §5.1.4), a routing layer that selects
ranking parameters based on query-type classification has measurable
potential upside. Estimated Lab Q1 item. Ref:
\texttt{specs/2026-05-21-per-method-benchmark-comparison.md},
\texttt{g10c-per-style-mutex-2026-05-21}.

\paragraph{F4 --- EverMemBench equivalence: honest comparison against
EverMind-AI
dataset}\label{f4-evermembench-equivalence-honest-comparison-against-evermind-ai-dataset}

\texttt{everos-honest-comparison-benchmark-gap} identifies that
EverMind-AI publishes standardized results on EverMemBench (EverCore
83\% LongMemEval / 93\% LoCoMo; HyperMem 92.73\% LongMemEval). Running
nox-mem on EverMemBench with the same evaluation protocol closes the
benchmark gap and provides a reviewer-grade comparison for the arXiv
submission. This is gated on EverMind-AI repository availability
(currently unavailable) or an alternative comparable dataset. Estimated
Lab Q1 priority if the repository returns. Ref:
\texttt{everos-benchmark-publisher-competitor},
\texttt{benchmark-gap-longmemeval-locomo}.

\paragraph{F5 --- Neural reranker: cross-encoder rerank
post-RRF}\label{f5-neural-reranker-cross-encoder-rerank-post-rrf}

The current retrieval stack terminates at RRF fusion (§4.1). A
cross-encoder reranker --- receiving the top-K RRF candidates and the
original query as a pair --- is the standard next step in multi-stage
retrieval and typically yields +3--8\% nDCG@10 over bi-encoder baselines
(see e.g., Nogueira \& Cho 2019 on MS MARCO). The Autonomy constraint
(\texttt{neural-reranker-evolution-vector}) favors a locally-runnable
cross-encoder (e.g., \texttt{cross-encoder/ms-marco-MiniLM-L-6-v2} via
sentence-transformers, \textasciitilde66MB) over a cloud inference call,
keeping the retrieval stack fully offline-capable. Estimated Lab Q1/Q2.
Ref: \texttt{neural-reranker-as-vetor-evolutivo-pos-rrf},
\texttt{docs/ROADMAP.md} Lab Q1.

\paragraph{F6 --- Lab Q1 scale validation: 250k chunk
corpus}\label{f6-lab-q1-scale-validation-250k-chunk-corpus}

The current production corpus is 94,936 chunks (as of 2026-06-04) ---
already approaching the \textasciitilde100k-vector threshold where exact
search latency degrades. The Lab Q1 roadmap targets a 250k chunk corpus
to validate: (a) sqlite-vec ANN recall at scale (current exact-search;
approximate search becomes necessary past \textasciitilde100k vectors at
reasonable latency targets); (b) salience formula stability (the recency
component decays over a longer history window); (c) FTS5 BM25 IDF
calibration (with more documents, rare-term IDF weights shift). The
\texttt{lab-q1-scale-250k} item has no committed spec yet; it is gated
on \texttt{NOX\_SALIENCE\_MODE=active} remaining stable through the GTM
Phase 2 feedback cycle. Ref: \texttt{docs/ROADMAP.md} Lab Q1,
\texttt{q-a-p-pillars-strategic-pivot-2026-05-17}.

\paragraph{F7 --- Multilingual corpus coverage: Portuguese and
Spanish}\label{f7-multilingual-corpus-coverage-portuguese-and-spanish}

The current evaluation corpus is English-dominant (LongMemEval and
LoCoMo are English datasets; the internal entity-eval golden set mixes
English and Portuguese). The FTS5 \texttt{unicode61} tokenizer with
porter stemmer does not stem Portuguese or Spanish tokens correctly
(e.g., ``decisão'' stems to ``decisa'' rather than ``decid-''; Spanish
gerunds lose morphological overlap). The Gemini semantic layer partially
compensates via cross-lingual embedding space, but there is no explicit
multilingual evaluation. A Portuguese golden set is a natural next step
given the production operational language of the corpus; \textbf{Quati}
(Bueno et al., 2024, arXiv:2404.06976), a Brazilian-Portuguese IR
dataset annotated by native speakers, is a strong candidate benchmark
for that evaluation. This is deferred to post-launch community feedback
intake.

\paragraph{F8 --- GTM Phase 2 launch and community feedback
intake}\label{f8-gtm-phase-2-launch-and-community-feedback-intake}

The Q/A/P roadmap (decision
\texttt{qap-pillars-strategic-pivot-2026-05-17}) defines GTM Phase 2 as
gated on D43 (nox-mem top-3 in at least 2 of 4 key metrics --- nDCG@10,
R@10, MRR, latency). With the canonical and rc4 runs complete, D43 is
satisfied; the public OSS launch is the remaining gate-dependent step,
with timing tracked in \texttt{docs/ROADMAP.md}. Post-launch, community
feedback from the OSS release is expected to surface real-world
limitation patterns not visible in the synthetic golden sets (e.g.,
corpora with high image-to-text OCR content, multi-language mixes, or
very short memory fragments \textless{} 20 words that the current
chunker merges). The feedback cycle directly informs Lab Q2 priorities.
Ref: \texttt{docs/ROADMAP.md} GTM Phase 2, \texttt{docs/gtm/},
\texttt{overnight-automode-push-pattern}.

\begin{center}\rule{0.5\linewidth}{0.5pt}\end{center}

\subsection{8. Knowledge Graph v2}\label{knowledge-graph-v2}

\subsubsection{8.1 Entity Extraction}\label{entity-extraction}

\textbf{v1 (Regex-based)}: Used hardcoded regular expressions for 3
entity types (person, project, agent) with a static alias map for name
normalization. Limited to predefined names, producing 26 entities.

\textbf{v2 (LLM-powered)}: Uses Ollama llama3.2:3b with a structured
extraction prompt. Each chunk is processed with temperature 0.1 for
deterministic output. The LLM returns JSON with entities (name + type)
and relations (source + relation + target).

Extraction results after processing 866 chunks:

{\def\LTcaptype{none} % do not increment counter
\begin{longtable}[]{@{}llll@{}}
\toprule\noalign{}
Metric & Regex v1 & LLM v2 & Improvement \\
\midrule\noalign{}
\endhead
\bottomrule\noalign{}
\endlastfoot
Entities & 26 & 384 & 14.8x \\
Relations & 59 & 529 & 9.0x \\
Entity Types & 3 & 11 & 3.7x \\
\end{longtable}
}

\textbf{Entity Type Distribution:}

{\def\LTcaptype{none} % do not increment counter
\begin{longtable}[]{@{}lll@{}}
\toprule\noalign{}
Type & Count & Description \\
\midrule\noalign{}
\endhead
\bottomrule\noalign{}
\endlastfoot
project & 109 & Software projects, products, repos \\
tool & 67 & Libraries, frameworks, CLI tools \\
concept & 54 & Abstract ideas, patterns, methodologies \\
person & 53 & Team members, contacts, stakeholders \\
organization & 50 & Companies, teams, departments \\
agent & 45 & AI agents in the fleet \\
location & 2 & Geographic references \\
other & 4 & Device, currency, date, computer \\
\end{longtable}
}

\subsubsection{8.2 Temporal Decay and TTL}\label{temporal-decay-and-ttl}

Relations have a 90-day time-to-live (TTL) from creation. The confidence
decay mechanism operates as follows:

\begin{enumerate}
\def\labelenumi{\arabic{enumi}.}
\tightlist
\item
  Relations start with confidence 0.8 (extracted) or 0.9 (confirmed)
\item
  Every 30 days without re-confirmation, confidence drops by 0.1
\item
  Relations below 0.3 confidence receive accelerated 7-day expiry
\item
  Expired relations are deleted during \texttt{kg-prune} execution
\item
  Re-confirmation (observing the same relation in new chunks) resets
  confidence to 0.9 and extends TTL by 90 days
\end{enumerate}

This mechanism ensures the knowledge graph naturally forgets stale
information while reinforcing actively observed patterns.

\subsubsection{8.3 Decision Versioning}\label{decision-versioning}

Architectural decisions are tracked with full version history in the
\texttt{decision\_versions} table. Each decision has a unique key (e.g.,
\texttt{dedup-strategy}, \texttt{fallback-chain}) and supports:

\begin{itemize}
\tightlist
\item
  Version chains with supersession tracking
\item
  Authorship attribution
\item
  Source file provenance
\item
  Current vs.~historical querying
\end{itemize}

10 decisions are currently tracked, covering API key management, LLM
fallback chains, embedding model selection, agent isolation strategy,
and synchronization schedules.

\subsubsection{8.4 Graph Traversal}\label{graph-traversal}

The \texttt{findPath()} function implements BFS (Breadth-First Search)
to discover shortest paths between any two entities. This enables
queries like ``How is Toto connected to nox-mem?'' which traverses
person → project → tool → agent relationships. Maximum depth is
configurable (default: 4 hops).

\begin{center}\rule{0.5\linewidth}{0.5pt}\end{center}

\subsection{9. Cross-Agent Intelligence}\label{cross-agent-intelligence}

\subsubsection{9.1 Agent Expertise
Profiling}\label{agent-expertise-profiling}

Each agent's memory is analyzed to determine its unique expertise based
on chunk type distribution. The dominant chunk type determines the
agent's strength category:

\begin{itemize}
\tightlist
\item
  \textbf{daily} → ``Daily operations \& activity logging''
\item
  \textbf{team} → ``Team coordination \& shared knowledge''
\item
  \textbf{decision} → ``Decision tracking \& rationale''
\item
  \textbf{lesson} → ``Lessons learned \& pattern recognition''
\end{itemize}

Profiles include chunk counts, type breakdowns, top topics (via FTS5
term frequency), and last activity dates.

\subsubsection{9.2 Knowledge Sharing}\label{knowledge-sharing}

The \texttt{pullInsightsFrom()} function enables any agent to query
lessons and decisions from other agents without direct database access.
This creates a knowledge transfer mechanism where, for example, Cipher
(Security) can learn from Forge's (Code Reviewer) past code review
decisions.

\texttt{pullAllInsights()} aggregates insights across all agents, sorted
by date, providing a fleet-wide learning feed.

\subsubsection{9.3 Cross-Agent Knowledge Graph
Merge}\label{cross-agent-knowledge-graph-merge}

\texttt{mergeCrossKnowledgeGraphs()} scans all agent databases for
kg\_entities and kg\_relations tables, merging them into a unified
entity view. Entities are matched by type + lowercase name. The output
shows which entities are known to which agents and their combined
mention counts, enabling identification of shared knowledge
vs.~agent-specific expertise.

\begin{center}\rule{0.5\linewidth}{0.5pt}\end{center}

\subsection{10. MCP Server Interface}\label{mcp-server-interface}

nox-mem exposes 14 tools via the Model Context Protocol (MCP) over stdio
(JSON-RPC 2.0):

{\def\LTcaptype{none} % do not increment counter
\begin{longtable}[]{@{}
  >{\raggedright\arraybackslash}p{(\linewidth - 4\tabcolsep) * \real{0.2069}}
  >{\raggedright\arraybackslash}p{(\linewidth - 4\tabcolsep) * \real{0.3448}}
  >{\raggedright\arraybackslash}p{(\linewidth - 4\tabcolsep) * \real{0.4483}}@{}}
\toprule\noalign{}
\begin{minipage}[b]{\linewidth}\raggedright
Tool
\end{minipage} & \begin{minipage}[b]{\linewidth}\raggedright
Category
\end{minipage} & \begin{minipage}[b]{\linewidth}\raggedright
Description
\end{minipage} \\
\midrule\noalign{}
\endhead
\bottomrule\noalign{}
\endlastfoot
nox\_mem\_search & Retrieval & Hybrid search (FTS5 + semantic + RRF) \\
nox\_mem\_stats & Monitoring & Database statistics and health \\
nox\_mem\_primer & Context & Session recovery summary
(\textasciitilde500 tokens) \\
nox\_mem\_ingest & Ingestion & Index a file into memory \\
nox\_mem\_cross\_search & Cross-Agent & Search across all 7 databases \\
nox\_mem\_cross\_stats & Cross-Agent & Chunk counts per agent \\
nox\_mem\_metrics & Monitoring & Daily observability metrics \\
nox\_mem\_kg\_build & KG & Build knowledge graph from chunks \\
nox\_mem\_kg\_query & KG & Query entity and its relations \\
nox\_mem\_kg\_stats & KG & Knowledge graph statistics \\
nox\_mem\_agent\_profiles & Intelligence & Agent expertise profiles \\
nox\_mem\_cross\_kg & Intelligence & Merged cross-agent knowledge
graph \\
nox\_mem\_kg\_path & Intelligence & BFS path between entities \\
nox\_mem\_self\_improve & Analysis & Contradiction detection, pattern
analysis \\
\end{longtable}
}

\begin{center}\rule{0.5\linewidth}{0.5pt}\end{center}

\subsection{11. HTTP API Server}\label{http-api-server}

A lightweight HTTP API (Node.js built-in \texttt{http} module, zero
dependencies) runs on port 18800, exposing memory data to the React
dashboard:

{\def\LTcaptype{none} % do not increment counter
\begin{longtable}[]{@{}
  >{\raggedright\arraybackslash}p{(\linewidth - 4\tabcolsep) * \real{0.3571}}
  >{\raggedright\arraybackslash}p{(\linewidth - 4\tabcolsep) * \real{0.2857}}
  >{\raggedright\arraybackslash}p{(\linewidth - 4\tabcolsep) * \real{0.3571}}@{}}
\toprule\noalign{}
\begin{minipage}[b]{\linewidth}\raggedright
Endpoint
\end{minipage} & \begin{minipage}[b]{\linewidth}\raggedright
Method
\end{minipage} & \begin{minipage}[b]{\linewidth}\raggedright
Response
\end{minipage} \\
\midrule\noalign{}
\endhead
\bottomrule\noalign{}
\endlastfoot
\texttt{/api/health} & GET & System health: chunks, consolidation,
vector coverage, services, KG stats, DB size \\
\texttt{/api/agents} & GET & Agent expertise profiles array \\
\texttt{/api/kg} & GET & Knowledge graph entities and relations \\
\texttt{/api/kg/path?from=X\&to=Y} & GET & BFS shortest path between
entities \\
\texttt{/api/search?q=QUERY\&limit=N} & GET & Hybrid search results \\
\texttt{/api/cross-kg} & GET & Merged cross-agent knowledge graph \\
\end{longtable}
}

CORS headers are set for cross-origin access from the Vercel-hosted
dashboard.

\begin{center}\rule{0.5\linewidth}{0.5pt}\end{center}

\subsection{12. Operational
Infrastructure}\label{operational-infrastructure}

\subsubsection{12.1 Cron Schedule}\label{cron-schedule}

24 cron jobs manage automated operations:

{\def\LTcaptype{none} % do not increment counter
\begin{longtable}[]{@{}
  >{\raggedright\arraybackslash}p{(\linewidth - 6\tabcolsep) * \real{0.1935}}
  >{\raggedright\arraybackslash}p{(\linewidth - 6\tabcolsep) * \real{0.3548}}
  >{\raggedright\arraybackslash}p{(\linewidth - 6\tabcolsep) * \real{0.1613}}
  >{\raggedright\arraybackslash}p{(\linewidth - 6\tabcolsep) * \real{0.2903}}@{}}
\toprule\noalign{}
\begin{minipage}[b]{\linewidth}\raggedright
Time
\end{minipage} & \begin{minipage}[b]{\linewidth}\raggedright
Frequency
\end{minipage} & \begin{minipage}[b]{\linewidth}\raggedright
Job
\end{minipage} & \begin{minipage}[b]{\linewidth}\raggedright
Details
\end{minipage} \\
\midrule\noalign{}
\endhead
\bottomrule\noalign{}
\endlastfoot
23:00-23:25 & Daily & Agent consolidation & 6 agents, 5-min stagger,
reindex→consolidate \\
23:30 & Daily & Workspace consolidation & Central workspace daily
notes \\
23:35 & Daily & Session wrap-up & SESSION-STATE.md, Notion sync, git
commit \\
04:00 & Weekly (Sun) & Vectorize & Gemini embeddings for new/changed
chunks \\
*/5 min & Continuous & Health check & Watcher heartbeat, service
liveness \\
02:00 & Daily & SQLite backup & Online backup API, 7-day retention
pruning \\
*/6 hours & Continuous & Git backup & Auto-commit memory file changes \\
09:00 & Weekly (Mon) & Token check & Forge CC token verification \\
\end{longtable}
}

\subsubsection{12.2 Backup Strategy}\label{backup-strategy}

Three backup mechanisms operate independently:

\begin{enumerate}
\def\labelenumi{\arabic{enumi}.}
\tightlist
\item
  \textbf{SQLite Online Backup}: Uses better-sqlite3's backup API for
  crash-consistent copies. Daily at 02:00, 7-day retention with
  automatic pruning.
\item
  \textbf{Git Auto-Commit}: Memory directory changes are committed every
  6 hours, providing full change history.
\item
  \textbf{File System}: WAL mode ensures database consistency during
  concurrent reads/writes.
\end{enumerate}

\subsubsection{12.3 LLM Fallback Chain}\label{llm-fallback-chain}

To ensure continuous operation regardless of provider availability:

\textbf{Paid Tier}: Claude Opus → Sonnet → Haiku → GPT-5.1 → Gemini 2.5
\textbf{Free Tier}: Nemotron → Groq Llama70B → Healer → Hunter → Trinity
→ Gemma27B

The fallback is configured in the environment and selected at runtime
based on task complexity and availability.

\begin{center}\rule{0.5\linewidth}{0.5pt}\end{center}

\subsection{13. Dashboard Integration}\label{dashboard-integration}

The TotoClaw Command Center (React 18 + TypeScript + Vite + shadcn/ui)
provides 11 pages including 4 nox-mem-specific views:

\begin{itemize}
\tightlist
\item
  \textbf{Memory Health} (\texttt{/memory}): Real-time system stats,
  vector coverage progress bar, service status indicators, agent
  breakdown table
\item
  \textbf{Knowledge Graph} (\texttt{/knowledge-graph}): Interactive
  force-directed canvas graph, entity type filters, BFS path finder
\item
  \textbf{Agent Intel} (\texttt{/agent-intel}): Agent expertise cards
  with type distribution bars, hybrid search interface, cross-agent
  knowledge entities
\item
  \textbf{System Paper} (\texttt{/system-paper}): Live technical
  analysis with Recharts visualizations (pie, bar, radar, area charts),
  auto-refresh every 60 seconds
\end{itemize}

All data is fetched from the nox-mem API server via TanStack React Query
with configurable polling intervals.

\begin{center}\rule{0.5\linewidth}{0.5pt}\end{center}

\subsection{14. Evolution History}\label{evolution-history}

{\def\LTcaptype{none} % do not increment counter
\begin{longtable}[]{@{}
  >{\raggedright\arraybackslash}p{(\linewidth - 4\tabcolsep) * \real{0.3214}}
  >{\raggedright\arraybackslash}p{(\linewidth - 4\tabcolsep) * \real{0.2143}}
  >{\raggedright\arraybackslash}p{(\linewidth - 4\tabcolsep) * \real{0.4643}}@{}}
\toprule\noalign{}
\begin{minipage}[b]{\linewidth}\raggedright
Version
\end{minipage} & \begin{minipage}[b]{\linewidth}\raggedright
Date
\end{minipage} & \begin{minipage}[b]{\linewidth}\raggedright
Key Changes
\end{minipage} \\
\midrule\noalign{}
\endhead
\bottomrule\noalign{}
\endlastfoot
v1.0 & Mar 14 & SQLite FTS5, basic search, consolidation, Notion sync \\
v2.0 & Mar 17 & MCP server, systemd services, watcher heartbeat,
primer \\
v2.2 & Mar 20 & Cross-agent search, KG v1 (regex), self-improve,
decision versioning \\
v2.5 & Mar 22 & Multi-agent workspace fix (OPENCLAW\_WORKSPACE), gateway
supervision \\
v2.6 & Mar 22 & Hybrid search default (FTS5+Gemini+RRF), 866/866
vectorized \\
v3.0 & Mar 23 & KG v2 (LLM, 384 entities), Cross-Agent Intelligence,
HTTP API, dashboard \\
v3.7 & Apr 23 & Schema V10 (\texttt{retention\_days} v8 + \texttt{pain}
v9 + \texttt{section} v10), entity file format, section\_boost \\
Wave A & May 19 & Additive salience formula, \texttt{tier\_boost}
off-by-default, \texttt{source\_type} backfill (67,949 chunks), G5 V3
ablation (PRs \#150 / \#151 / \#153) \\
G10 Hard Mutex & May 20 & \texttt{section/source\_type} mutex deployed
against \texttt{g9.db} 69,495 chunks (PRs \#181 / \#182) \\
G10d ACTIVE-T2 & May 21 & Conditional mutex gated by
\texttt{query\_entity\_count\ \textless{}=\ 2}, deployed via systemd
drop-in \texttt{NOX\_MUTEX\_QUERY\_ENTITY\_THRESHOLD=2} (PR \#198,
decision D51); multi-hop +1.58\% nDCG / adversarial +3.04\% nDCG
recovered \\
F10 Phase A + B & May 21 & Foundation observability dashboards
(\texttt{/observability/health.html} +
\texttt{/observability/evals.html}) deployed via Tailscale tunnel (PRs
\#207 / \#212, decision D53) \\
\end{longtable}
}

\begin{center}\rule{0.5\linewidth}{0.5pt}\end{center}

\subsection{15. Conclusion}\label{conclusion}

nox-mem demonstrates that persistent, searchable, and shareable memory
for AI agent fleets is achievable with commodity infrastructure (single
VPS, SQLite, local LLM). The hybrid search system consistently
outperforms single-method retrieval, particularly for multilingual
content and compound technical terms. The LLM-powered knowledge graph
provides 15x richer entity extraction compared to regex approaches,
while temporal decay ensures the graph stays current without manual
curation. The Wave A empirical evaluation (§5) established nDCG@10 =
0.6237 on the entity-flavored golden set (+78.8\% relative over the G3
baseline), with \texttt{section\_boost} identified as the dominant
driver (99.85\% of the lift recovered by A3 alone) and the additive
salience formula validated by the
\texttt{active\ \textgreater{}\ shadow} reversal. The G10d conditional
mutex evolution (§5.1.4, deployed 2026-05-21) consolidates the canonical
boost stack
\texttt{section\_boost\ ×\ source\_type\_boost\ (Hard\ Mutex\ gated\ by\ query\_entity\_count\ \textless{}=\ 2)\ ×\ salience\ v2\ additive}
in production, recovering multi-hop and adversarial regressions with
contained dilution on single-hop. The F10 layer (§5.7.4, decision D53)
makes production state verifiable at any time via the Phase A dashboard
(\texttt{/observability/health.html}) + Phase B dashboard
(\texttt{/observability/evals.html}).

The §6 Q4 COMPARISON is pre-registered
(\texttt{specs/2026-05-23-Q4-comparison-execution-plan.md}). After a
methodology smoke (2026-05-24, n=20) and an infrastructure abort
(incident D76, §7.1 L5), the \textbf{canonical run executed 2026-06-15}
(n=100/dataset, 6 systems: 3 produced numbers, 3 documented gaps ---
§6.3), and the \textbf{all-Gemini controlled variant (rc4) executed
2026-06-29} (full n=2,482 --- §6.3.2 / §6.4). Under each system's native
embedder the two leaders split --- Mem0 wins LoCoMo, nox-mem wins
LongMemEval (§6.3); with the embedding provider equalized, nox-mem
outperforms Mem0 on both datasets and all five represented categories
(§6.3.2). Both readings are reported side by side per §6.6, with three
residual confounds declared and the task-type asymmetry ablated away.
The D43 gate (top-3 on at least 2 of the 4 key metrics) is satisfied,
unlocking the GTM Phase 2 pre-launch defense.

The cross-agent intelligence layer transforms isolated agent memories
into a collaborative knowledge base, enabling institutional learning
across the fleet. Combined with the live dashboard, the system provides
full observability into the collective memory of the agent organization.

\textbf{Repository:} github.com/totobusnello/memoria-nox
\textbf{Dashboard:} github.com/totobusnello/agent-hub-dashboard
\textbf{Spec:}
Projetos/memoria-nox/specs/2026-03-14-nox-memory-system-design.md

\begin{center}\rule{0.5\linewidth}{0.5pt}\end{center}

\subsection{References and Footnotes}\label{references-and-footnotes}

\subsubsection{Related-systems
references}\label{related-systems-references}

\subsubsection{Internal references}\label{internal-references}

\subsubsection{Autonomy table (Table 2)
sources}\label{autonomy-table-table-2-sources}

\end{document}
